% DER_XVIII-e: Grothendieck's text with the commentary set in, then the overview, inside a group that applies the chunk's own definitions
\dgrpart{1}{XVIII-e}{Part 5 of 5, manuscript pages 419--19}{XVIII-e\enspace\textperiodcentered\enspace Chapter XVIII, part 5 of 5\enspace\textperiodcentered\enspace manuscript pages 419--19}{Chapter XVIII \textperiodcentered\ \textup{XVIII-e}}
\dgrchunkbegin
% DER_XVIII-e: the definitions of the chunk's preamble, applied inside its group
\let\printedpage\relax\newcommand{\printedpage}[1]{\marginpar{\footnotesize p.~#1}}
\let\msp\relax\newcommand{\msp}[1]{\marginpar{\footnotesize [#1]}}
\let\cE\relax\newcommand{\cE}{\mathcal{E}}
\let\cC\relax\newcommand{\cC}{\mathcal{C}}
\let\cM\relax\newcommand{\cM}{\mathcal{M}}
\let\Ob\relax\newcommand{\Ob}{\operatorname{Ob}}
\let\Fl\relax\newcommand{\Fl}{\operatorname{Fl}}
\let\Hom\relax\newcommand{\Hom}{\operatorname{Hom}}
\let\End\relax\newcommand{\End}{\operatorname{End}}
\let\Ind\relax\newcommand{\Ind}{\operatorname{Ind}}
\let\Ens\relax\newcommand{\Ens}{\mathrm{Ens}}
\let\pr\relax\newcommand{\pr}{\mathrm{pr}}
\let\id\relax\newcommand{\id}{\mathrm{id}}
\let\card\relax\newcommand{\card}{\operatorname{card}}
\let\cardred\relax\newcommand{\cardred}{\operatorname{card}\,\text{\upshape réd}}
\let\defeq\relax\newcommand{\defeq}{\overset{\text{\upshape déf}}{=}}
\let\red\relax\newcommand{\red}{\operatorname{\text{réd}}}
\let\circled\relax\newcommand{\circled}[1]{\tikz[baseline=(c.base)]\node[draw,circle,inner sep=1pt](c){#1};}
\let\ux\relax\newcommand{\ux}{\underline{x}}
\let\Img\relax\newcommand{\Img}{\operatorname{Im}}
\let\Ker\relax\newcommand{\Ker}{\operatorname{Ker}}
\let\redop\relax\newcommand{\redop}{\mathop{\textnormal{réd}}\nolimits}
\let\into\relax\newcommand{\into}{\lhook\joinrel\longrightarrow}
\let\cQ\relax\newcommand{\cQ}{\mathcal{Q}}
\let\uS\relax\newcommand{\uS}{\underline{S}}
\let\uT\relax\newcommand{\uT}{\underline{T}}
\let\U\relax\newcommand{\U}{\mathfrak{U}}
\let\cN\relax\newcommand{\cN}{\mathcal{N}}
\let\cF\relax\newcommand{\cF}{\mathcal{F}}
\let\cA\relax\newcommand{\cA}{\mathcal{A}}
\let\Max\relax\newcommand{\Max}{\operatorname{Max}}
\let\Ann\relax\newcommand{\Ann}{\operatorname{Ann}}
\let\incl\relax\newcommand{\incl}{\lhook\joinrel\longrightarrow}
\let\limind\relax\newcommand{\limind}{\mathop{\underrightarrow{\mathrm{lim}}}}
\let\limproj\relax\newcommand{\limproj}{\mathop{\underleftarrow{\mathrm{lim}}}}

% ---- from der-ch18_p161-165.tex
\setcounter{footnote}{174}

%=====================================================================
\printedpage{159}
% head: A. Grothendieck : Dérivateurs, Ch. XVIII, Édition des Universités de Montpellier II et Paris VII

\msp{page 419}
est $\pi$-accessible, quand $x_0$ est dans $\cE_{0\,\pi}$. Or par hypothèse sur $\pi$, il est $\pi$-accessible si $x_0$ est dans $\cE_{0\,\pi_0}$. Mais si $x$ est dans $\cE_{0\,\pi}$, on aura % [sic] x -> x_0, the formula that follows concerns x_0
\[
x_0 \;=\; \varinjlim_{i\in I} x_{0\,i} \qquad
\left\{
\begin{array}{l}
I\ \ \text{grande}\ \ \text{devant}\ \ \pi_0,\\
\text{les}\ \ x_{0\,i}\ \ \text{dans}\ \ \cE_{0\,\pi_0},\\
\card I \leqslant \pi
\end{array}
\right.
\]
(car $\pi_0$ est utile pour $\cE_0$ par hypothèse, et $\pi$ satisfait $\pi^{\pi_0} = \pi \geqslant \cardred \cE_{\pi_0}$). Comme $x_0 \longrightarrow H(x_0,x_1) : \cE_0 \longrightarrow \Ens$ satisfait à $\mathrm{L}_{\pi_0}$, pour tout $x_1$ dans $\cE_1$, on aura donc % [sic] x_0 \longrightarrow H(x_0,x_1) -> x_0 \longmapsto H(x_0,x_1), elsewhere such functors are written with \longmapsto
\[
H(x_0) \;\simeq\; \varprojlim_{i\in I} H(x_{0\,i})
\]
(où pour $x_0$ dans $\cE_0$, $H(x)$ désigne le foncteur $x_1 \longmapsto H(x_0,x_1) : \cE_1 \longrightarrow (\Ens)$). Comme % [sic] H(x) -> H(x_0), the functor depends on x_0
les foncteurs $H(x_{0\,i})$ sont $\pi$-accessibles, et $\card I \leqslant \pi$, il s'ensuit que leur $\varprojlim\limits_{I}$ $H(x_0)$ est $\pi$-accessible (SGA 4 I 9.6 (ii)). Cela achève de prouver donc 2°).

3°). Pour prouver que $\Ind_\pi(\cE(\pi)) \longrightarrow \cE$ est essentiellement surjectif, il suffit de voir que tout $x \in \Ob\cE_0$ et tout $x \in \Ob\cE_1$
\msp{page 420}
est dans l'image essentielle. Or cela résulte de l'hypothèse que $\pi$ est utile pour $\cE_0$ et $\cE_1$, et de 1°) qui implique que $\cE_0 \longrightarrow \cE$ et $\cE_1 \longrightarrow \cE$ satisfont $\mathrm{L}_\pi$.

\dgrblock{dgrb1110}{The two cardinal bounds in the gluing argument}
\begin{dcomment}{exposition}{p419}{The two cardinal bounds in the gluing argument}
Grothendieck is finishing a construction of an accessible category from an open piece, a closed piece, and the arrows crossing from the first to the second. The point of the cardinal calculation is to make presentability in a piece agree with presentability in the whole category. The latter has an additional test: an object of the open piece must commute with sufficiently filtered colimits whose values lie in the closed piece. That test is measured by the gluing functor $H$.

We use the modern convention that a regular cardinal $\kappa$ controls diagrams with fewer than $\kappa$ arrows. Thus, for an infinite $\pi$, a diagram ``grand devant $\pi$'' is $\pi^+$-filtered, and an object called $\pi$-accessible here is $\pi^+$-presentable. A useful $\pi$ gives a small family of such objects generating by those colimits. The shift by one successor is essential when interpreting the inequalities $\operatorname{card}I\leq\pi$ in the argument.

There are two distinct bounds. First, $\pi_0$ controls the presentation of an open object and the continuity of $H$ in its contravariant variable. Second, $\pi$ controls the number of terms in a presentation and the continuity of the mixed hom functors in their covariant variable. For a presentation $x_0=\mathop{\rm colim}_{j\in J}a_j$ with $|\operatorname{Mor}J|\leq\pi$, the calculation is
\[
 \begin{aligned}
 H(x_0,\mathop{\rm colim}_{i\in I}y_i)
 &\cong\lim_{j\in J}H(a_j,\mathop{\rm colim}_{i\in I}y_i)\\
 &\cong\lim_{j\in J}\mathop{\rm colim}_{i\in I}H(a_j,y_i)\\
 &\cong\mathop{\rm colim}_{i\in I}H(x_0,y_i).
 \end{aligned}
\]
The first and last identifications use continuity in the first variable at the rank of the $J$-presentation. The middle interchange uses the $\pi^+$-filteredness of $I$. To see the latter directly in sets, choose representatives for the at most $\pi$ components of a compatible family, then choose one upper bound for their indices and for the at most $\pi$ equations expressing compatibility. The same common-bound argument proves injectivity. This is the mechanism of SGA~4, I.9.6(ii) and I.9.8 \citep{SGA4I1972}.

This distinction will matter again on manuscript page~430. Increasing the rank at which a functor is known to preserve filtered colimits does not give continuity at the smaller rank used to present an object. Here the required smaller-rank continuity is explicitly assumed in Corollary~4.14.25(2); later it is inadvertently lost from the numerical hypotheses.
\end{dcomment}


Cela achève de prouver la proposition 4.14.24, et de même le résultat suivant, plus précis que l'implication c) $\Longrightarrow$ a).

\textbf{Corollaire 4.14.25.} \textit{Soient $\cE$ une catégorie, $\cE_0$ une sous-catégorie ouverte, $\cE_1$ la sous-catégorie fermée complémentaire,}
\[
H \;=\; ((x_0,x_1) \longmapsto \Hom_{\cE}(x_0,x_1)) \;:\; \cE_0^\circ \times \cE_1 \;\longrightarrow\; \Ens
\]
\textit{le foncteur recollement.}
\begin{enumerate}
\item[\textit{1°)}] \textit{Soit $\pi$ un cardinal infini. Pour ce que $\cE$ et $\cE_0$ satisfassent à $L_\pi$, et le foncteur d'inclusion $\cE_0 \hookrightarrow \cE$ itou, il faut et il suffit que $\cE_0$ satisfasse à $L_\pi$, et que pour tout $x_1$ dans $\cE_1$, le foncteur $x_0 \longmapsto H(x_0,x_1) : \cE_0 \longrightarrow \Ens^\circ$ satisfasse à $L_\pi$. Alors $\cE_1$ et $\cE_1 \hookrightarrow \cE$ satisfont aussi à $L_\pi$,}
\msp{page 421}
\textit{et on a}
\[
\cE_\pi \;\subset\; \cE(\pi) \;\defeq\;
\begin{array}{l}
\text{\itshape sous-catégorie\ \ strictement\ \ pleine\ \ de\ \ $\cE$\ \ telle\ \ que}\\
\Ob\cE(\pi) = \Ob\cE_{0\,\pi} \cup \Ob\cE_{1\,\pi}.
\end{array}
\]
\item[\textit{2°)}] \textit{Supposons que la condition c) de 4.14.24 soit satisfaite, et soit $\pi_0$ un cardinal infini tel que $\pi_0$ soit utile pour $\cE_0$, et que les foncteurs $x_0 \longmapsto H(x_0,x_1) : \cE_0 \longrightarrow \Ens^\circ$ définis par les $x_1 \in \Ob\cE_1$ soient $\pi_0$-accessibles. Soit $\pi$ un cardinal tel que $\pi^{\pi_0} = \pi \geqslant \sup(\pi_0, \cardred\cE_{0\,\pi_0})$, et tel que pour tout $x_0 \in \Ob\cE_{\pi_0}$, le foncteur $x_1 \longmapsto H(x_0,x_1) : \cE_1 \longrightarrow \Ens$ soit $\pi$-accessible. Alors on a}
\[
\cE_\pi \;=\; \cE(\pi)\ .
\]
% foot: Version 17 décembre 2025

\dgrblock{dgrb1111}{The closed piece must have its colimits}
\begin{dcomment}{examples}{p420}{The closed piece must have its colimits}
The purpose of the first part of Corollary~4.14.25 is to separate existence of colimits from the presentability calculation. As printed, however, its right-hand condition says nothing about the existence of colimits in $\mathcal E_1$. This is a genuine omission in the edition. It cannot follow from continuity of the mixed hom functor.

Here is a finite example detecting the omission at every infinite rank. Let $B$ have one object $b$ and endomorphism monoid $\{1,p\}$, where $p^2=p\ne1$. Put
\[
 \mathcal E=\{a\}\amalg B,\qquad
 \mathcal E_0=\{a\},\quad\mathcal E_1=B.
\]
The gluing functor is constantly empty. For any nonempty filtered diagram, the limit of the constant empty-set diagram is empty, so the printed condition on $H$ is satisfied. The terminal category $\mathcal E_0$ has all the required colimits.

To exhibit the failure in the closed piece, index a diagram by the ordinal $\pi^+$, send every index to $b$, and send every strict transition to $p$. This is a functor because $p^2=p$. The index is $\pi^+$-filtered. A cocone to $b$ must have every leg equal to $p$: for each index $i$, choose $j>i$ and use $c_i=c_jp=p$. Thus there is exactly one cocone to $b$, whereas $B(b,b)$ has two elements. No object represents the cocone functor, and the diagram has no colimit in $B$ or in $\mathcal E$.

The corrected criterion requires both pieces to admit $L_\pi$-colimits, and requires $H(-,x_1)$ to carry those in the open piece to limits. Its proof also explains why no further compatibility is needed. For a filtered diagram $D:I\to\mathcal E$, let $I_1$ consist of the indices whose values lie in the closed piece. If $I_1$ is nonempty, choose one of its objects. Every finite set of indices has an upper bound receiving an arrow from that chosen object, so $I_1$ is cofinal; the same argument with at most $\pi$ indices gives the required higher filteredness. Its colimit is computed in $\mathcal E_1$. There are no arrows from the closed piece to the open one, so this is also a colimit in $\mathcal E$. If $I_1$ is empty, compute the colimit in $\mathcal E_0$; its universal property against targets in $\mathcal E_1$ is precisely the asserted continuity of $H$. Conversely, an existing total colimit of a nonempty diagram in $\mathcal E_1$ must lie in $\mathcal E_1$, since its cocone cannot point into the open piece.
\end{dcomment}


%=====================================================================
\printedpage{160}
% head: A. Grothendieck : Dérivateurs, Ch. XVIII, Édition des Universités de Montpellier II et Paris VII

\item[\textit{3°)}] \textit{Si $\pi_0$ et $\pi$ sont comme dans 2°), et si de plus $\pi$ est utile pour $\cE_1$, alors $\pi$ est utile pour $\cE$, donc (en vertu de 2°)}
\[
\cE \;\xleftarrow{\ \sim\ }\; \Ind_\pi(\cE(\pi))\ .
\]
\end{enumerate}
Nous pouvons maintenant énoncer la généralisation cherchée de 4.14.17 :

\msp{page 422}
\textbf{Théorème 4.14.26} (généralisation de 4.14.13). \textit{Soit $p : \cE \longrightarrow \cC$ un foncteur, avec $\cC$ petite. On suppose que $p$ satisfait à l'hypothèse H) de la page 409, qui équivaut (moyennant 4.14.24) à la suivante :}
\begin{enumerate}
\item[\textit{H$'$)}] \textit{$\forall\ s \in \Ob\cC$, la catégorie fibre $\cE_s$ est accessible, et pour tout $\lambda : s \longrightarrow t$ dans $\Fl\cC$, le foncteur}
\[
\left\{
\begin{array}{lcl}
H_\lambda : \cE_s^\circ \times \cE_t & \longrightarrow & \Ens\\[1ex]
\qquad (x_0,x_1) & \longmapsto & \Hom_\lambda(x_0,x_1)
\end{array}
\right.
\]
\textit{satisfait les deux conditions suivantes :}
\begin{enumerate}
\item[\textit{a)}] \textit{$\forall\ x_0 \in \Ob\cE_s^\circ$, $x_1 \longmapsto H_\lambda(x_0,x_1) : \cE_t \longrightarrow \Ens$ est accessible.}
\item[\textit{b)}] \textit{Il existe un cardinal infini $\pi_\lambda$, tel que pour tout $x_1 \in \Ob\cE_t$, le foncteur $x_0 \longmapsto H_\lambda(x_0,x_1) : \cE_s \longrightarrow (\Ens)^\circ$ soit $\pi_\lambda$-accessible.}
\end{enumerate}
\end{enumerate}
\textit{Ceci posé, la catégorie $\underline{\Gamma}(\cE/\cC)$ des sections de $\cE$ sur $\cC$ est accessible, et les foncteurs $\varepsilon_s : \underline{\Gamma}(\cE/\cC) \longrightarrow \cE_s$ sont accessibles. On peut préciser de façon similaire à 4.14.18 (énoncé en forme laissé au lecteur).}

\msp{page 423}
\textsc{Démonstration} de 4.14.26. Elle se fait exactement comme celle de 4.14.17 (p. 400 ff.), mais est conceptuellement plus simple. On suppose $\cC$ petite, et on introduit les catégories
\[
\Pi \;=\; \prod_{s\,\in\,\Ob\cC} \cE_s\ , \qquad \Pi' \;=\; \prod_{\lambda\,\in\,\Fl\cC} \cE_\lambda\ ,
\]
où
\[
\cE_\lambda \;=\; \cE \times_{\cC} (\mathbf{\Delta}_1, i_\lambda)\ ,
\]
et les deux foncteurs canoniques
\[
\Pi \;\underset{G}{\overset{F}{\rightrightarrows}}\; \Pi'
\]
de composantes
\[
\begin{aligned}
F_\lambda \;&=\; i_\lambda \circ \pr_{s(\lambda)}\\
G_\lambda \;&=\; j_\lambda \circ \pr_{b(\lambda)}\ ,
\end{aligned}
\]
où pour $\lambda : s \longrightarrow t$, on pose
\[
s(\lambda) \;=\; s\ , \qquad b(\lambda) \;=\; t\ ,
\]
et $i_\lambda$, $j_\lambda$ sont les inclusions canoniques
% In the print, the vertical arrow p_lambda and Delta_1 are enclosed together in a pair of large parentheses.
\[
\begin{tikzcd}
\cE_s \arrow[r,"i_\lambda"] \arrow[d,dash,"\wr"] & \cE_\lambda \arrow[dd,"p_\lambda"] & \cE_t\ . \arrow[l,"j_\lambda"'] \arrow[d,dash,"\wr"] \\
(\cE_\lambda)_0 & & (\cE_\lambda)_1 \\
& \mathbf{\Delta}_1 &
\end{tikzcd}
\]
% foot: Version 17 décembre 2025

\dgrblock{dgrb1112}{Sections retain the equations between arrows}
\begin{dcomment}{exposition}{p422}{Sections retain the equations between arrows}
Theorem~4.14.26 changes the task from assembling objects to assembling a diagram. A section chooses $X_s\in\mathcal E_s$ for every base object and a lift $u_\lambda:X_s\to X_t$ for every $\lambda:s\to t$. These choices must satisfy the unit and composition equations. Grothendieck first constructs all the choices and only then imposes the equations; this order makes the accessibility argument possible.

For each $\lambda$, the pullback $\mathcal E_\lambda$ is the collage of its two fibres with crossing hom functor $H_\lambda$. An arrow from the zero copy of $X_s$ to the one copy of $X_t$ is therefore exactly a lift of $\lambda$. The product $\Pi$ records all objects $X_s$, while the inserter of $F,G:\Pi\rightrightarrows\Pi'$ records these objects together with all arrows $u_\lambda$. A morphism in the inserter is a family $a_s:X_s\to Y_s$ satisfying
\[
 a_tu_\lambda=v_\lambda a_s
 \quad\text{for every }\lambda:s\to t.
\]
These are already the naturality equations for a map of sections. What is still missing is functoriality of the chosen arrows themselves.

The asymmetry in H$'$ has a structural reason. A single cardinal must control the contravariant variable for all possible targets, because it ensures that colimits in a fibre remain colimits when one tests them against every object elsewhere. In the covariant variable one first tests only a small family of generators; their individual accessibility bounds can then be dominated by one larger cardinal. Confusing these two quantifiers would replace a genuine uniformity assumption by a pointwise statement.

This construction continues the older SGA~4 treatment of accessible section categories, particularly I.9.22--9.26, while allowing crossing hom functors which need not come from chosen cartesian transport \citep{SGA4I1972}. In current terminology, the argument belongs to the closure of accessible categories under products, inserters, and equifiers. Makkai--Par\'e's monograph appeared in 1989: its formalism was contemporary with this manuscript, rather than a subsequent solution. The relevant result is Theorem~5.1.6; an explicit modern use of precisely this sequence of constructions appears in the appendix of \citet{ChingRiehl2014}. The use of a two-categorical limit theorem is substantive: arbitrary strict equalizers of functors between accessible categories do not enjoy the same general closure property.
\end{dcomment}


%=====================================================================
\printedpage{161}
% head: A. Grothendieck : Dérivateurs, Ch. XVIII, Édition des Universités de Montpellier II et Paris VII

\msp{page 424}
On désigne par $\cM'$ la catégorie
\[
\cM' \;=\;
\begin{array}{l}
\text{catégorie des couples } (X,u)\text{, avec}\\
X \in \Ob\Pi,\ u : F(x) \longrightarrow G(x),
\end{array}
\] % [sic] F(x), G(x) -> F(X), G(X), the object is named X
qui est constructible par 4.14.22.1 est accessible (car $\Pi$ et $\pi'$ et les foncteurs $F$ et $G$ le sont, par l'hypothèse H) faite sur $p$). On peut l'interpréter comme la catégorie des sections % [sic] \pi' -> \Pi', the category introduced is \Pi'
de $\mathrm{Diag}(\cE)$ sur $\mathrm{Diag}(\cC)$. On a un foncteur strictement pleinement fidèle
\[
\cM \;\defeq\; \underline{\Gamma}(\cE/\cC) \;\hookrightarrow\; \cM'\ ,
\]
qui permet d'interpréter $\cM$ comme une sous-catégorie strictement pleine de $\cM'$. On voit alors comme dans le cas de 4.14.17 que cette sous-catégorie est accessible dans $\cM'$ -- c'est essentiellement sorital. Il s'ensuit que $\cM$ est accessible et que les foncteurs $\cM \longrightarrow \cE_s$ le sont aussi, comme composés de
\[
\cM \;\overset{\text{\upshape acc.}}{\hookrightarrow}\; \cM' \;\overset{\text{\upshape acc.}}{\hookrightarrow}\; \Pi \;\xrightarrow{\ \ \pr_s\ \ \text{\upshape acc.}\ \ }\; \cE_s\ ,
\]

\hfill q.e.d.

\msp{page 425}

\dgrblock{dgrb1113}{Writing out the omitted equifiers}
\begin{dcomment}{later}{p424}{Writing out the omitted equifiers}
Grothendieck's description of the last step as ``essentiellement sorital'' can be made precise without adding an assumption to the theorem. For each object $s$, impose
\[
 u_{1_s}=1_{X_s},
\]
and, for each composable pair $s\xrightarrow{\lambda}t\xrightarrow{\mu}v$, impose
\[
 u_{\mu\lambda}=u_\mu u_\lambda.
\]
Both sides are natural transformations between the same accessible functors on the inserter. For example, the two transformations in the second equation take values in the mixed-arrow collage $\mathcal E_{\mu\lambda}$, from its zero copy of $X_s$ to its one copy of $X_v$. Naturality follows from the commutative squares already built into morphisms of the inserter. The unit equations have the same description using $\mathcal E_{1_s}$.

An equifier is the full subcategory on the objects where two specified natural transformations agree. The smallness of $\mathcal C$ makes the collection of unit and composition equations small; a product gathers their target categories, and equifiers impose all the equations. Accessibility of these constructions and of their structural functors gives both assertions of the theorem. The notation $\mathcal M'\hookrightarrow\Pi$ in the last display should be read as a forgetful arrow: different arrow assignments may have the same underlying family of objects, and a family of object maps need not satisfy the naturality squares. The projection is accessible, which is all the proof uses.

For a concrete illustration, take $\mathcal C=B\mathbf N$, the one-object category freely generated by one endomorphism, and take $\mathcal E=\mathcal C\times\mathbf{Set}$. A section is a set $X$ with an endomorphism $T$. Before imposing the equations, one has an unrelated family of self-maps $u_n$ for $n\in\mathbf N$; afterward $u_0=1$ and $u_n=T^n$. The free section on one element is $(\mathbf N,n\mapsto n+1)$: an equivariant map from it is uniquely determined by the image of zero. Hence this section is finitely presentable, although its underlying set is infinite. Evaluation preserves filtered colimits, but it need not preserve finitely presentable objects at that same rank. The distinction between accessibility and strong accessibility therefore appears even in this elementary example.

The product--inserter--equifier method is part of the broader theory of flexible limits developed by \citet{BirdKellyPowerStreet1989}. Its relevance here is exact: the equations constrain specified arrows between already chosen objects, rather than requiring unrelated object-valued functors to become literally equal.
\end{dcomment}


\subsection*{4.15\quad Catégorie accessible sur une autre}

\textbf{Proposition 4.15.1.} \textit{Soit $f : \cE \longrightarrow \cC$ un foncteur avec $\cC$ petite. On suppose de plus satisfaite la condition H (p. 409), qu'on explicite ici sous la forme équivalente H$'$ :}
\begin{enumerate}
\item[\textit{H 1)}] \textit{Pour tout $s \in \Ob\cC$, $\cE_s$ est accessible.}
\item[\textit{H 2)}] \textit{Pour toute $\lambda \in \Fl\cC$, $\lambda : s \longrightarrow t$, existe un cardinal infini $\pi_\lambda$ tel que pour tout $x_1 \in \Ob\cE_t$, le foncteur}
\[
x_0 \;\longmapsto\; \Hom_\lambda(x_0,x_1) \;:\; \cE_s \;\longrightarrow\; (\Ens)^\circ
\]
\textit{est $\pi_\lambda$-accessible ($\cE_s$ satisfaisant de plus à $L_{\pi_\lambda}$).}
\item[\textit{H 3)}] \textit{Pour toute $\lambda \in \Fl\cC$, $\lambda : s \longrightarrow t$, et pour tout $x_0 \in \Ob\cE_s$, le foncteur}
\[
x_1 \;\longmapsto\; \Hom_\lambda(x_0,x_1) \;:\; \cE_t \;\longrightarrow\; (\Ens)
\]
\textit{est accessible.}
\end{enumerate}
\textit{Alors les foncteurs d'inclusion $i_s : \cE_s \longrightarrow \cE$ sont accessibles. De plus, si $\cC$ ``est sans projecteurs'' (i.e. si tout endomorphisme $p$ d'un objet $s$ de $\cC$ tel que $p^2 = p$, est l'identité), alors $\cE$ est accessible, et $f : \cE \longrightarrow \cC$ est accessible.}

Cela résulte du résultat plus précis :
% foot: Version 17 décembre 2025

%=====================================================================
\printedpage{162}
% head: A. Grothendieck : Dérivateurs, Ch. XVIII, Édition des Universités de Montpellier II et Paris VII

\msp{page 426}
\textbf{Corollaire 4.15.2} (\footnote{Pour a) et b) on n'utilise que H 1 et H 2, à l'exclusion de H 3.}). \textit{Les hypothèses sur $f : \cE \longrightarrow \cC$ sont celles de 4.15.1 ($\cC$ petite, et $f$ satisfait $H' =$ H 1) 2) 3)).}
\begin{enumerate}
\item[\textit{a)}] \textit{Soit $s \in \Ob\cC$ et soit $\pi_1$ un cardinal infini tel que pour toute $\lambda : s \longrightarrow t$ (flèche de $\cC$ de source $t$), on ait $\pi_1 \geqslant \pi_\lambda$ (où $\pi_\lambda$ est défini dans l'énoncé de H 2). Alors le foncteur $i_s : \cE_s \longrightarrow \cE$ est $\pi_1$-accessible.} % [sic] source t -> source s, lambda : s -> t has source s (cf. p. 163 "de source s")
\item[\textit{b)}] \textit{On suppose $\cC$ sans projecteurs. Soit $\pi_1$ un cardinal $\geqslant \sup(\pi_\lambda, \card\cC)$. Alors $\cE$ satisfait $L_{\pi_1}$, et $f : \cE \longrightarrow \cC$ est $\pi$-accessible.}
\item[\textit{c)}] \textit{Supposons encore que $\cC$ soit sans projecteurs (p.ex. une catégorie ``rigide'', p.ex. une catégorie préordonnée). Soient $\pi_0$, $\pi_1$ des cardinaux tels que :}
\begin{enumerate}
\item[\textcircled{\footnotesize $1^o$}] $\forall\ s \in \Ob\cC$, $\pi_0$ est utile pour $\cE_s$ \textit{(un tel $\pi_0$ existe grâce à H 1 et au fait que $\cC$, donc $\Ob\cC$ petits).}
\item[\textcircled{\footnotesize $2^o$}] $\pi_1 \geqslant \sup(\pi_0, \card\cC, \sup_{\lambda\in\Fl\cC}\pi_\lambda, \sup_{s\in\Ob\cC} \cardred \cE_{s\,\pi_0})$.
\item[\textcircled{\footnotesize $3^o$}] \textit{$\forall\lambda \in \Fl\cC$, et $x_0 \in \Ob\cE_{s,\pi_0}$, le foncteur $H_\lambda(x_0)$ de H 3, $x_1 \longmapsto \Hom_\lambda(x_0,x_1) : \cE_t \longrightarrow \Ens$ est $\pi_1$-accessible.}
\end{enumerate}
\textit{Soit $\pi$ un cardinal tel que $\pi^{\pi_0} = \pi \geqslant \pi_1$. Alors par b), $\cE$ satisfait $L_{\pi_1}$, donc $L_\pi$, et $f : \cE \longrightarrow \cC$ est $\pi_1$-accessible. De plus,}
\[
i_s(\cE_{s\,\pi}) \;\subset\; \cE_\pi\ ,
\]
\textit{ou encore}
\[
\cE(\pi) \;\subset\; \cE_\pi\ , \tag*{$(*)$}
\]
\textit{où $\cE(\pi)$ est la sous-catégorie pleine de $\cE$ telle que}
\[
\Ob\cE(\pi) \;=\; \bigcup_{s\in\Ob\cC} \Ob\,(\cE_{s\,\pi})\ .
\]
\textit{Comme $\cE_\pi$ satisfait $L_{\pi_1}$ par b), donc a fortiori $L_\pi$, on a donc par $(*)$ un foncteur} pleinement fidèle
\[
\Ind_\pi(\cE_{s\,\pi}) \;\xrightarrow{\ \sim\ }\; \cE\ ,
\]
\textit{et ce foncteur est une équivalence de catégories (donc $\cE$ est accessible).}
\item[\textit{d)}] \textit{En fait, sous les hypothèses de c) sur $\cC$, $\pi_0$, $\pi_1$, $\pi$, on a $\cE(\pi) = \cE_\pi$ (non seulement $\cE(\pi) \subset \cE_\pi$).}
\end{enumerate}

\msp{page 427}
\textsc{Démonstration de 4.15.1 et 4.15.2.} Il est clair que 4.15.1 est contenu dans 4.15.2 a) et c). Nous allons donc prouver 4.15.2.

a) Soit $(x_i)_{i\in I}$ un système inductif dans $\cE_s$, avec $I$ grand devant $\pi_0$, et $x$ sa limite inductive dans $\cE_s$ (qui existe, car $\cE_s$ satisfait à $\mathrm{L}_{\pi_\lambda}$, donc à $\mathrm{L}_{\pi_0}$). Il faut prouver que c'est une limite dans $\cE$, donc que pour tout $y \in \Ob\cE$, la flèche canonique
\[
\Hom_{\cE}(x,y) \;\longrightarrow\; \varprojlim_{I} \Hom_{\cE}(x_i,y) \tag*{$(*)$}
\]
% foot: Version 17 décembre 2025

%=====================================================================
\printedpage{163}
% head: A. Grothendieck : Dérivateurs, Ch. XVIII, Édition des Universités de Montpellier II et Paris VII
est un isomorphisme. Or $\Hom_{\cE}(x,y)$ et les $\Hom_{\cE}(x_i,y)$ sont des objets de $\Ens/\Hom_{\cC}(s,t)$, où $s = f(x)$, $t = f(y)$, et les flèches de transition sont des flèches dans $\Ens/\Hom_{\cC}(s,t)$. Cela montre que la bijectivité de $(*)$ revient à celle de
\[
\Hom_\lambda(x,y) \;\longrightarrow\; \varprojlim_{I} \Hom_\lambda(x_i,y) \tag*{$(*)$}
\]
pour tout $\lambda \in \Hom_{\cC}(x,y)$ \texttt{[plutôt} $\Hom_{\cC}(f(x),f(y))$\texttt{]}. Pour $\lambda$ fixé, et $y$ fixé dans $\Ob\cE_t$, la bijectivité dans $(*)$, pour tout système inductif grand devant $\pi_1$ dans $\cE_s$, de limite $x$, signifie que le foncteur
\[
x \;\longmapsto\; \Hom_\lambda(x,y) \;:\; \cE_s \;\longrightarrow\; (\Ens)^\circ
\]
est $\pi_0$-accessible. Or c'est bien le cas, par hypothèse sur $\pi_1$, quel que soit la flèche $\lambda : s \longrightarrow t$ de source $s$, et quel que soit $y \in \Ob\cE_t$.

\msp{page 428}
Soit $(x_i)_{i\in I}$ un système inductif grand devant $\pi_1$, et $(s_i)_{i\in I}$ l'image de ce système dans $\cC$. Nous prouverons plus bas le

\dgrblock{dgrb1114}{Colimits in a fibre and colimits in the total category}
\begin{dcomment}{exposition}{p427}{Colimits in a fibre and colimits in the total category}
Grothendieck now asks whether the same local data control the total category itself. The first step, preservation of colimits by a fibre inclusion, is more elementary than the final theorem and does not need H3. Let $x=\mathop{\rm colim}_i x_i$ be a sufficiently filtered colimit in $\mathcal E_s$, and let $y\in\mathcal E_t$. The decomposition
\[
 \mathcal E(x,y)=
 \coprod_{\lambda\in\mathcal C(s,t)}
       \operatorname{Hom}_\lambda(x,y)
\]
is respected by every transition map in the comparison with $\mathcal E(x_i,y)$. Since the index category is nonempty and connected, a compatible family in the limit has one fixed base arrow $\lambda$. Thus its universal property reduces to the separate universal properties of $H_\lambda(-,y)$, which H2 supplies. The occurrences of $\pi_0$ in this proof are to be read as $\pi_1$, the bound specified in part~(a).

The second step explains the hypothesis on the base. If the image of a sufficiently filtered diagram can be reindexed to become constant with identity transitions, the diagram lies in one fibre, so the preceding calculation supplies its total colimit. A small base with no nonidentity idempotents has this stabilization property at a sufficiently large rank. The proof of the required small-category statement is deferred until manuscript pages~440--463. For the existential conclusion of Proposition~4.15.1 it is enough to use the strict cardinal bound proved there; no claim about an optimal boundary cardinal is needed.

The phrase ``sans projecteurs'' must be kept distinct from ``karoubienne''. A groupoid has no nonidentity idempotents, although it may have many automorphisms. A Karoubian category may have nonidentity idempotents, provided their images exist. A constant-object diagram whose strict transitions are one such idempotent need not become a diagram of identity maps inside its original fibre. The later difficulty with a Karoubian base is therefore a difficulty of transport to the image of an idempotent, rather than a failure of the fibrewise universal property just proved.
\end{dcomment}


% In the left margin, beside the statement of Lemme 4.15.3 and beside the following paragraph "Mais attention ...", the edition prints a wavy marginal sign (one for each).
\textbf{Lemme 4.15.3} (\footnote{Cf. démonstration page 452.},\footnote{sous réserve de la validité de 4.8.13 (Mais c'est à présent établi, sous une forme suffisante, dans 4.15.10).}). \textit{Soit $\cC$ une catégorie, $\pi$ un cardinal, tels que} $\card\cC \leqslant \pi$\textit{, et que $\cC$ soit sans projecteurs. Alors pour tout système inductif $(s_i)_{i\in I}$ grand devant $\pi$ dans $\cC$, existe une partie cofinale $J$ dans $I$ telle que le système induit $(s_i)_{i\in J}$ soit} constant, \textit{i.e. tel que les $s_i$ soient égaux à un $s \in \cC$, et que les morphismes de transition soient tous égaux à} $\id_s$.

Mais attention, on s'appuie pour prouver 4.15.3 sur le théorème 4.8.13 (page 212-216), que je n'ai toujours pas prouvé. Dans le cas où $\cC$ est rigide (i.e. $\End_{\cC}(s) = \{\id_s\}\ \forall\ s \in \Ob\cC$), la démonstration est cependant immédiate, cf. plus bas.

\msp{page 429}
Pour prouver que $x = \varinjlim\limits_{I} x_i$ existe dans $\cE$, et commute à $f$, on peut remplacer $I$ par une sous-catégorie cofinale. Cela nous ramène, grâce au lemme, au cas où $(s_i)_{i\in I}$ est constant, donc au cas où le système inductif provient d'un système inductif dans une catégorie fibre $\cE_s$. Mais par a), on sait que sa limite $x$ dans $\cE_s$ est une limite dans $\cE$, d'où l'existence de cette dernière, et la commutation à $f$, puisque $\varinjlim s_i = s$ ( $\varinjlim$ d'un système inductif filtrant constant).

c) Il faut prouver que si $x \in \Ob\cE_s$ est $\pi$-accessible dans $\cE_s$, il l'est dans $\cE$, i.e. que pour tout système inductif $(y_i)_{i\in I}$ grande devant $\pi$, la flèche canonique
\[
\varinjlim_{i} \Hom(x,y_i) \;\longrightarrow\; \Hom(x,y) \tag*{$(**)$}
\]
est bijective. Or comme par 2°) $\pi \geqslant \pi_1 \geqslant \sup(\card S, \sup_\lambda \pi_\lambda)$, il résulte encore du lemme % [sic] card S -> card \mathcal{C}, no S is defined; the condition 2° bounds card \mathcal{C}
que, quitte à remplacer $I$ par une sous-catégorie cofinale, on peut supposer que le système inductif $(y_i)_{i\in I}$ provient d'un système dans une catégorie fibre $\cE_t$, et que $y$ est sa $\varinjlim$ dans $\cE_t$. Mais alors les ensembles $\Hom(x,y)$ et $\Hom(x,y_i)$, et les morphismes de transition entre celles-ci, sont encore dans $\Ens/\Hom(s,t)$, et on est ramené
% foot: Version 17 décembre 2025

% ---- from der-ch18_p166-170.tex
\printedpage{164}
% head: A. Grothendieck : Dérivateurs, Ch. XVIII, Édition des Universités de Montpellier II et Paris VII

\msp{page 430}à prouver que pour tout $\lambda : s \longrightarrow t$, la flèche
\[
\varinjlim_{i} \Hom_{\lambda}(x, y_i) \ \longrightarrow \ \Hom_{\lambda}(x, y)
\]
est bijective. Pour $x$ et $\lambda$ fixés, et des systèmes inductifs variables dans $\cE_t$, grands devant $\pi$, cela signifie que le foncteur
\[
H_{\lambda}(x) \ = \ (y \longmapsto \Hom_{\lambda}(x, y)) \ : \ \cE_t \ \longrightarrow \ \Ens
\]
est $\pi$-accessible, et ceci quel que soit
\[
x \ \in \ \Ob \cE_{s\,\pi} \ .
\]
Or c'est le cas si $s \in \Ob \cE_{s\,\pi_0}$, en vertu de l'hypothèse $3^{\circ}$) (alors $H_{\lambda}(x_0)$ est $\pi_1$-accessible, % [sic] $s \in \Ob \cE_{s\,\pi_0}$ -> $x_0 \in \Ob \cE_{s\,\pi_0}$, the condition concerns an object of $\cE_{s\,\pi_0}$
et a fortiori $\pi$-accessible). Il faut passer de là au cas où $x \in \Ob \cE_{s\,\pi}$. Or comme $\pi_0$ est utile pour $\cE_s$ par $1^{\circ}$, on a
\[
x \ = \ \varinjlim_{j \in J} x_j \ , \qquad x_j \in \Ob \cE_{s\,\pi_0},\ J \text{ grande devant } \pi_0 \ ,
\]
et comme par $2^{\circ}$, $\pi^{\pi_0} = \pi \geqslant \pi_1 \geqslant \sup(\pi_0, \card \red \cE_{s\,\pi_0})$, on en conclut qu'on peut prendre $J$ de sorte que
\[
\card J \ \leqslant \ \pi \ .
\]
D'autre part, en vertu de $\pi \geqslant \pi_1 \ \geqslant \ \pi_{\lambda}$, le foncteur
\[
\left\{
\begin{array}{lcl}
x_0 & \longmapsto & H_{\lambda}(x_0) \\[1ex]
\cE_s & \longrightarrow & (\underline{\Hom}(\cE_s, \Ens))^{\circ}
\end{array}
\right.
\]
est $\pi_0$-accessible, donc on a
\[
H_{\lambda}(x_0) \ \simeq \ \varprojlim_{j \in J} H_{\lambda}(x_j) \ .
\]

\msp{page 431}Comme $x_j \ \in \ \Ob \cE_{s\,\pi_0}$, $H(x_j)$ est un foncteur $\pi_1$-accessible, et comme $\card J \leqslant \pi$, la $\varprojlim\limits_{j \in J} H_{\lambda}(x_j)$ est aussi $\pi$-accessible (cf. SGA 4 I 9.8).

\dgrblock{dgrb1115}{The numerical rank assertion needs another hypothesis}
\begin{dcomment}{development}{p430}{The numerical rank assertion needs another hypothesis}
Grothendieck wants to pass from continuity on a small family of generators to continuity on every object of the chosen larger rank. This is a sound strategy, but the displayed implication on this page uses first-variable continuity at $\pi_0$ although H2 only provides it at $\pi_\lambda$. The inequality $\pi\geq\pi_1\geq\pi_\lambda$ does not repair that mismatch. The following example shows that the numerical assertion of Corollary~4.15.2(c), as stated, is false.

\textbf{A bounded-powerset counterexample.}\ours\quad
Put $\pi_0=\aleph_0$ and $\pi_1=\pi=2^{\aleph_0}$. Let both fibres over $\Delta_1$ be copies of $\mathbf{Set}$. For a set $A$, write $P_{\leq\pi}(A)$ for the set of its subsets of cardinal at most $\pi$, with direct image as its action on maps. Define the total category by the usual hom sets within each fibre, no arrows from the second fibre to the first, and
\[
 \operatorname{Hom}_{0\to1}(A,B)
       =\mathbf{Set}(P_{\leq\pi}(A),B).
\]
Functoriality of bounded powerset defines both compositions involving a crossing arrow, so this is a category over $\Delta_1$.

We check every numerical assumption. Both fibres have useful cardinal $\aleph_0$ in Grothendieck's convention: every set is the $\aleph_1$-filtered union of its countable subsets. A skeleton of the category of countable sets has $2^{\aleph_0}$ morphisms, so the reduced-size bound holds. Cardinal arithmetic gives $\pi^{\pi_0}=\pi$.

The functor $P_{\leq\pi}$ preserves $\pi^+$-filtered colimits. Indeed, a subset of the colimit of size at most $\pi$ can be lifted by choosing at most $\pi$ representatives and one common stage. If two lifted subsets have the same image, choose matching elements and witnesses for their equalities; at most $\pi$ witnesses suffice, and a common upper stage makes the two subsets equal. These two observations prove surjectivity and injectivity of the canonical comparison. Consequently H2 holds with $\pi_\lambda=\pi$ for the crossing arrow. H3 holds for every $A$, since the mixed hom functor is representable by the set $P_{\leq\pi}(A)$. If $A$ is countable, this representing set has cardinal at most $2^{\aleph_0}=\pi$, so the stronger H3 bound demanded in part~(c) also holds. The identity-arrow conditions are automatic.

Now choose $A$ of cardinal $\pi$ in the first fibre. It is $\pi$-accessible there. But $P_{\leq\pi}(A)=P(A)$ has cardinal $2^\pi>\pi$. Index a diagram in the second fibre by all subsets $B\subseteq P(A)$ with $|B|\leq\pi$, ordered by inclusion. This poset is $\pi^+$-filtered, its union is $P(A)$, and this is also its total colimit. The identity map of $P(A)$ belongs to $\operatorname{Hom}_{0\to1}(A,P(A))$ and factors through no such $B$. Therefore $A$ is not $\pi$-accessible in the total category. This disproves both the inclusion in part~(c) and the equality in part~(d).\par

The repair is simple and preserves Proposition~4.15.1. Choose $\pi_0$ to dominate all $\pi_\lambda$ before making the subsequent choices of $\pi_1$ and $\pi$. Then the $\pi_0$-filtered generator presentation can be carried through $H_\lambda(-,y)$, and the limit--colimit calculation on manuscript page~419 applies. Such a common initial choice is available because the base is small. We checked the standard uniformization and two-categorical closure results in the sources cited here; we did not find this particular bounded-powerset diagnosis of the printed rank statement. The example is supplied to isolate the missing hypothesis, without a claim of priority.
\end{dcomment}


Ainsi on a prouvé $\cE_{s\,\pi} \subset \cE(\pi)$ pour tout $s \in \Ob \cC$, i.e. l'inclusion
\[
\cE(\pi) \ \subset \ \cE_{\pi} \ ,
\]
d'où le foncteur pleinement fidèle
\[
\Ind_{\pi}(\cE(\pi)) \longrightarrow \cE \ , \tag*{$(*)$}
\]
comme $\cE$ satisfait à $\mathrm{L}_{\pi}$. Il reste à prouver que c'est essentiellement surjectif. Mais si $x \in \Ob \cE$, $s = f(x)$, on a $x \in \Ob \cE_s$, et comme $\pi_0$ est utile pour $\cE_s$ par $1^{\circ}$), et $\pi^{\pi_0} = \pi > \pi_0$,
$\pi$ est également utile pour $\pi$, donc on a \textit{dans} $\cE_s$ % [sic] utile pour $\pi$ -> utile pour $\cE_s$, usefulness is a property relative to the category $\cE_s$
\[
x \ = \ \varinjlim_{I} x_i \ , \qquad \text{les } x_i \in \cE_{s\,\pi} \subset \cE(\pi),\ I \text{ grande devant } \pi \ ,
\]
\printedpage{165}
% head: A. Grothendieck : Dérivateurs, Ch. XVIII, Édition des Universités de Montpellier II et Paris VII
et d'après a) (comme $\pi \geqslant \pi_1 \geqslant \sup_{\lambda \in \Fl \cC} \pi_{\lambda}$) $x$ est aussi limite des $x_i$ dans $\cE$, donc il est dans l'image essentielle de $(*)$.

d) Reste à prouver que l'on a $\cE_{\pi} \ \subset \ \cE(\pi)$. Mais on sait (par 4.8.~\dots), comme $\cE \ \simeq \ \Ind_{\pi}(\cE(\pi))$, que $\cE_{\pi} \ = \ \cE(\pi)^{\natural}$, i.e. que pour $x \ \in \ \Ob \cE_{\pi}$, $\exists \ y \ \in \ \cE(\pi)$ et un projecteur $p$ dans $x$, tel que $x$ soit isomorphe à l'image de $p$ dans $\cE$. Or $f(p)$ est un projecteur dans $f(y) = s$, donc par hypothèse sur $\cC$, $f(p) = \id_s$, i.e. $p$ est un projecteur dans $\cE_s$. Comme donc $x$ est dans $\cE_s$, et

\msp{page 432}même dans $\cE_{s\,\pi}$, puisque $\cE_{s\,\pi}$ est stable dans $\cE_s$ par facteurs directs. Cela prouve que $x \in \Ob \cE(\pi)$, q.e.d.

On peut présenter 4.15.1 sous une forme plus sympathique.

\textbf{Corollaire 4.15.4.} \textit{Soit $f : \cE \longrightarrow \cC$ un foncteur de $\mathfrak{U}$-catégories. Les conditions suivantes sont équivalentes :}
\begin{enumerate}
\item[\textit{a)}] \textit{Pour tout foncteur $\cC' \longrightarrow \cC$, avec $\cC'$ petite et sans projecteurs, $\cE' \overset{\text{déf}}{=} \cE \times_{\cC} \cC'$ est accessible, et $f' : \cE' \longrightarrow \cC'$ est accessible.}
\item[\textit{b)}] \textit{Comme a), mais avec $\cC' = \mathbf{\Delta}_1$, en d'autres termes, c'est la condition $H$ p.~409 : pour toute $\lambda \in \Fl \cC$, $\cE_{\lambda} = \cE \times_{\cC} (\mathbf{\Delta}_1, i_{\lambda})$ est accessible, et $\cE_{\lambda} \longrightarrow \mathbf{\Delta}_1$ est accessible.}
\item[\textit{c)}] \textit{La condition $H' = H\ \mathit{1} + H\ \mathit{2} + H\ \mathit{3}$, explicitée dans l'énoncé de 4.15.1.}
\end{enumerate}
\textsc{Démonstration}. Elle se fait circulaire,
\[
\mathrm{a} \ \Longrightarrow \ \mathrm{b} \ \Longrightarrow \ \mathrm{c} \ \Longrightarrow \ \mathrm{a} \ ,
\]
l'implication a $\Longrightarrow$ b étant tautologique, b $\Longrightarrow$ c l'étant presque (cf. 4.14.24, a $\Longrightarrow$ b $\Longrightarrow$ c). En fait, les conditions b et c sont des conditions sur les catégories $\cE_s \xrightarrow{f_{\lambda}} \mathbf{\Delta}_1$ pour $\lambda \in \Fl \cC$, donc stables par changement de base. Pour prouver

\msp{page 433}c $\Longrightarrow$ a, on peut donc supposer que $\cC' = \cC$, $\cC' \longrightarrow \cC$ le foncteur identique, et prouver que si c) est satisfaite et si $\cC$ est petite et sans projecteurs, alors $\cE$ est accessible et $f$ accessible. Or c'est ce qu'affirme 4.15.1.

\dgrblock{dgrb1116}{Recovering the generators without choosing a splitting}
\begin{dcomment}{exposition}{p432}{Recovering the generators without choosing a splitting}
The conclusion Grothendieck needs is that the small generators of the total category can be recognized in the fibres. After correcting the uniformity assumption, the difficult inclusion is $\mathcal E(\pi)\subseteq\mathcal E_\pi$. The reverse inclusion has a short proof which also avoids a strict-fibre ambiguity in the printed retract argument.

Suppose $x\in\mathcal E_s$ is $\pi$-accessible in the total category. Let $y=\mathop{\rm colim}_i y_i$ be an $L_\pi$-colimit inside $\mathcal E_s$. The inclusion preserves it by part~(a), so
\[
 \mathop{\rm colim}_i\mathcal E(x,y_i)
       \xrightarrow{\sim}\mathcal E(x,y).
\]
Both sides map to the fixed set $\mathcal C(s,s)$, and the transition maps preserve that map. Taking the fibre over $1_s$ commutes with the colimit: a representative has base arrow $1_s$ precisely when its image does, because every transition lies over the identity. We obtain the required bijection for $\mathcal E_s(x,-)$. Thus $x$ is already $\pi$-accessible in its own fibre.

In the printed proof the projector must be on the chosen object $y$, not on $x$. Moreover, a splitting of that vertical projector in the total category may have an image over an object merely isomorphic to $s$. One can replace it by the splitting in the accessible fibre, but the identity-component argument proves the assertion directly and works without selecting an isomorphic representative.

Corollary~4.15.4 then expresses the result in a useful local form. H$'$ is unchanged by base change, because each new crossing hom functor is one of the original ones. Testing the walking arrow detects exactly this data, and the corrected existential theorem assembles it over any small base without nonidentity idempotents. The restriction on the testing bases is part of the assertion; it does not license replacement of an arbitrary strict pullback by a pseudopullback.
\end{dcomment}


\textbf{Remarque 4.15.5.} Si $f : \cE \longrightarrow \cC$ est un foncteur pleinement fidèle, il est immédiat que la condition H est satisfaite (car $\cE_{\lambda} \longrightarrow \mathbf{\Delta}_1$ sera pleinement fidèle, donc $\cE_{\lambda}$ est équivalente
à $e$ où à $\mathbf{\Delta}_1$, donc accessible, et tout foncteur de $\cE_{\lambda}$ dans une catégorie quelconque est % [sic] où -> ou, disjunction "e ou Δ1"
accessible \dots). Mais si on suppose $\cC$ petite \textit{sans} supposer $\cC$ sans projecteurs, il ne s'ensuit pas que $\cE$ soit accessible, ou ce qui revient au même (étant essentiellement petite), que $\cE$ soit karoubienne. (L'exemple le plus trivial est celui où $\cE = \cC$, $f = \id_{\cC}$, $\cE$ n'est karoubienne que si $\cC$ l'est !) Mais même en supposant $\cC$ une petite catégorie karoubienne, il ne s'ensuit pas forcément que $\cE$ l'est. En effet, \textit{toute} petite catégorie $\cE$ (karoubienne ou non) se plonge comme sous-catégorie pleine dans son enveloppe karoubienne $\cC$ -- il suffit donc de prendre $\cE$ non karoubienne. L'exemple

\msp{page 434}le plus économique est
\[
\Ob \cE \ = \ \{0\} \ , \qquad \Fl(\cE) \ = \ \{\id_0, p\} \ , \qquad \text{avec } p^2 = p \neq \id_0 \ ,
\]
\printedpage{166}
% head: A. Grothendieck : Dérivateurs, Ch. XVIII, Édition des Universités de Montpellier II et Paris VII
l'enveloppe de Karoubi de $\cE$ est donnée par
\[
P \ = \ \left\{
\begin{tikzcd}
1 \arrow[r, shift right, "\sigma"'] & 0 \arrow[l, shift right, "\rho"'] \arrow[loop right, "p"]
\end{tikzcd}
\qquad
\begin{aligned}
\rho\sigma \ &= \ \id_1 \\
\sigma\rho \ &= \ p \ \neq \ \id_0
\end{aligned}
\quad \Longrightarrow \ p^2 = p
\right\} .
\]
Cela soulève la question de trouver un énoncé similaire à 4.15.1, donnant des conditions serviables simples, plus fortes que $\mathrm{H}' = \mathrm{H}$ (1+2+3), sur $f : \cE \longrightarrow \cC$, pour que $\cE$ soit accessible et $f$ accessible, en supposant $\cC$ \textit{karoubienne}. Ou encore, des conditions nécessaires et
suffisants serviables (plus fortes que 4.15.4) sur $f : \cE \longrightarrow \cC$, pour que toute \textit{petite} catégorie % [sic] suffisants -> suffisantes, agreement with "conditions"
\textit{karoubienne} $\cC'$, et tout foncteur $\cC' \longrightarrow \cC$, la catégorie $\cE' = \cE \times_{\cC} \cC'$ soit accessible, et le foncteur $\cE' \longrightarrow \cC'$ accessible. Une idée naturelle, c'est qu'il suffirait que cette conclusion soit vraie quand $\cC'$ est une des deux catégories $\mathbf{\Delta}_1$, $P$ (où $P$ est la catégorie à deux objets, à trois flèches non identiques, définie plus haut). Mais je doute qu'il en soit ainsi (\footnote[178]{Ça devrait marcher par contre si on prend les catégories $0 \longrightarrow 1$,
\[
\begin{array}{c}
2 \longrightarrow \underbrace{\begin{tikzcd}[ampersand replacement=\&]
1 \arrow[r, shift right, "\sigma"'] \& 0 \arrow[l, shift right, "\rho"'] \arrow[loop right, "p"]
\end{tikzcd}}_{P}
\quad \text{et} \quad
2 \longleftarrow \underbrace{\begin{tikzcd}[ampersand replacement=\&]
1 \arrow[r, shift right, "\sigma"'] \& 0 \arrow[l, shift right, "\rho"'] \arrow[loop right, "p"]
\end{tikzcd}}_{P}
\qquad (\text{\footnotemark[179]}) \ , \\
\nwarrow \ \begin{array}{c} \text{3 objets et} \\ \text{5 flèches non} \\ \text{identiques} \end{array} \ \nearrow
\end{array}
\]
grâce à 4.15.7, page 440.}\footnotetext[179]{ces catégories représentent respectivement les foncteurs $\mathrm{Cat} \longrightarrow \Ens$ associant à une catégorie $\cC$ les diagrammes dans $\cC$ formés de $(x, y, p, u)$, avec $p$ projecteur dans $p$ \texttt{[plutôt dans $\cC$]}, et $u : y \longrightarrow x^{\natural}$ (i.e. $u : y \longrightarrow x$ avec $u = pu$), resp. $u : x^{\natural} \longrightarrow y$ (i.e. $u : x \longrightarrow y$ avec $u = up$).}).

On va quand même essayer un énoncé un peu très particulier dans ce sens :

% Verified on PDF leaf8: a non-textual wavy marginal mark beside items a), b); no mathematical token omitted.
\msp{page 435}\textbf{Proposition 4.15.6} (\footnote[180]{\textit{Faux ?}}). \textit{Soit $f : \cE \longrightarrow \cC$ un foncteur de $\mathfrak{U}$-catégories. On suppose :}
\begin{enumerate}
\item[\textit{a)}] \textit{$f$ est fibrant ou cofibrant} (\footnote[181]{Le cas fibrant n'est pas prouvé, peut-être \textit{faux} -- et de même pour le cas cofibrant !}).
\item[\textit{b)}] \textit{Les catégories fibres $\cE_s$ ($s \in \Ob \cC$) sont accessibles, et pour toute flèche $\lambda : s \longrightarrow t$ dans $\cC$, le foncteur changement de base $\lambda^* : \cE_t \longrightarrow \cE_s$ (resp. cochangement de base $\lambda_* : \cE_s \longrightarrow \cE_t$) est accessible.}
\item[\textit{c)}] \textit{$\cC$ est essentiellement petite et karoubienne.}
\end{enumerate}
\textit{Alors $\cE$ est accessible, et $f$ est accessible.}

\dgrblock{dgrb1117}{Why a retract has to be lifted}
\begin{dcomment}{examples}{p434}{Why a retract has to be lifted}
Grothendieck's smallest example makes visible what the walking arrow does not test. Let $B$ be the one-object category with $p^2=p\ne1$. Its Karoubi envelope $P$ has objects $0,1$, arrows $\sigma:1\to0$ and $\rho:0\to1$, and relations
\[
 \rho\sigma=1_1,\qquad\sigma\rho=p.
\]
The inclusion $B\hookrightarrow P$ is fully faithful. Every pullback along a walking arrow has total category equivalent to an empty category, a point, or $\Delta_1$, and therefore satisfies the local accessibility condition H. Nevertheless $B$ is not accessible: the long constant-object diagram with transition $p$ has no colimit, as computed above. The base $P$ is Karoubian, but its retract object has not been lifted into $B$.

This is why the next proposed test categories attach an incoming or outgoing arrow to a retract. The retract alone tests existence of an image; an incoming arrow tests maps into that image, and an outgoing arrow tests maps from it. Concretely, once a splitting $x\xrightarrow{r}z\xrightarrow{i}x$ with $ri=1$ and $ir=p$ is fixed, an arrow $y\to z$ corresponds to an arrow $u:y\to x$ satisfying $pu=u$, by $v\mapsto iv$ and $u\mapsto ru$. The two composites are identities because $ri=1$ and $pu=u$. Dually, arrows $z\to y$ correspond to arrows $u:x\to y$ satisfying $up=u$. These are the two three-object diagrams in footnote~178. Their later reformulation in \S4.16.6 separates the global accessibility problem from the stronger requirement concerning all strict base changes.

The counterexample also explains why a (co)fibration is a promising extra assumption. Cartesian or cocartesian transport along the two arrows of a retract gives a way to move the problem into the fibre over its image. The proof below makes that move explicit; it resolves the doubts attached to Proposition~4.15.6 without pretending that its written proof is complete.
\end{dcomment}


\textsc{Démonstration}. On a déjà vu que les conditions a) et b) impliquent $\mathrm{H}'$, i.e. H (1+2+3). D'ailleurs, si $\cC'$ est une catégorie et $\cC' \xrightarrow{\ \sim\ } \cC$ une équivalence, donc alors $\cE' \overset{\text{déf}}{=} \cE \times_{\cC} \cC' \longrightarrow \cE$ est une équivalence, donc la conclusion voulue pour $\cE$, $f$ équivaut à la même conclusion pour $\cE'$, $f' : \cE' \longrightarrow \cC'$. Donc on peut supposer que $\cC$ est petite (et même réduite, si on y tient). En vertu de 4.15.2 a), on voit donc que les foncteurs d'inclusion $\cE_s \longrightarrow \cE$ sont accessibles (et même $\pi_1$-accessibles) dès que $\pi_1 \geqslant \sup_{\lambda \in \Fl \cC} \pi_{\lambda}$. On voit d'ailleurs

\msp{page 436}par AQT \texttt{[âne qui trotte]} qu'on peut trouver des cardinaux $\pi_0$, $\pi_1$ ayant les propriétés suivantes :

\printedpage{167}
% head: A. Grothendieck : Dérivateurs, Ch. XVIII, Édition des Universités de Montpellier II et Paris VII
\begin{enumerate}
\item[\circled{1\textsuperscript{o}}] $\pi_0$ est utile pour $\cE_s$, $\forall \ s \in \Ob \cC$.
\item[$2^{\circ}$)] $\pi_1 \geqslant \sup(\pi_0, \card \cC, \sup_{\lambda \in \Fl \cC} \pi_{\lambda}, \sup_{s \in \Ob \cC} \card \red \cE_s)$.
\item[\circled{3\textsuperscript{o}}] $\forall \ (\lambda : s \longrightarrow t) \in \Fl \cC$ et $x_0 \in \Ob \cE_{s\,\pi_0}$, le foncteur
\[
H_{\lambda}(x_0) \ : \ (x_1 \longmapsto \Hom_{\lambda}(x_0, x_1)) \ : \ \cE_t \ \longrightarrow \ (\Ens)
\]
est $\pi_1$-accessible. (\textbf{N.B.} Si $f$ est cofibrant, cela signifie que $\forall \ \lambda \in \Fl \cC$, $\lambda : s \longrightarrow t$,
\fbox{$\lambda_*(\cE_{s\,\pi_0}) \ \subset \ \cE_{s\,\pi_1}$}. Si $f$ est fibrant, il suffit d'exiger que $\forall \ \lambda \in \Fl \cC$, $\lambda^* : \cE_t \longrightarrow \cE_s$ % [sic] $\cE_{s\,\pi_1}$ -> $\cE_{t\,\pi_1}$, $\lambda_*$ goes to $\cE_t$, cf. the later display "(contenu dans 3°)"
soit $\pi_1$-accessible.) Mais dans ce cas, on veut aussi \fbox{$\lambda^*(\cE_{s\,\pi_0}) \ \subset \ \cE_{s\,\pi_1}$}. % [sic] $\lambda^*(\cE_{s\,\pi_0})$ -> $\lambda^*(\cE_{t\,\pi_0})$, $\lambda^*$ is defined on $\cE_t$, cf. the later display "(contenu dans 3°)"
\item[\circled{4\textsuperscript{o}}] Les foncteurs $\lambda^*$ (cas $f$ fibrant) resp. $\lambda_*$ (cas $f$ cofibrant) sont $\pi_1$-accessibles.
\end{enumerate}
Soit alors $\pi$ un cardinal tel que
\[
\pi \ = \ \pi^{\pi_0} \ \geqslant \ \pi_1 \ , \tag*{$(*)$}
\]
et soit $\cE(\pi)$ la sous-catégorie (strictement) pleine de $\cE$ définie par
\[
\Ob \cE(\pi) \ = \ \bigcup_{s \, \in \, \Ob \cC} \Ob \cE_{s\,\pi} \ .
\]
Grâce à $(*)$, et à $1^{\circ}$, $4^{\circ}$ et à
\[
\left\{
\begin{array}{l}
\pi_1 \ \geqslant \ \sup_s(\card \red \cE_{s\,\pi_0}) \qquad \text{(contenu dans $2^{\circ}$)} \\[1ex]
\lambda^*(\cE_{t\,\pi_0}) \ \subset \ \cE_{s\,\pi_1} \qquad \text{resp.} \qquad \lambda_*(\cE_{s\,\pi_0}) \ \subset \ \cE_{t\,\pi_1} \qquad \text{(contenu dans $3^{\circ}$)} \ ,
\end{array}
\right.
\]
on voit que $\cE(\pi)$ est une sous-catégorie

\msp{page 437}fibrée, resp. une sous-catégorie cofibrée, sur $\cC$ (le foncteur $\cE(\pi) \longrightarrow \cE$ étant donc cartésien resp. cocartésien). À $\cC$-équivalence près, $\cE$ se reconstitue alors à partir de $\cE(\pi)$, en passant au $\Ind_{\pi}(\cE_{s\,\pi})$, qui forment un système fibré resp. cofibré sur $\cC$ par les foncteurs $\Ind_{\pi}(\lambda^*_{\pi})$ resp. $\Ind_{\pi}(\lambda^{\pi}_*)$ \dots

Prouvons que $\cE$ satisfait à $\mathrm{L}_{\pi}$ et que $f : \cE \longrightarrow \cC$ est $\pi$-accessible (\footnote[182]{Ici on utilise seulement $\pi \geqslant \card \red \cC$.}). Comme $\cC$ est karoubienne, et $\pi \geqslant \pi_1 \geqslant \card \cC$, $\cC$ satisfait à $\mathrm{L}_{\pi}$ et les ind-objets grands devant $\pi$ de $\cC$ sont essentiellement constants (4.8.13). Soit alors $I$ grande devant $\pi$, et $I \longrightarrow \cE$, $i \longmapsto x_i$, un foncteur, il faut construire une $\varinjlim$ dans $\cE$, en utilisant la $\varinjlim$ déjà connue $s$ des $s_i = f(x_i)$. On va prendre un $x$ dans $\cE_s$. Prenons d'abord le cas où $f$ est cofibrant. On voit que les $\beta_i(x_i)$ forment un système inductif dans $\cE_s$, puisque pour $i = j$, on a
\[
\beta_{i\,*}(x_j) \ = \ \beta_{j\,*} \underbrace{\beta_{ji\,*}(x_i)}_{\longrightarrow \, x_j} \ \longrightarrow \ \beta_{j\,*}(x_j) \ . % [sic] $\beta_{i\,*}(x_j)$ -> $\beta_{i\,*}(x_i)$, the object transported from $s_i$ is $x_i$
\]
\[
\begin{tikzcd}[row sep=small]
x_i \arrow[rr, "\alpha_i"] & & x \, ? \\
s_i \arrow[rr, "\beta_i"] \arrow[dr, "\beta_{ji}"'] & & s \\
& s_j \arrow[ur, "\beta_j"'] &
\end{tikzcd}
\]
On pose

\printedpage{168}
% head: A. Grothendieck : Dérivateurs, Ch. XVIII, Édition des Universités de Montpellier II et Paris VII
\msp{page 438}
\[
x \ = \ \varinjlim_{I} \beta_{i\,*}(x_i) \ , \qquad \varinjlim \ \textit{dans } \cE_s \ .
\]
On sait que c'est une $\varinjlim$ dans $\cE$, donc que pour tout $y$ dans $\cE$, on a
\[
\Hom(x, y) \xrightarrow{\ \sim\ } \varprojlim_{I} \Hom(\beta_{i\,*}(x_i), y) \ ,
\]
ou ce qui revient au même, pour toute flèche $\lambda : s \longrightarrow t$, et $y$ dans $\Ob \cE_t$, on a
\[
\underbrace{\Hom_{\lambda}(x, y)}_{= \, \Hom_{\cE_t}(\lambda_* x, y)} \ \simeq \ \varprojlim \Hom_{\cE_t}(\underbrace{\lambda_* \beta_{i\,*}(x_i)}_{(\lambda_i \beta_i)_*(x_i)}, y) \ .
\]
Cela résulte aussi du fait que $\lambda_*$ est $\pi$-accessible (étant $\pi_1$-accessible) \dots{} On vérifie facilement que c'est une $\varinjlim$ des $x_i$ dans $\cE$. Et bien sûr, cette limite commute avec $f$ par construction.

Mais dans le cas où $f$ est fibrant, je ne vois comment faire pour construire une limite dans $\cE_s$. Il est vrai qu'on peut supposer que tous les $s_i$ sont égaux à $s$, mais les transitions $\beta_{ji}$ ne sont pas l'identité, ni même des isomorphismes. Il faudrait prendre $\cC = P$ pour essayer de comprendre -- je ne vais pas insister (\footnote[183]{J'ai vérifié que si $f$ est fibrant, alors $\cE$ est \textit{karoubienne}. Il y a donc de l'espoir \dots}).

\dgrblock{dgrb1118}{Constructing the missing colimit in a fibration}
\begin{dcomment}{exposition}{p438:Il faudrait prendre $\cC = P$ pour essayer de comprendre}{Constructing the missing colimit in a fibration}
Grothendieck has already constructed total colimits in the opfibration case. For a fibration the direction of transport initially seems wrong, but the splitting of the stabilized base idempotent supplies exactly the missing direction. We give a direct proof, using the constant-idempotent reindexing established later in \S4.15.7 and the sufficient strict bound in Corollary~4.15.11. This makes the proof independent of an unresolved optimal cardinal estimate.

Replace the fibration by an equivalent split fibration and choose a sufficiently large regular successor $\kappa$ so that every fibre is $\kappa$-accessible, every reindexing functor preserves $\kappa$-filtered colimits, and the small-base stabilization theorem applies. This simultaneous choice uses only a small family of accessible functors. Reindex a given $\kappa$-filtered diagram cofinally by a $\kappa$-directed poset so that every base object is $s$ and every strict transition is one idempotent $p$. If this poset has a greatest element, its value is the total colimit and the base colimit, and every hom comparison is immediate. Otherwise every index has a strictly larger index, since the poset is directed. Work in this latter case and split
\[
 s\xrightarrow{r}t\xrightarrow{i}s,
 \qquad ri=1_t,\quad ir=p.
\]
Write $U=i^*:\mathcal E_s\to\mathcal E_t$ and $R=r^*:\mathcal E_t\to\mathcal E_s$, so $UR=1$ and $RU=p^*$. A strict transition $X_j\to X_k$ is represented by
\[
 a_{kj}:X_j\longrightarrow RU X_k.
\]
Composition says $a_{\ell j}=RU(a_{\ell k})a_{kj}$ for $j<k<\ell$. Therefore $Z_j=UX_j$, with transition $z_{kj}=U(a_{kj})$, is a genuine diagram in $\mathcal E_t$. Form $Z=\mathop{\rm colim}_j Z_j$ there and denote its structural maps by $c_j$.

For $k>j$, the composite $R(c_k)a_{kj}:X_j\to RZ$ is independent of $k$: two choices admit a common later index, where the composition identity and $c_\ell z_{\ell k}=c_k$ give equality. These composites define a cocone from the original diagram to the total object $Z$, lying over $r$.

To prove universality, take a cocone to $Y\in\mathcal E_v$. Its base maps $\lambda_j:s\to v$ satisfy $\lambda_j=\lambda_kp$ for $j<k$. Choosing one further index shows that they are all equal to a map $\lambda$ with $\lambda p=\lambda$. Hence $\lambda=\mu r$ for the unique $\mu=\lambda i:t\to v$. Its legs are maps $b_j:X_j\to R\mu^*Y$. Applying $U$ gives a compatible family $Ub_j:Z_j\to\mu^*Y$. Conversely that compatible family recovers $b_j$ from any later index by
\[
 b_j=R(Ub_k)a_{kj}\qquad(k>j).
\]
These operations are inverse, by the cocone equation and $UR=1$. Fibre cocones on $Z_j$ correspond to maps $Z\to\mu^*Y$, so the original cocones correspond precisely to total maps $Z\to Y$. This proves both existence of the total colimit and its preservation by the projection.

The same transport also proves the Karoubian assertion in footnote~183. An idempotent on $X$ over $p$ is represented by $a:X\to RUX$ with $RU(a)a=a$. Then $b=Ua$ is a vertical idempotent on $UX$. Split $b=jq$, $qj=1_Y$ in the accessible fibre $\mathcal E_t$. The maps represented by $Rq\,a:X\to RY$ and $j:Y\to UX$ give a splitting in the total category: their composites are $1_Y$ and the original idempotent. Thus the source's observation is correct, and transport explains it.
\end{dcomment}


\msp{page 439}Dans le cas où $f$ est fibrant, admettons (comme hypothèse supplémentaire au besoin) que $\cE$ satisfasse $L$ et que $f$ soit accessible, et prenons $\pi$ assez grand pour que $\cE$, $f$ satisfassent $\mathrm{L}_{\pi}$. Nous pouvons alors terminer la démonstration de 4.15.6. Il faut prouver
\[
\cE(\pi) \ \subset \ \cE_{\pi} \ ,
\]
donc que si $x \in \Ob \cE_{s\,\pi}$, et $(y_i)_{i \in I}$ est un système inductif dans $\cE$ grand devant $\pi$, de limite $y$, on a
\[
\varinjlim \Hom_{\cE}(x, y_i) \xrightarrow{\ \sim\ } \Hom_{\cE}(x, y) \ .
\]
Donc il faut prouver que toute flèche $x \xrightarrow{u} y$ se factorise, de façon essentiellement unique $u_i$, par un $y_i$. Mais la flèche $s \xrightarrow{\ \lambda \, = \, f(u)\ } t$ dans $\cC$ se factorise, de façon essentiellement unique, par des $\lambda_i \, : \, s \longrightarrow t_i$. Donc on est ramené à trouver un \textit{$\lambda_i$-morphisme} $x \longrightarrow y_i$ pour $i$ grand, qui factorise. Dans le cas fibrant, cela signifie qu'on cherche $x \longrightarrow \lambda^*_i(y_i)$ qui factorise $x \longrightarrow \lambda^*(y)$, or les $\lambda^*_i(y_i)$ forment un système inductif dans $\cE_s$, et sa limite \textit{dans} $\cE_s$ est $\lambda^*(y)$ (\textit{à vérifier}). Comme $x$ est dans $\cE_{s\,\pi}$, l'hom\texttt{[omorphisme]} donné $x \longrightarrow \lambda^*(y)$ se factorise de

\msp{page 440}façon essentiellement unique par un $\lambda^*_i(y_i)$.

\dgrblock{dgrb1119}{Presentability after splitting the base idempotent}
\begin{dcomment}{exposition}{p439}{Presentability after splitting the base idempotent}
The preceding construction identifies the colimit whose existence was assumed on this page. It also verifies the pullback comparison marked ``\`a v\'erifier'', and proves the presentability statement needed for accessibility.

Keep the notation $Z_j=UX_j$, $Z=\mathop{\rm colim}Z_j$ and the splitting $s\xrightarrow r t\xrightarrow i s$. Let $x\in\mathcal E_u$ be $\kappa$-presentable in its fibre. A total arrow $x\to Z$ over $\nu:u\to t$ is a fibre arrow
\[
 x\longrightarrow\nu^*Z
       \cong\mathop{\rm colim}_j\nu^*Z_j.
\]
The displayed isomorphism follows from the chosen accessibility rank of $\nu^*$. Fibre presentability factors this map through $\nu^*Z_j=(i\nu)^*X_j$, giving a total arrow $x\to X_j$ over $i\nu$. The cocone leg $X_j\to Z$ over $r$ has fibre representative $d_j=R(c_k)a_{kj}$ for any $k>j$, and
\[
 U(d_j)=c_kU(a_{kj})=c_kz_{kj}=c_j.
\]
Consequently composition with that leg applies $\nu^*c_j$ to the chosen factorization, giving exactly the prescribed arrow to $Z$; its base is $ri\nu=\nu$. This proves surjectivity of the total hom comparison.

For injectivity, take two representatives whose composites to $Z$ agree. Passing to a common index makes their base arrows have the same composite with $r$. Passing once more through a transition over $p=ir$ makes both base arrows equal to $i\nu$ for that common $\nu$. They can now be interpreted as maps into the one fibre diagram $\nu^*Z_j$. Their equality at the colimit is witnessed at a later stage by the injectivity of the fibre hom comparison. This is exactly equality of the original total representatives after that transition.

Every fibrewise $\kappa$-presentable object is therefore $\kappa$-presentable in $\mathcal E$. Fibre inclusions preserve the relevant colimits: for a diagram in $\mathcal E_u$, test its fibre colimit against $Y\in\mathcal E_v$ separately over every $\lambda:u\to v$ using $\mathcal E_u(-,\lambda^*Y)$. Since the base is small, the union of small generating families in the fibres is small, and it generates all total objects. The fibration case of Proposition~4.15.6 follows, including accessibility of its projection.

Makkai--Par\'e's general theory treats accessible indexed categories and their total categories in \S5.3 \citep{MakkaiPare1989}. The proof here records directly the special small Karoubian-base situation Grothendieck needs; we do not use an unverified identification of a strict pullback with a two-categorical limit. For a (co)fibration, however, pullback along an equivalence of bases is safe because isomorphisms lift. This explains the reduction to a small skeleton at the start of his argument.
\end{dcomment}


Cette fois par contre, c'est le cas $f$ cofibrant qui pose problème. En fait, je n'arrive pas à m'en tirer.

Ainsi, on trouve bien un foncteur
\[
\Ind_{\pi}(\cE(\pi)) \ \longrightarrow \ \cE \ ,
\]

% ---- from der-ch18_p171-175.tex
\setcounter{footnote}{183}

\printedpage{169}
% head: A. Grothendieck : Dérivateurs, Ch. XVIII, Édition des Universités de Montpellier II et Paris VII
qui est essentiellement surjectif, mais faute de savoir que $\mathcal{E}(\pi) \subset \mathcal{E}_\pi$ dans le cas cofibrant, il n'est pas clair que ce foncteur soit une équivalence de catégories.

\dgrblock{dgrb1120}{Completing the opfibration argument}
\begin{dcomment}{exposition}{p440:il n'est pas clair que ce foncteur soit une équivalence de catégories.}{Completing the opfibration argument}
The opfibration proof needs an additional transition after a fibre factorization. Choose a common strong rank $\kappa$ for the direct-image functors: they then preserve $\kappa$-filtered colimits and $\kappa$-presentable objects. This is the uniformization recalled by \citet[Introduction and proof of Theorem~2.1]{PareRosicky2013}. Reindex and split the base idempotent as above, and put
\[
 \begin{gathered}
 R=r_*:\mathcal E_s\to\mathcal E_t,
 \qquad U=i_*:\mathcal E_t\to\mathcal E_s,\\
 RU=1,\qquad UR=p_*.
 \end{gathered}
\]
A transition is represented by $a_{kj}:URX_j\to X_k$, with
$a_{\ell j}=a_{\ell k}UR(a_{kj})$ for $j<k<\ell$. The diagram $Z_j=RX_j$ has transitions $z_{kj}=R(a_{kj})$, and its fibre colimit $Z$ is the total colimit. To check this last assertion, base cocones again have one form $\mu r$. Maps out of $Z$ over $\mu$ are maps $\mu_*Z\to Y$; preservation of the fibre colimit by $\mu_*$ identifies them with exactly the cocones $\mu_*RX_j\to Y$ arising from the original diagram. This completes the universal-property verification sketched on page~438.

Let $x\in\mathcal E_u$ be fibrewise $\kappa$-presentable. A total map $x\to Z$ over $\nu:u\to t$ is a map $h:\nu_*x\to Z$. Since $\nu_*x$ is $\kappa$-presentable, it factors through some $Z_j=RX_j$, say by $h_j:\nu_*x\to RX_j$. Applying $U$ gives a map into $URX_j$, rather than into $X_j$. One later transition supplies the missing arrow:
\[
 (i\nu)_*x=U\nu_*x
       \xrightarrow{Uh_j}URX_j
       \xrightarrow{a_{kj}}X_k\qquad(k>j).
\]
Its composite to $Z$ is the desired total map. This proves surjectivity.

For injectivity, normalize two representatives to the same index and the same base arrow $i\nu$ by passing to a further transition. Write their fibre representatives as $b,b':U\nu_*x\to X_j$. Equality at $Z$ means equality of the two maps $Rb,Rb':\nu_*x\to RX_j$ in the colimit. Fibre presentability gives a stage $k$ where their composites with $z_{kj}$ agree. After one additional transition to $\ell>k$, the original representatives become
\[
 a_{\ell k}\,U(z_{kj}Rb)
 \quad\hbox{and}\quad
 a_{\ell k}\,U(z_{kj}Rb'),
\]
which are equal. The composition law above shows these are exactly their images in $X_\ell$. Thus the total hom comparison is bijective.

Fibre generation and preservation of its colimits by the inclusions now give accessibility of the total category, and the displayed Ind functor is an equivalence at the chosen rank. The doubts in footnotes~180--181 concern an unfinished proof, not a false conclusion. The abstract result is consistent with Makkai--Par\'e's Theorem~5.4.7, which computes lax colimits of strong diagrams of accessible categories by taking the idempotent completion of the categorical lax colimit; this formulation is explicitly recalled by \citet[Introduction]{PareRosicky2013}. We checked that attribution there rather than claiming to have inspected the unavailable book proof. The calculations above establish both orientations needed in this manuscript directly.
\end{dcomment}


\textbf{4.15.7.} Il me faut reporter la démonstration du lemme 4.15.3 (p. 428), en admettant 4.8.13. (Non encore démontré -- je vais revenir plus bas sur la démonstration de 4.8.13.) On va utiliser le

\textbf{Lemme 4.15.7.1.} \textit{Soient $\cC$ une catégorie, $\ux = (x_i)_{i\in I}$ un ind-objet de $\cC$. Pour que $\ux$ soit essentiellement constant, i.e. isomorphe, dans} $\Ind(\cC)$\textit{, au ind-objet constant défini par un objet}

\msp{page 441}
\textit{de $\cC$, il faut et il suffit qu'il existe un foncteur cofinal $\varphi : J \longrightarrow I$, avec $J$ filtrant, de telle façon que le système réindexé $(x_j)_{j\in J}$, avec $x_j = x_{\varphi(j)}$, ait la propriété suivante : pour tout $j \in J$, on peut trouver un projecteur $p_j$ dans $x_j$, de telle façon que :}
\begin{itemize}
\item[\textit{1°)}] \textit{Pour toute flèche $\lambda : j \longrightarrow j'$ dans $J$, on ait commutativité dans}
\[
\begin{tikzcd}
x_j \arrow[r, "x_\lambda"] \arrow[d, "p_j"] & x_{j'} \arrow[d, "p_{j'}"] \\
x_j \arrow[r, "x_\lambda"] & x_{j'}
\end{tikzcd}
\]
\textit{(où les $x_\lambda$ sont les flèches de transition), i.e. $p_{j'}x_\lambda = x_\lambda p_j$.}
\item[\textit{2°)}] \textit{Pour tout $j \in \Ob J$, il existe une flèche $\lambda : j \longrightarrow j'$ dans $J$ telle que $p_{j'}x_\lambda = x_\lambda p_j = x_\lambda$.}
\item[\textit{3°)}] \textit{Pour tout $j \in \Ob J$, le projecteur $p_j$ dans $x_j$ admet un image $x_j^\natural$.} % [sic] un image -> une image, gender agreement
\item[\textit{4°)}] \textit{Pour toute flèche $\lambda : j \longrightarrow j'$ dans $J$, la flèche $x_\lambda^\natural : x_j^\natural \longrightarrow x_{j'}^\natural$ induite par $x_\lambda$ (grâce à 1°) est un isomorphisme.}
\end{itemize}
\textit{De plus, si on peut supposer $J$ de la forme $i_0\backslash I$, le foncteur cofinal $J \longrightarrow I$ étant le foncteur canonique $i_0\backslash I \longrightarrow I$ ($i_0$ étant un objet fixé convenable de $I$)} \texttt{[phrase incomplète]}

\msp{page 442}
Notons tout de suite deux corollaires.

\dgrblock{dgrb1121}{The retract inside an essentially constant system}
\begin{dcomment}{exposition}{p441}{The retract inside an essentially constant system}
Grothendieck wants to recover an actual stage of a large diagram from the ind-object it represents. The obstruction is already visible in a finite calculation: the limiting object may be a retract of a stage without being the whole stage. The projectors in 4.15.7.1 identify precisely the part which survives. We write $h_x=\operatorname{Hom}_{\mathcal C}(-,x)$ and
\[
X=\varinjlim_{i\in I}h_{x_i}
\quad\hbox{in}\quad [\mathcal C^{\mathrm{op}},\mathrm{Ens}].
\]
\emph{Essentially constant} means that $X$ is representable. It does not mean that the displayed diagram was already constant, or that its colimit inside $\mathcal C$, if one exists, is represented by the same presheaf.

Here is the proof with the two inverse maps kept in their correct order. Suppose $X\simeq h_z$, and write
\[
\alpha:h_z\longrightarrow X,\qquad
\beta:X\longrightarrow h_z,\qquad
\beta\alpha=1_{h_z},\quad\alpha\beta=1_X.
\]
Yoneda gives a representative $\alpha_{i_0}:z\to x_{i_0}$ of $\alpha$, while $\beta$ gives compatible maps $\beta_i:x_i\to z$. The compatibility equation is
$\beta_j\tau_{ji}=\beta_i$, for $\tau_{ji}:x_i\to x_j$.
Several equations on manuscript pp.~449--450 reverse these composites; the reversals are printed so in the edition.

The upper slice $J=i_0\backslash I$ makes the chosen representative coherent. An object of $J$ is an arrow $u:i_0\to i$, and we put $\alpha_u=x(u)\alpha_{i_0}$. Then
\[
\beta_i\alpha_u=1_z,\qquad p_u=\alpha_u\beta_i,\qquad p_u^2=p_u.
\]
Thus $p_u$ splits with image $z$, and a morphism of the slice acts as the identity on that image. Its naturality follows from the compatibility of both families. Finally $\alpha\beta=1_X$ says that the class of $p_u$ equals the class of $1_{x_i}$ in the appropriate filtered colimit of Hom sets. Equality there has a witness: some arrow $u\to v$ satisfies $x(u\to v)p_u=x(u\to v)$. Naturality gives the other absorbed composite. These are exactly the four conditions of the lemma.

For completeness, the slice has the required size of filteredness. A small diagram in $i_0\backslash I$ becomes a diagram in $I$ after adjoining $i_0$ and its specified arrows. A cocone in $I$ therefore gives a cocone in the slice. The same construction shows that the projection is final. In particular, if $I$ admits cocones for diagrams of cardinal at most $\pi_0$, so does $J$.

Conversely, choose splittings $x_j\xrightarrow{r_j}z_j\xrightarrow{s_j}x_j$ of the stated projectors. Condition~1 makes the $z_j$ a diagram, and condition~4 makes its arrows isomorphisms. A filtered category has trivial groupoid completion: common upper objects compare two paths, and equalization of parallel arrows removes their ambiguity. Hence the $z_j$ are coherently isomorphic to one object $z$. Transporting $r_j$ and $s_j$ through this isomorphism produces maps $\beta:X\to h_z$ and $\alpha:h_z\to X$. The splitting identities give $\beta\alpha=1$, while condition~2 gives $\alpha\beta=1$ in the colimit. This also repairs the incomplete final sentence of the statement: one \emph{can} choose the upper slice described above.
\end{dcomment}


\textbf{Corollaire 4.15.7.2.} \textit{Supposons que $\cC$ soit sans projecteurs. Alors $\ux = (x_i)_{i\in I}$ est un ind-objet essentiellement constant si et seulement s'il existe un foncteur cofinal, $J \longrightarrow I$, avec $J$ filtrant (\footnote{et on peut supposer $J = i_0\backslash I \xrightarrow{\text{can.}} I$}), de telle façon que le système inductif réindexé $(x_j)_{j\in J}$ soit isomorphe} (dans la catégorie $\underline{\mathrm{Hom}}(J,\cC)$) \textit{à un système inductif} constant \textit{(i.e. défini par un foncteur constant $J \longrightarrow \cC$ : tous les $x_j$ égaux à un même $x$, et les morphismes de transition des identités).}

La suffisance est claire, pour la nécessité on applique 4.15.7.1, qui nous dit qu'on peut réindexer de façon à satisfaire aux conditions 1°) à 4°) (pour des projecteurs $p_j$ convenables dans les $x_j$). Mais l'hypothèse sur $\cC$ implique que $p_j = \id_{x_j}\ \forall\, j \in \Ob J$, et les conditions 1°, 2°, 3° sont satisfaites tautologiquement pour un tel choix, la condition 4° signifie que les morphismes de transition sont des isomorphismes. Mais il est immédiat que pour un système inductif $(x_j)_{j\in J}$ donné, avec $J$ 1-\textit{connexe} (p.ex. $J$ filtrant), la condition que les $x_j$ sont des isomorphismes % [sic] x_j -> x_\lambda, the transition morphisms are meant
est nécessaire et suffisante pour que ce système soit isomorphe à un système inductif constant, d'où la conclusion.
% foot: Version 17 décembre 2025

\printedpage{170}
% head: A. Grothendieck : Dérivateurs, Ch. XVIII, Édition des Universités de Montpellier II et Paris VII
\msp{page 443}
\textbf{Corollaire 4.15.7.3.} \textit{Supposons que $I$ soit grande devant $\pi_0$, où $\pi_0$ est un cardinal infini tel que $\pi_0 \geqslant \card \cC$. Alors $(x_i)_{i\in I}$ est essentiellement constant si et seulement si on peut trouver un foncteur cofinal $J \longrightarrow I$ avec $J$ grande devant $\pi_0$ (et ordonnée si on veut) (\footnote{La démonstration (légèrement réarrangée) montre qu'on peut prendre pour $J$ une sous-catégorie (pas pleine) d'une catégorie $i_1\backslash I$, pour un objet convenable $i_1$ de $I$. On a donc $\card J \leqslant \pi$ si $\card I \leqslant \pi$, $\pi$ étant un cardinal infini donné.}), de telle façon que le système réindexé $(x_j)_{j\in J}$ satisfasse les conditions suivantes :}
\begin{itemize}
\item[\textit{1°)}] \textit{L'application $j \longmapsto x_j$ de $\Ob J$ dans $\Ob\cC$ est constante, i.e. il existe un $x_0 \in \Ob\cC$ tel que $x_j = x_0\ \forall\, j \in \Ob J$.}
\item[\textit{2°)}] \textit{Il existe un projecteur $p$ dans $x_0$ tel que pour $j < j'$ dans $J$, on ait $\tau_{j',j} = p$.}
\item[\textit{3°)}] \textit{Le projecteur $p$ admet une image dans $\cC$.}
\end{itemize}
\textsc{Démonstration}. Ces conditions impliquent évidemment celles de 4.15.7.1, en prenant $p_j = p\ \forall\, j \in \Ob J$, du moins si $J$ n'admet de plus grand élément (ce qui fait une difficulté pour 2°). Mais si $J$ admet un objet final, il en est de même de $I$, et il est clair alors que tout système inductif indexé par $I$ est essentiellement constant. Donc, dans ces cas, les conditions 1°) 2°) 3°) ci-dessus impliquent que $\ux$ est essentiellement constant. Inverse\msp{page 444}ment, supposons que $\ux$ soit essentiellement constant, et prouvons que l'on peut réindexer comme indiqué. Le cas où $I$ admet un objet final est trivial, on prendra $J = e$, $p = \id$. Supposons donc que $I$ n'ait pas d'objet final. Comme $I$ grande devant $\pi_0$ et $\pi_0 > \card\Ob\cC$, on voit qu'il existe un $x_0 \in \Ob\cC$ tel que la sous-catégorie pleine $I' = I_{x_0}$ de $I$, formée des $i$ tels que $x_i = x_0$, soit cofinale dans $I$. Comme $(x_i)_{i\in I}$ est essentiellement constant, il en est de même de $(x_{i'})_{i'\in I'}$, et on peut lui appliquer le critère 4.15.7.1, avec \marginpar{$\Vert$}% double vertical bar printed in the left margin beside the two following lines
$J \longrightarrow I'$ cofinal. La démonstration montrera de plus que si $I'$ est grand devant $\pi_0$ (cardinal infini sans autres conditions), on peut supposer $J$ grande devant $\pi_0$ (\footnote{en tenir compte dans la démonstration de 4.15.7.1.}). De plus, il est clair qu'on peut supposer $J$ ordonnée. Comme $\pi_0 \geqslant \card\Hom(x_0,x_0)$, on voit encore qu'il existe un projecteur $p$ dans $x_0$ tel que la partie $J_1 = J'$ de $J$, formée des $j$ tels que $p_j = p$, soit cofinale. Étant cofinale, elle est également grande devant $\pi_0$. Donc on peut supposer que tous les $p_j$ sont égaux à un même $p$. Les conditions 1°-4° deviennent :
\begin{itemize}
\item[1°)] $p\tau_{j'j} = \tau_{j'j}p$ pour $j \leqslant j'$.
\item[2°)] Pour tout $j$, existe $j' \geqslant j$ avec
\[
p\tau_{j'j} \;=\; \tau_{j'j}p \;=\; \tau_{jj'}\ . % [sic] \tau_{jj'} -> \tau_{j'j}, index order used everywhere else
\]
\item[3°)] L'image $x_0^\natural$ du projecteur $p$ dans $x_0$ existe.
\item[4°)] L'endomorphisme $\tau^\natural_{j'j}$ de $x_0^\natural$ est un isomorphisme, pour $j \leqslant j'$.
\end{itemize}
\msp{page 445}
Or on a le

\dgrblock{dgrb1122}{Constancy in the ind-category and constancy of a diagram}
\begin{dcomment}{context}{p442}{Constancy in the ind-category and constancy of a diagram}
The distinction Grothendieck is using belongs to the original ind-object formalism. In SGA~4, Expos\'e~I, 8.5.3, the quotation marks around a formal inductive limit indicate its formation in the ind-category; they prevent its identification with a colimit calculated in the original category. The same expos\'e, 8.7.8, locates the Karoubi envelope among retracts of representables \citep{SGA4I1972}. These pages return to that distinction because a small base category is supposed to stop a sufficiently filtered diagram from wandering. A surviving projector is the precise residue of that wandering.

Corollary~4.15.7.2 has the stronger hypothesis that there are \emph{no nonidentity idempotents}. This is not another expression for idempotent completeness. A category of vector spaces is idempotent complete but has many projections. A groupoid has only identity idempotents and is therefore automatically idempotent complete. With the stronger hypothesis, every $p_j$ in the preceding lemma is $1_{x_j}$, and the image diagram is the original diagram.

There remains a small distinction even then. An isomorphism-valued diagram is naturally isomorphic to a constant diagram; its arrows need not literally be identities. To see the trivialization directly, fix $j_0$. For $j$, choose a common upper object $k$ and put
\[
\theta_j=x(j\to k)^{-1}x(j_0\to k):x_{j_0}\longrightarrow x_j.
\]
Filteredness makes this independent of the common upper object and of the chosen arrows, and gives $x(\lambda)\theta_j=\theta_{j'}$. The family $\theta$ is the required natural isomorphism. The later thinning argument obtains literal identity arrows by passing to a further subsystem. Keeping these two operations apart will matter in Theorem~4.15.9, where the printed condition~b) incorrectly demands identity arrows on an entire tail.
\end{dcomment}


\dgrblock{dgrb1123}{A two-element monoid which creates one missing retract}
\begin{dcomment}{examples}{p443}{A two-element monoid which creates one missing retract}
The projectors in the criterion are necessary even when the category is as small as possible. Let $\mathcal C$ have one object $c$ and endomorphism monoid
\[
M=\{1,p\},\qquad p^2=p,\qquad p\ne1.
\]
Choose any infinite regular ordinal $\kappa$, regard it as a directed poset, and form the diagram which has value $c$ at every stage and transition $p$ at every strict inequality. Composition holds because $p^2=p$; identity arrows still have value $1$.

The representable $h_c$ has value $M$ at the unique object. Every strict transition in its presheaf diagram is left multiplication by $p$, sending both elements to $p$. Consequently
\[
X(c)=\varinjlim_{\alpha<\kappa}M=\{*\}.
\]
The two elements of each stage become equal after one transition, and the resulting class is nonempty. The presheaf $X$ is therefore the one-element right $M$-set, with both elements of $M$ acting identically. It is a retract of $h_c$: the unique map $M\to\{*\}$ has the equivariant section $*\mapsto p$. Their composite on $M$ is precisely left multiplication by $p$.

Yet $X$ is not representable in $\mathcal C$, since its value has one element whereas the only representable has two. The ind-object is not essentially constant. The projectors commute, eventually absorb every transition, and induce identity maps on their images in the presheaf category; condition~3 of 4.15.7.1 fails exactly because that image does not belong to $\mathcal C$.

The Karoubi envelope repairs the defect by adjoining an object $d=(c,p)$. Explicitly,
\[
\operatorname{Hom}(c,d)=\{p\},\quad
\operatorname{Hom}(d,c)=\{p\},\quad
\operatorname{End}(d)=\{p\},
\]
where the last $p$ is the identity of $d$. The composite $d\to c\to d$ is $1_d$, and the reverse composite is the original endomorphism $p$ of $c$. In this enlarged category the same diagram represents $d$.

Taking $\kappa=\pi^+$ makes the indexing poset grand devant $\pi$ for any chosen infinite $\pi$. Thus increasing the cardinal cannot remove the need to split idempotents. It is the conclusion \emph{Karoubi envelope}, rather than the original category, which survives without a hypothesis on projectors.
\end{dcomment}


\textbf{Lemme 4.15.7.4.} \textit{Soient $\cC$ une catégorie, $x_0$ un objet de $\cC$, $p$ un projecteur dans $x_0$, $\tau$ un endomorphisme de $x_0$ satisfaisant les relations}
\[
\tau p \;=\; p\tau \;=\; \tau\ , \tag*{$(*)$}
\]
% foot: Version 17 décembre 2025
\printedpage{171}
% head: A. Grothendieck : Dérivateurs, Ch. XVIII, Édition des Universités de Montpellier II et Paris VII
\textit{et tel que l'automorphisme de $x_0^\natural \defeq$} $\Img p$ \textit{(défini p.ex. comme objet de $\cC^\wedge$, $x_0^\natural =$} $\Ker(\id_{x_0},p)$ % [sic] "(défini ... Ker(id_{x_0},p)" -> a closing parenthesis is missing, the parenthesis opened before "défini" is never closed
\textit{soit l'identité. Alors on a $\tau = p$.}

\textsc{Démonstration de} 4.15.7.4. On est ramené aussitôt, par le procédé standard, au cas où $\cC = \mathrm{Ens}$, où la vérification est immédiate : les automorphismes $\tau$ de l'ensemble $E$ $(= x)$ qui satisfait $(*)$ correspondent biunivoquement aux endomorphismes $u$ de $E^\natural = pE$, via
\[
\begin{cases}
u \;\longmapsto\; \sigma u \rho \\
\tau \;\longmapsto\; \tau|E^\natural
\end{cases}
\qquad
\text{(où }
\begin{tikzcd}[column sep=small]
E^\natural \arrow[r, hook, shift left, "\sigma"] & E \arrow[l, shift left, "\rho"]
\end{tikzcd}
\text{ sont les flèches canoniques)}\ ,
\]
donc $\tau|E^\natural = \id \Longrightarrow u = \sigma\rho = p$.

Choisissons un $j_0 \in J$, et considérons les homomorphismes
\[
\tau^\natural_{j,j_0} \;:\; x_0^\natural \longrightarrow x_0^\natural
\]
induits par les $\tau_{j,j_0}$ avec $j \geqslant j_0$. Ici encore, comme l'ensemble des isomorphismes possibles est de

\msp{page 446}
cardinal $\leqslant \pi_1$, et $J/j_0$ est grande devant $\pi_0$, on voit qu'il existe $J' \subset j_0\backslash J$, cofinale dans $j_0\backslash J$, donc dans $J$, telle que pour $j \in J'$, les $\tau^\natural_{j,j_0}$ soient égaux à un même $\tau : x_0^\natural \xrightarrow{\sim} x_0^\natural$. Donc si $j', j'' \in J'$, avec $j' \leqslant j''$, on a
\[
\underbrace{\tau_{j''j_0}}_{\tau} \;=\; \tau_{j''j'}\,\underbrace{\tau_{j'j_0}}_{\tau}\ ,
\]
d'où
\[
\tau_{j''j'} \;=\; \id_{x_0^\natural}\ ,
\]
puisque la flèche $\tau$ est un isomorphisme. Donc le système inductif réindexé par $J'$, la condition 4° est remplacée par la condition plus stricte \fbox{4° bis) $\tau^\natural_{j',j} = \id_{x_0^\natural}$.} Mais la condition 2° devient alors, grâce à 4.15.7.4,
\[
\boxed{\text{2° bis) Pour tout } j\text{, existe } j' \geqslant j \text{ avec } \tau_{j',j} = p}\ .
\]
Donc, on peut supposer 1°, 2° bis, 3°, 4° bis (on désigne à nouveau par $J$, au lieu de $J'$, l'ensemble ordonné grand devant $\pi_0$ indexant).

Soit $K$ l'ensemble ordonné défini par

\msp{page 447}
\[
\begin{cases}
\Ob K \;=\; \Ob J \\
\text{si } j,j' \in \Ob J\text{, on a } j \overset{K}{\leqslant} j' \text{ si et seulement si } j \overset{\mathrm{J}}{=} j'\text{, ou } j \overset{J}{<} j' \text{ et } \tau_{j',j} = p\,.
\end{cases}
\]
Il est immédiat que c'est une \textit{relation d'ordre} sur $J$ (car $p^2 = p$), plus fin que la relation d'ordre initiale. Le système inductif induit sur $K$ satisfait aux conditions 1°, 2°, 3° de 4.15.7.3. D'autre part, $K \longrightarrow J$ est cofinal, donc aussi le composé $K \longrightarrow J \longrightarrow I$. Il reste à prouver que $K$ est grande devant $\pi_0$, donc que toute partie $K' \subset K$ telle que $\card K' \leqslant \pi_0$ admet un majorant dans $K$. Elle admet un majorant \textit{dans} $J$, soit $j$. Par hypothèse 2 bis, il existe $j' \geqslant j$ tel que $\tau_{j'j} = p$, et on peut supposer $j' > j$. (En effet, si $\tau_{j'j} = p$, on voit que
% foot: Version 17 décembre 2025
\printedpage{172}
% head: A. Grothendieck : Dérivateurs, Ch. XVIII, Édition des Universités de Montpellier II et Paris VII
\msp{page 448}
la forme initiale 2°) de la condition 2° bis que l'on a $\tau_{j''j} = p$ pour tout $j'' \geqslant j'$, donc si $j$ n'est par objet final de $J$, % [sic] par -> pas, misprint
on peut supposer $j' > j$. Mais $J$ n'a pas d'objet final, car sinon $I$ en aurait un, cas que l'on avait éliminé. Donc il existe $j'$ tel que $j' \overset{K}{>} j$.) Or on a ceci (qu'on applique au cas $i \in K'$) :
\[
i \leqslant j \ \text{ et } \ j \overset{K}{<} j' \qquad \Longrightarrow \qquad i \overset{K}{<} j'\ .
\]
\[
\underbrace{x_i}_{= x_0} \xrightarrow{\tau_{j,i}} x_j \xrightarrow{\tau_{j',j}=p} x_{j'}\ ,
\]
\[
\tau_{j'i} \;=\; \tau_{j'j}\tau_{ji} = \underbrace{p\tau_{ji}}_{\overset{\text{1°}}{=}\,\tau_{ji}p} \;\overset{?}{=}\; p\ ,
\]
on a en effet
\[
\left.
\begin{aligned}
\tau_{j',i} &\;=\; p\tau_{j',i} \\
\tau_{j',i} &\;-\; \tau_{j',i}p % [sic] - -> =, the second relation of (*) is meant
\end{aligned}
\right\}
\quad \overset{\text{par 4.15.7.4}}{\Longrightarrow} \quad % Doubt resolved against PDF leaf 14: no underline; the two strokes belong to the implication arrow.
\tau_{j',i} \;=\; p\ .
\]
Cela achève la démonstration du corollaire 4.15.7.3.

\textbf{Corollaire 4.15.7.5} (de 4.15.7.3). \textit{Soit $(x_i)_{i\in I}$ un système inductif essentiellement constant grand devant $\pi_0 \geqslant$} $\card\cC$\textit{, dans $\cC$} sans projecteurs. \textit{Alors il existe $J \longrightarrow I$ cofinal, avec $J$ ordonné grand devant $\pi_0$, et tel que le système inductif induit sur $J$ soit} constant.

\textsc{Démontrons le lemme} 4.15.7.1. Il faut expliciter l'existence de deux flèches dans $\Ind(\cC)$,
\[
z \underset{\beta}{\overset{\alpha}{\rightleftarrows}} \text{"}\varinjlim_{i\in I}\text{"} x_i \;=\; \ux\ ,
\]
avec $\beta\alpha = \id_z$, $\alpha\beta = \id_{\ux}$. La donnée

\msp{page 449}
de $x$ équivaut à celle d'un élément de $\varinjlim\limits_{i\in I} \Hom(z,x_i)$, et est donc donné par un % [sic] x -> \alpha, parallel with "la donnée de \beta" below
\[
\alpha_{i_0} \;:\; z \longrightarrow x_{i_0}\ ,
\]
la donnée de $\beta$ équivaut à celle d'un élément de $\varprojlim\limits_{i\in I} \Hom(x_i,z)$, i.e. de flèches
\[
\beta_i \;:\; x_i \longrightarrow z\ ,
\]
avec la condition de compatibilité habituelle. Quitte à remplacer $I$ par $I'$, avec $I' \longrightarrow I$ cofinal, on peut supposer $I$ ordonnée (et grande devant $\pi_0$, si $I$ initiale l'était), et la condition de compatibilité sur les $\beta_i$ s'écrit
\[
\beta_i\tau_{ji} \;=\; \beta_j \qquad \text{si } i \leqslant j \text{ dans } I,
\]
où $\tau_{ji} : x_i \longrightarrow x_j$ désigne la flèche de transition. Il faut expliciter, en termes de ces données, les conditions
\[
\alpha\beta \;=\; \id_z\ , \qquad \beta\alpha \;=\; \id_{\ux}\ .
\]
% foot: Version 17 décembre 2025

\dgrblock{dgrb1124}{A cofinal subsystem with one projector}
\begin{dcomment}{exposition}{p448}{A cofinal subsystem with one projector}
The purpose of 4.15.7.3 is stronger than representability. Grothendieck needs a presentation whose arrows in the small base are literally controlled, so that a diagram upstairs can be studied in one fibre. The argument combines a cardinal pigeonhole principle with a careful deletion of arrows. Both operations preserve the filteredness needed later.

Throughout this argument, grand devant $\pi_0$ means that diagrams with at most $\pi_0$ morphisms admit a cocone. For a poset this says that every subset of cardinal at most $\pi_0$ has an upper bound; in modern terminology it is $\pi_0^+$-directed. Suppose such a poset $I$ is coloured by a set $A$ of at most $\pi_0$ colours. If no colour fibre were cofinal, we could choose, for every $a\in A$, an element $i_a$ above which colour $a$ never occurs. A common upper bound of all $i_a$ would have no colour. Hence one fibre is cofinal. A cofinal full subposet remains $\pi_0^+$-directed, since an upper bound in $I$ can be moved into that subposet.

Replace the original index by a final $\pi_0^+$-directed poset, using the usual cofinal-poset theorem (also \citealp[Proposition~5.3.1.18]{Lurie2009HTT} in its higher form). First apply this principle to the objects $x_i$, and then to the projectors supplied by 4.15.7.1. We obtain a fixed object $x_0$ and a fixed split idempotent $p$. Fixing a base index and applying the principle once more makes the isomorphisms from its image to the later images equal. Cancellation of these isomorphisms makes all induced transitions on $\operatorname{Im}p$ identities. There are at most $\pi_0$ possible image automorphisms: if $p=sr$ with $rs=1$, the map $u\mapsto sur$ embeds their set into $\operatorname{End}_{\mathcal C}(x_0)$.

The elementary identity behind the remaining step is
\[
\tau p=p\tau=\tau,\quad r\tau s=1
\quad\Longrightarrow\quad
\tau=sr\tau sr=s(r\tau s)r=p.
\]
Thus the eventual absorption condition says that every index eventually maps by $p$.
Define $K$ on the same objects by retaining identities and only those strict arrows whose value is $p$. Transitivity follows from $p^2=p$. To see that $K$ remains $\pi_0^+$-directed, take at most $\pi_0$ objects and choose a common upper bound $j$ in the old order. Choose $j'>j$ with $\tau_{j'j}=p$; the terminal-object case is treated separately by a singleton index. For each chosen $i$,
\[
\tau_{j'i}=p\tau_{ji}.
\]
This arrow is absorbed by $p$ on both sides, and its action on the image is the identity. The displayed identity therefore gives $\tau_{j'i}=p$. Hence $j'$ is an upper bound in $K$.

The inclusion $K\to J$ is final as a functor, not merely surjective on objects. For a fixed $j\in J$, the comma category $j/K$ is nonempty; any two of its objects acquire a common upper object in $K$ by the same construction. It is a directed poset and therefore connected. This is the missing explicit finality verification. The conclusion has exactly the form required: one object, one idempotent on every strict transition, and a sufficiently directed indexing category. The printed $\pi_1$ and lower slice $J/j_0$ on manuscript p.~446 should be read as $\pi_0$ and $j_0\backslash J$; they are retained in the source text and registered separately.
\end{dcomment}


\printedpage{173}
% head: A. Grothendieck : Dérivateurs, Ch. XVIII, Édition des Universités de Montpellier II et Paris VII
\[
\begin{tikzcd}
z \arrow[r, "\alpha_{i_0}"] \arrow[dr] & x_{i_0} \arrow[r, "\tau_{i,i_0}"] \arrow[d, "\beta_{i_0}"] & x_i \arrow[dl, "\beta_i"] \\
 & z &
\end{tikzcd}
\]
La relation $\alpha\beta = \id_z$ équivaut à
\[
\beta_{i_0}\alpha_{i_0} \;=\; \id_z\ , \tag*{$(*_{i_0})$}
\]
elle implique d'ailleurs, en posant pour $i \geqslant i_0$,
\[
\alpha_i \;\defeq\; \tau_{i,i_0}\alpha_{i_0} \;:\; z \longrightarrow x_i\ ,
\]
que l'on a pour $i \geqslant i_0$
\[
\beta_i\alpha_i \;=\; \id_z \qquad \forall\, i \geqslant i_0\ . \tag*{$(*_i)$}
\]

\msp{page 450}
Posons alors
\[
p_i \;=\; \alpha_i\beta_i \;:\; x_i \longrightarrow x_i\ ,
\]
il est clair qu'on a
\[
\begin{cases}
p_i^2 \;=\; p_i \qquad \text{($p_i$ un projecteur)} \\
\alpha_i \;:\; z \longrightarrow x_i \qquad \text{induit un isomorphisme} \\
\qquad z \xrightarrow{\sim} x_i^\natural \;\defeq\; \Img p_i \;=\; \Ker(\id_{x_i},p_i)\ .
\end{cases}
\]
L'ensemble ordonné cofinal cherché $J$ est $i_0\backslash I$, il est filtrant, et est grand devant $\pi_0$ si $I$ l'est. À vrai dire, la réduction au cas $I$ ordonné est inutile, on peut prendre directement $J = i_0\backslash I$, il est immédiat que $J$ est filtrante ($I$ l'étant), et $i_0\backslash I \longrightarrow I$ cofinal, enfin que si $I$ est grande devant un cardinal infini $\pi_0$, il en est de même de $i_0\backslash I$.

Exprimons que $\beta\alpha = \id_{\ux}$ ; or on voit que $\beta\alpha$ est défini par l'endomorphisme $(p_j)_{j\in J}$ du système inductif $(x_j)_{j\in J}$ -- le fait que ce soit bien un endomorphisme de ce système, i.e. qu'on ait la relation 1°) de 4.15.7.1,
\[
p_{j'}x_\lambda \;=\; x_\lambda p_\lambda \qquad \text{si } \lambda : j \longrightarrow j' \text{ dans } J, % [sic] p_\lambda -> p_j, cf. relation 1°) of 4.15.7.1
\]
est immédiat, ça s'écrit
\[
\alpha_{j'}\underbrace{\beta_{j'}x_\lambda}_{\beta_j} \;=\; \underbrace{x_\lambda\alpha_j}_{\alpha_{j'}}\beta_j\ , \qquad \text{or les deux membres sont égaux à } \alpha_{j'}\beta_j.
\]
% foot: Version 17 décembre 2025

% ---- from der-ch18_p176-180.tex
\printedpage{174}
% head: A. Grothendieck : Dérivateurs, Ch. XVIII, Édition des Universités de Montpellier II et Paris VII

\msp{page 451}
Dire que c'est l'identité, est exprimé par la condition 2°) sur $(x_j,p_j)_{j\in J}$. La condition 3° est clair par construction, $p_j$ admettant une image isomorphe à $z$ par $\alpha_j$. D'autre part, on voit \texttt{[tout]} de suite que pour une flèche $\lambda : j \longrightarrow j'$ de $J$, on a commutativité dans % [sic] est clair -> est claire, agreement with "condition"
\[
\begin{tikzcd}
x_j \arrow[r, "x_\lambda"] & x_{j'} \\
x_j^\natural \arrow[u, hook] \arrow[r, "x_\lambda"] & x_{j'}^\natural \arrow[u, hook] \\
z \arrow[u, "\wr", "\overline{\alpha_j}"'] \arrow[r, "\mathrm{id}_z", "\sim"'] & z \ , \arrow[u, "\wr", "\overline{\alpha_{j'}}"']
\end{tikzcd}
\]
où $\overline{\alpha_j}$, $\overline{\alpha_{j'}}$ sont les isomorphismes définis par $\alpha_j$, $\alpha_{j'}$ respectivement. Cela prouve que $x_\lambda^\natural$ est un isomorphisme.

Cela prouve que si $(x_i)_{i\in I}$ est essentiellement constant, il existe $J\longrightarrow I$ comme indiqué. L'inverse est également vrai, car on peut supposer $J = I$, puisque $(x_i)_{i\in I}$ et $(x_j)_{j\in J}$ définissent des ind-objets isomorphes. Alors le système inductif des $x_j^\natural$ a des morphismes de transition des isomorphismes, et comme $J$ est filtrante donc 1-connexe, on en déduit facilement qu'il est isomorphe
\msp{page 452}
à un système inductif constant, de valeur $z$. En choisissant un tel isomorphisme
\[
z \xrightarrow[\sim]{\overline{\alpha_i}} x_i^\natural \ ,
\]
on trouve $z \xrightarrow{\alpha} \underline{x}$, et le système des $p_i$ définit des
\[
\begin{tikzcd}
x_i \arrow[r, "\overline{p_i}"] \arrow[rr, bend right=12, "\beta_i \overset{\text{déf}}{=} \overline{\alpha_j}^{-1}\overline{p_i}"'] & x_i^\natural \arrow[r, "\overline{\alpha_j}^{-1}"] & z \ ,
\end{tikzcd}
\] % [sic] \overline{\alpha_j} (index j) -> \overline{\alpha_i}, the arrow goes from x_i^\natural, cf. z -> x_i^\natural above
et on voit facilement que l'on aura
\[
\beta\alpha \ = \ id_z \ , \qquad \alpha\beta \ = \ \mathrm{id}_{\underline{x}} \ ,
\]

\hfill q.e.d.

\textsc{Je peux maintenant prouver} 4.15.3. C'est en effet essentiellement la conjonction des 4.15.7.5 (p. 448), et de 4.8.13, qui nous dit essentiellement que si $\mathcal{C}$ est \textit{karoubienne} (condition vérifiée si $\mathcal{C}$ est sans projecteurs), alors tout système inductif dans $\mathcal{C}$, grand devant $\pi_0$, où $\pi_0 \geqslant \card\,\text{réd}\,\mathcal{C}$, est essentiellement constant,

\hfill q.e.d.

\printedpage{175}
% head: A. Grothendieck : Dérivateurs, Ch. XVIII, Édition des Universités de Montpellier II et Paris VII

\msp{page 453}
\textbf{4.15.8.} Je veux revenir ici sur la démonstration inachevée de 4.8.13 (p. 212-216). Je reprends la démonstration où je l'avais laissée (p. 215). Il faut prouver que
\[
F(\xi) \ = \ \varinjlim_{i\in I} M_i
\]
est de cardinal $\leqslant \pi_0$. J'ai annoncé qu'on peut se ramener au cas où $M$ est un monoïde de monomorphismes, i.e. que $uv = uv' \Longrightarrow v = v'$ (pour $u,v,v' \in M$). Je vais expliciter cet argument.

Pour tout $i\in I$, soit
\[
N_i \ = \ \left\{ u \in M \ \middle| \ \begin{array}{l}
\text{l'ensemble } I_{i,u} \text{ des } j\in I \text{ tels}\\
j \geqslant i,\ \alpha_{ji} = u \text{ est cofinal}\\
\text{dans } I\text{, i.e. pour tout } j_0\in I,\\
\text{existe } j \in I_{i,u}\text{, avec } j \geqslant j_0
\end{array} \right\} .
\tag*{(4.15.8.1)}
\]
Comme $I$, donc $i\backslash I$ est grande devant $\pi_0$, et $\card M \leqslant \pi_0$, on voit aussitôt que
\[
N_i \ \neq \ \emptyset \ .
\tag*{(4.15.8.2)}
\]
$\Big($ Si d'autre part $i' \geqslant i$, on trouve
\[
N_i \ = \ N_{i'}\alpha_{i'i} \ .
\tag*{(4.15.8.3)}
\]
L'inclusion $N_{i'}\alpha_{i'i} \subset N_i$ est $\pm$ tautologique, prouvons $N_i \subset N_{i'}\alpha_{i'i}$. Soit $u \in N_i$, l'ensemble $J = I_{i,u} \cap (i'\backslash I)$ des $j \geqslant i'$ tels que $\alpha_{ji} = u$ est cofinal dans $I$, donc grand devant $\pi_0$.
\[
\begin{tikzcd}
i \arrow[r] \arrow[rr, bend right=12, "u = \alpha_{ji}"'] & i' \arrow[r, "\alpha_{ji'}"] & j
\end{tikzcd}
\]
Donc il existe un $v \in M$ tel que l'ensemble $J'$ de $j\in J$ tels que $\alpha_{ji'} = v$
\msp{page 454}
soit cofinal dans $J$, donc dans $I$. Comme cet ensemble $J'$ est contenu dans $I_{i',v}$, cela implique que $I_{i',v}$ est cofinal dans $I$, i.e. $v \in N_{i'}$. On a d'autre part, pour $j \in J'$,
\[
\underbrace{\alpha_{ji'}}_{v}\alpha_{i'i} \ = \ \underbrace{\alpha_{ji}}_{u} \ ,
\]
d'où $u = v\alpha_{i'i}$ avec $v\in N_{i'}$, d'où $u \in N_i\alpha_{i'i}$, q.e.d. $\Big)$ % [sic] N_i\alpha_{i'i} -> N_{i'}\alpha_{i'i}, the inclusion to be proved is N_i \subset N_{i'}\alpha_{i'i}

Pour $i \in I$, soit
\[
\alpha_i \ : \ \underbrace{\mathrm{Hom}(\xi, \underbrace{x_i}_{=\,\xi})}_{=\,M} \ \longrightarrow \ \underbrace{\mathrm{Hom}(\xi,F) \ = \ F(\xi)}_{\simeq\ \varinjlim\limits_{i} \mathrm{Hom}(\xi,x_i)}
\]
\printedpage{176}
% head: A. Grothendieck : Dérivateurs, Ch. XVIII, Édition des Universités de Montpellier II et Paris VII
le morphisme canonique. On trouve ceci :
\begin{center}
\makebox[0.20\linewidth][l]{(4.15.8.4)}%
\begin{minipage}[c]{0.78\linewidth}
Si $u \in N_i$, alors la relation d'équivalence $R_i$ dans $M$ définie par $\alpha_i$ ($v \overset{R_i}{\sim} v' \Longleftrightarrow \alpha_i(v) = \alpha_i(v')$) est égale à celle définie par l'application $v \longmapsto uv$ de $M$ dans $M$, i.e.
\[
\alpha_i(v) \ = \ \alpha_i(v') \qquad \Longleftrightarrow \qquad uv \ = \ uv' \ .
\]
\end{minipage}
\end{center}

C'est immédiate, puisque $I_{i,u}$ étant cofinal dans $I$, on a $\alpha_i(v) = \alpha_i(v')$ si et seulement si $\exists\ j \in I_{i,u}$ avec $\alpha_{ji}v = \alpha_{ji}v'$, or $\alpha_{ji} = u$. % [sic] C'est immédiate -> C'est immédiat, agreement

Pour $u \in M$, notons $R_u$ la relation d'équivalence dans $M$ définie par $v \longmapsto uv$. L'ensemble des relations d'équivalence dans \texttt{[}$M$\texttt{]}
\msp{page 455}
obtenu ainsi est de cardinal $\leqslant \card M \leqslant \pi_0$, donc l'ensemble des relations d'équivalences $R_i$ ($i \in I$) est lui aussi de cardinal $\leqslant \pi_0$. Cela montre, $\pi_0$ étant grand devant $\pi_0$, qu'il existe une relation d'équivalence de la forme $R_u$ ($u \in M$), et un ensemble cofinal $I'$ dans $I$, tel que $\forall\ i \in I'$ on ait $R_i = R_u$. Cela montre que, quitte à remplacer $I$ par l'ensemble ordonné cofinal $I'$, qu'on peut supposer que l'on a
\[
R_i \ = \ R_u \qquad \forall\, i \in I \ .
\tag*{(4.15.8.5)}
\]
Soit alors
\[
E \ = \ M/R_u \ = \ M/R_i \qquad (\forall\ i \in I) \ ,
\]
on voit donc que les $\alpha_{ji} \in M$ (pour $i \leqslant j$) sont tels que l'application
\[
\tau_{\alpha_{ji}} \ : \ v \ \longmapsto \ \alpha_{ji}v \ : \ M \ \longrightarrow \ M
\]
est compatible avec la relation d'équivalence $R_u$, donc passe au quotient en une application
\[
\overline{\alpha_{ji}} \ : \ E \ \longmapsto \ E \ ,
\]
laquelle est d'ailleurs \textit{injective.}

Soit $M'$ le sous-monoïde de $M$ formé des $u \in M$ qui ont la propriété de stabilité précédente, ainsi que la propriété d'injectivité qu'on vient de dire, et $\overline{M'}$ le monoïde image de $M'$ dans le monoïde
\[
\mathrm{Mon}(E) \ \subset \ \mathrm{End}_{\mathrm{Ens}}(E)
\]
\msp{page 456}
des applications injectives de $E$ dans lui-même. On voit donc que le système inductif donné, défini par les $\alpha_{ji}$, prend ses valeurs dans le sous-monoïde $M'$ de $M$ (i.e. de la sous-catégorie correspondante de la catégorie $B_M$), et que la limite inductive $F(\xi)$ qui nous intéresse (pour prouver $\card F(\xi) \leqslant \pi_0$) est isomorphe à la limite inductive du système inductif d'un système d'ensembles tous égaux à $E$, les morphismes de transition étant les $\overline{\alpha_{ji}} \in \overline{M'} \subset \mathrm{Mon}(E)$. De plus, on aura $\card \overline{M'} \leqslant \card M \leqslant \pi_0$. Le résultat $\card F(\xi) \leqslant \pi_0$ voulu résultera alors du

\dgrblock{dgrb1125}{Removing identifications before counting new elements}
\begin{dcomment}{exposition}{p455}{Removing identifications before counting new elements}
The obstruction to a cardinal estimate is that a transition may both identify old elements and introduce new ones. Grothendieck separates these effects. The sets $N_i$ record the transition maps which recur cofinally; each such map already detects every identification made after stage $i$.

For $u\in N_i$ and $v,v'\in M$, equality of the two classes in $\varinjlim M_j$ is witnessed at some later stage. Moving further into the cofinal fibre where the transition from $i$ equals $u$ gives
\[
[v]_i=[v']_i\quad\Longleftrightarrow\quad uv=uv'.
\]
The converse follows at any stage in that fibre. Thus the kernel relation $R_i$ is one of the at most $|M|$ relations $R_u$. The cofinal-colour argument gives a subsystem with $R_i=R$ independent of $i$.

Passing from $M$ to $E=M/R$ now removes exactly the identifications which the final colimit would make. The induced transition maps of $E$ are injective: equality of the images of $[v]$ and $[v']$ implies equality of their final classes, hence $vRv'$. The total colimit is unchanged, because no additional equivalence has been imposed. What remains to be counted is actual growth by inclusions.

There is a useful bookkeeping point in this construction. The elements of $M$ which preserve $R$ and act injectively on $E$ form a \emph{submonoid} $M_1$. Its image
\[
\overline{M_1}\subseteq\operatorname{Inj}(E,E)
\]
is a quotient of $M_1$, has cardinal at most $|M|$, and is left cancellative. It need not be a quotient of the whole monoid $M$. Consequently the reduction on these pages proves a sufficient criterion involving such subquotients. The restricted quotient wording in 4.15.9~b) is not justified by this reduction alone.

The distinction does not affect the strict-cardinal conclusion. Its counting argument applies to \emph{every} left-cancellative monoid below the chosen bound, including these subquotients. It does matter if one tries to read 4.15.9 as a sharp equivalence for one fixed category and one fixed cardinal. We retain the useful reduction while recording exactly what its proof supplies.
\end{dcomment}


\textbf{Lemme-clef 4.15.8.1.} \textit{Soit $M$ un monoïde de monomorphismes (i.e. tel que $uv = uv'$ $\Longrightarrow v = v'$, pour $u,v,v' \in M$), et soit $(\alpha_{ji})_{i\leqslant j}$ un système inductif dans $B_M$, indexé par un ensemble ordonné $I$ grand devant $\pi_0$, où $\pi_0$ est un cardinal}
\printedpage{177}
% head: A. Grothendieck : Dérivateurs, Ch. XVIII, Édition des Universités de Montpellier II et Paris VII
\msp{page 457}
\textit{infini tel que} $\card M \leqslant \pi_0$. \textit{Alors il existe un $i_0 \in I$ tel que pour tout $i \geqslant i_0$, $\alpha_{i_0,i}$ soit un isomorphisme.}

\dgrblock{dgrb1126}{An ordinary filtered sequence which really grows}
\begin{dcomment}{examples}{p456}{An ordinary filtered sequence which really grows}
The increased filteredness assumption is doing real work. To measure it against a familiar sequence, take the one-object category with endomorphism monoid $(\mathbf N,+)$, whose identity is $0$. At every integer stage put the unique object, and give the transition $i\to j$ the value $j-i$. Every strict transition is a nonunit.

Evaluating its ind-object at the unique object gives the sequence of regular right $\mathbf N$-sets
\[
\mathbf N\xrightarrow{\,n\mapsto n+1\,}\mathbf N
\xrightarrow{\,n\mapsto n+1\,}\mathbf N\longrightarrow\cdots .
\]
Its colimit is $\mathbf Z$, with the usual action by addition of nonnegative integers. An explicit isomorphism sends the class of $a$ at stage $i$ to $a-i$. The transition relation preserves this difference. Every integer occurs by choosing $i$ large enough, and two representatives with the same difference become equal at any stage above their two indices. This verifies both surjectivity and injectivity of the description.

The resulting right $\mathbf N$-set is not representable. The sole representable is $\mathbf N$, and an equivariant map $\mathbf Z\to\mathbf N$ cannot exist: if its value at $0$ were $b$, equivariance would force the value at $-n$ to be $b-n$, which is negative for $n>b$. In particular $\mathbf Z$ is not a retract of the representable. Equivalently, the only idempotent of the additive monoid $\mathbf N$ is $0$, so its Karoubi envelope adds nothing.

This sequence is filtered in the ordinary sense, but its index $\mathbf N$ is not grand devant $\aleph_0$: the whole countable set of indices has no upper bound. The next proposition shows that a diagram in the same monoid indexed by an $\aleph_1$-directed poset must have a constant tail. There is no contradiction with the explicit colimit above. The regular cardinal attached to the word filtered, and the source's successor convention, determine which of the two phenomena occurs.
\end{dcomment}


Si on applique ce lemme au cas de $\overline{M'}$ et du $\overline{\alpha_{ji}}$, on trouve donc qu'il existe un $i_0 \in I$ tel que les $\overline{\alpha_{ji_0}}$ soient des isomorphismes. Cela implique que $F(\xi) \simeq \varinjlim\limits_{i} \underbrace{E_i}_{E}$ est isomorphe à $E_{i_0}$, donc est de cardinal égal à $\card E_{i_0} = \card E \leqslant \pi_0$, ce qui achèvera la démonstration.

D'ailleurs, l'énoncé 4.8.13 implique 4.15.8.1, car dans le système inductif des $(M_i = M, \alpha_{ji})$, de $\varinjlim$ égale à $F(\xi)$, s'il n'existait \textit{pas} de $i_0$ comme stipulé dans 4.15.8.1, alors on en conclurait aisément que l'on aurait $\card F(\xi) > \pi_0$, contrairement à 4.8.13. Nous allons dégager ceci dans un lemme :

\msp{page 458}
\textbf{Lemme 4.15.8.2.} \textit{Soit $(E_i)_{i\in I}$ un système inductif d'ensembles, indexé par un ensemble ordonné $I$ grand devant $\pi_0$. On suppose qu'il existe $i_0 \in I$ tel que}
\begin{enumerate}[nosep]
\item[\textit{a)}] \textit{$i \geqslant i_0$ implique} $\card E_i = \pi_0$, \textit{et}
\item[\textit{b)}] \textit{$j \geqslant i \geqslant i_0$ implique que $u_{ji} : E_i \longrightarrow E_j$ est un monomorphisme. Supposons de plus}
\item[\textit{c)}] \textit{$\forall\ i \geqslant i_0$ $\exists\ j > i$ tel que $u_{ji}$ ne soit} pas \textit{un isomorphisme (i.e. ne soit pas surjectif). Soit $E_\infty = \varinjlim\limits_{i\in I} E_i$. Alors} $\card E_\infty > \pi_0$.
\end{enumerate}
Ce lemme implique donc l'\textit{équivalence} de 4.8.13 et du “lemme-clef” 4.15.8.1. Pour prouver 4.15.8.2, nous allons prouver quelque chose d'un peu plus précis :

\dgrblock{dgrb1127}{A boundary case controlled by right cancellation}
\begin{dcomment}{development}{p457}{A boundary case controlled by right cancellation}
Grothendieck has reduced his remaining difficulty to a monoid of injections. The source obtains a strict cardinal bound below. At the boundary, a recurring transition supplies an additional way to conclude if it can also be cancelled on the right. We record that limited extension, since it identifies the algebraic obstacle more precisely than the cardinal count alone.

\textbf{Proposition: a cancellable recurring transition.}\ours\
Let $\pi$ be infinite, let $I$ be a $\pi^+$-directed poset, and let $(a_{ji})$ be a diagram in a left-cancellative monoid $M$ of cardinal at most $\pi$. Suppose that, for some $i_*$, a cofinally recurring value $u$ of $a_{ji_*}$ is right cancellable: $vu=wu$ implies $v=w$. Then there is $j_0$ such that every $a_{ij_0}$, $i\ge j_0$, is a unit. On a cofinal full subposet all transitions are identities.

\emph{Proof.} Let $J=\{j\ge i_*:a_{ji_*}=u\}$, which is cofinal by assumption. If $j\le k$ belong to $J$, the cocycle identity gives
\[
u=a_{ki_*}=a_{kj}a_{ji_*}=a_{kj}u.
\]
Right cancellation gives $a_{kj}=1$. Choose $j_0\in J$. Given $i\ge j_0$, choose $k\in J$ with $k\ge i$. Then
$a_{ki}a_{ij_0}=a_{kj_0}=1$.
In a left-cancellative monoid, $ab=1$ also implies $ba=1$, since
$a(ba)=a=a1$. Thus $a_{ij_0}$ is a unit. Moreover, for $j\ge i\ge j_0$,
\[
a_{ji}=a_{jj_0}a_{ij_0}^{-1}
\]
is a unit. The full subsystem on $J$ already has identity transitions. This proves all assertions.

\textbf{Corollary: two-sided cancellation.}\ours\
The key lemma 4.15.8.1 holds with $|M|\le\pi$, including equality, if $M$ is cancellative on both sides. More generally, if $|\operatorname{Mor}\mathcal C|\le\pi$ and every endomorphism monoid of $\mathcal C$ is cancellative on both sides, every $\pi^+$-filtered diagram in $\mathcal C$ is essentially constant.

\emph{Proof.} In the monoid case the cofinal-colour principle supplies a recurring $u$, and every $u$ is right cancellable. For a general category, first replace the index by a final $\pi^+$-directed poset. Apply that principle to its at most $\pi$ object values, obtaining a cofinal full subposet with one value $x$. The diagram now takes values in $\operatorname{End}(x)$, and the proposition supplies a further cofinal subsystem with identity transitions. Its presheaf colimit is $h_x$.

We checked the small-category and Cauchy-completion treatments cited below, together with searches for the left-cancellative stabilization formulation; we did not find this exact recurring-value statement in the literature consulted. The proof is an elementary continuation of Grothendieck's cofinal-fibre argument, with no claim of priority. It does not settle the general left-cancellative case: from $vu=u$ one cannot cancel $u$ on the right without the extra hypothesis.
\end{dcomment}


\textbf{Corollaire 4.15.8.3.} \textit{Sous les conditions de 4.15.8.2, soit $\alpha_0$ le premier ordinal tel que} $\card \alpha_0 > \pi_0$, \textit{et soit $I_{\alpha_0}$ l'ensemble des ordinaux $\beta < \alpha$ (ensemble ordonné qui est grand devant $\pi_0$). Il existe alors une application strictement croissante}
\[
I_{\alpha_0} \ \longrightarrow \ I \ , \qquad \alpha \ \longmapsto \ i_\alpha \ , \qquad \textit{appliquant } 0 \textit{ dans } i_0 \ ,
\]
\msp{page 459}
\textit{tel que pour tout $\alpha \in I_{\alpha_0}$, on ait}
\[
\underbrace{u_{i_{\alpha+1},i_\alpha}}_{\substack{\textit{c'est un } \text{monomorphisme,}\\ \textit{d'après b, vu que } i_{\alpha+1} > i_\alpha \geqslant i_0}} \qquad\qquad \text{\textit{n'est pas un isomorphisme} (\footnotemark[187]).}
\]
\footnotetext[187]{\textit{mieux} : $\forall\ \alpha \in I_{\alpha_0}$ $\varinjlim\limits_{\beta\in I_\alpha} E_{i_\beta} \longrightarrow E_{i_\alpha}$ n'est pas surjectif.}
\textit{On a alors}
\[
E'_\infty \ \overset{\text{déf}}{=} \ \underbrace{\varinjlim_{\alpha\in I_{\alpha_0}} E_{i_\alpha}}_{} \ \hookrightarrow \ \varinjlim_{i\in I} E_i \ \overset{\text{déf}}{=} E_\infty
\]
\textit{et}
\[
\card E'_\alpha \ \geqslant \ \underbrace{\card I_{\alpha_0}}_{=\ \textit{successeur de } \pi_0} \ > \ \pi_0 \ ,
\] % [sic] E'_\alpha -> E'_\infty, the following line has \card E'_\infty
\textit{d'où} $\card E_\infty \geqslant \card E'_\infty > \pi_0$.

\textsc{Démonstration}. La construction des $i_\alpha$ par récurrence transfinie est immédiate, en utilisant c) et le fait que $I$ est grande devant $\pi_0$, et que pour tout $\alpha \in I_{\alpha_0}$, l'ensemble $I_\alpha$
\printedpage{178}
% head: A. Grothendieck : Dérivateurs, Ch. XVIII, Édition des Universités de Montpellier II et Paris VII
des $\beta < \alpha$ est de cardinal $\leqslant \pi_0$. Le fait que $E'_\infty \longrightarrow E_\infty$ soit injectif résulte de b). Le fait que
\[
\card E'_\infty \ \geqslant \ \card I_{\alpha_0}
\]
résulte du fait que $E'_\infty$ est réunion croissante de sous-ensembles identifiés à $E_{i_\alpha}$, et pour tout
\msp{page 460}
$\alpha$ on peut trouver un
\[
\xi_\alpha \ \in \ E_{i_\alpha}\backslash \bigcup_{\beta<\alpha} E_{i_\beta} \ .
\]
L'application
\[
\begin{aligned}
\alpha \ &\longmapsto \ \xi_\alpha\\
I_{\alpha_0} \ &\longrightarrow \ E'_\infty
\end{aligned}
\]
est injective, donc $\card E'_\infty \geqslant \card I_{\alpha_0}$.

% vertical bar printed in the left margin beside the whole of Corollaire 4.15.8.4
\textbf{Corollaire 4.15.8.4.} \textit{Pour un $M$ et un $\pi_0$ donnés, pour que le lemme-clef 4.15.8.1 soit valable, il faut et il suffit qu'il le soit dans le cas où $I = I_{\alpha_0}$, $\alpha_0$ étant le premier ordinal tel que} $\card \alpha_0 > \alpha_0$. \textit{Il revient au même qu'on ne puisse trouver de système inductif $(\alpha_{ji})_{j\geqslant i}$ indexé par $I_{\alpha_0}$, à valeurs dans $M$, tel que pour $i < j$, on ait $\alpha_{ji}$ non inversible dans $M$.} % [sic] \card \alpha_0 > \alpha_0 -> \card \alpha_0 > \pi_0, cf. Corollaire 4.15.8.3

\msp{page 461}
Je vais résumer ce qui a été obtenu dans un énoncé récapitulatif :

\dgrblock{dgrb1128}{The transfinite growth argument, including limit stages}
\begin{dcomment}{exposition}{p460}{The transfinite growth argument, including limit stages}
The point of 4.15.8.2 is to turn the failure of stabilization into too many distinct elements. Once transitions are injective, an element genuinely added at one stage can never disappear. The only care needed is at limit ordinals, where successive strict growth by itself does not guarantee that the next stage adds something beyond the whole preceding union.

Put $\kappa=\pi_0^+$. Starting at the prescribed $i_0$, choose a chain $i_\alpha$ for $\alpha<\kappa$. At a successor stage use hypothesis~c) to choose a nonsurjective transition. At a nonzero limit $\alpha$, the earlier indices form a set of cardinal at most $\pi_0$, so they have an upper bound $b$. Pass from $b$ to a stage $i_\alpha$ by a nonsurjective transition. Its image contains the image of every earlier stage, and is a proper subset of $E_{i_\alpha}$. This produces precisely the strengthened assertion in footnote~187:
\[
\varinjlim_{\beta<\alpha}E_{i_\beta}\longrightarrow E_{i_\alpha}
\quad\hbox{is not surjective}.
\]
At $\alpha=0$ the domain is empty and the source hypothesis $|E_{i_0}|=\pi_0$ makes the assertion valid.

Choose $\xi_\alpha$ outside the preceding union. Injectivity ensures that all these elements remain distinct in the chain colimit. That colimit injects into the original colimit: two representatives which become equal in the latter are equal at a common later stage, and injectivity then identifies them already at a stage of the chain above their two indices. Consequently the original colimit contains at least $\kappa$ elements.

The cardinal equality $|E_i|=\pi_0$ is not used to prove this lower bound, apart from the harmless nonemptiness at the first stage. This makes the lemma applicable to a monoid of cardinal \emph{at most} $\pi_0$, as required just above it. One can alternatively start counting at successor stages and omit the initial nonemptiness assumption altogether.

For a left-cancellative monoid, left multiplication by $a$ is surjective exactly when $a$ is a unit: a preimage of $1$ is a right inverse, which is also a left inverse by cancellation. Thus nonunits give precisely the nonsurjective injections needed here. A product containing a nonunit cannot be a unit in such a monoid. The selected chain therefore has every strict transition noninvertible, proving the reduction in 4.15.8.4 to the ordinal $\pi_0^+$. Its printed condition $|\alpha_0|>\alpha_0$ is impossible; the intended condition is $|\alpha_0|>\pi_0$.
\end{dcomment}


\textbf{Théorème 4.15.9.} \textit{Soient $\mathcal{C}$ une petite catégorie, $\pi$ un cardinal tel que $\pi \geqslant$} $\card\mathcal{C}$. \textit{Les conditions suivantes sont équivalentes (et je conjecture qu'elles sont toujours satisfaites) :}
\begin{enumerate}[nosep]
\item[\textit{a)}] \textit{$\mathcal{M} \overset{\text{déf}}{=} \mathrm{Ind}_\pi\mathcal{C}$ est réduit à $\mathcal{M}_\pi$.}
\item[\textit{a$'$)}] \textit{(Si $\mathcal{C}$ est karoubienne.) Tout système inductif $\underline{x} = (x_i)_{i\in I}$ dans $\mathcal{C}$, grand devant $\pi$, est essentiellement constant.}
\item[\textit{a$''$)}] \textit{Pour tout système inductif $(x_i)_{i\in I}$ dans $\mathcal{C}$, grand devant $\pi$, il existe un foncteur cofinal $J \longrightarrow I$, avec $J$ ensemble ordonné grand devant $\pi$, et tel que le système réindexé $(x_j)_{j\in J}$ satisfasse les conditions suivantes :}
\begin{enumerate}[nosep]
\item[\textit{1)}] \textit{Tous les $x_j$ sont égaux à un même élément $x$ de $\mathcal{C}$.}
\item[\textit{2)}] \textit{Il existe un projecteur $p$ dans $x$, tel que l'on ait $\alpha_{j'j} = p$ pour $j < j'$ ($j,j' \in J$), où $\alpha_{j'j} : x_j \longrightarrow x_{j'}$ est la flèche de transition.}
\end{enumerate}
\item[\textit{b)}] \textit{Pour tout $x \in \mathrm{Ob}\,\mathcal{C}$, et pour tout monoïde quotient $M'$ de $M = \mathrm{End}_{\mathcal{C}}(x)$, tel que $M'$ soit un monoïde de monomorphismes (i.e. $uv = uv' \Longrightarrow v = v'$ si $u,v,v' \in M'$), et pour tout système inductif grand devant $\pi$ à valeurs dans $B_{M'}$ (définie par un système d'éléments $(\alpha_{ji})_{i\leqslant j,\ i,j\in I}$ d'un ensemble ordonné $I$ grand devant $\pi$, avec $\alpha_{kj}\alpha_{ji} = \alpha_{ki}$ pour $i \leqslant j \leqslant k$), il existe $i_0 \in I$ tel}
\msp{page 462}
\textit{que $\alpha_{ii_0} = 1$ pour $i \geqslant i_0$, i.e. $\alpha_{ji} = 1$ pour $j \geqslant i \geqslant i_0$.}
\item[\textit{b$'$)}] \textit{$M'$ étant un monoïde de monomorphismes comme dans b) (donc} $\card M' \leqslant \card\mathcal{C} \leqslant$ \textit{$\pi$), soit de plus $\alpha_0$ le plus petit ordinal tel que} $\card\alpha_0 > \pi$, \textit{et soit $I = I_{\alpha_0}$ l'ensemble ordonné des ordinaux $\alpha < \alpha_0$ (lequel ensemble est grand devant $\pi$). Il n'existe} pas \textit{de système inductif dans $B_{M'}$, tel que pour $i < j$ dans $I$, $\alpha_{ji}$ soit un élément non inversible dans $M'$ (donc $u \longmapsto \alpha_{ji}u : M' \longrightarrow M'$ est injectif mais non surjectif).}
\end{enumerate}

% ---- from der-ch18_p181-185.tex
\printedpage{179}
% head: A. Grothendieck : Dérivateurs, Ch. XVIII, Édition des Universités de Montpellier II et Paris VII
% foot: Version 17 décembre 2025

Signalons tout de suite le

\textbf{Corollaire 4.15.10.} \textit{Les conditions équivalentes précédentes sont satisfaites, si $\pi > \card\mathcal{C}$.}

En effet, soit $\alpha_1$ le plus petit ordinal tel que $\card\alpha_1 > \card\mathcal{C}$. Donc $\pi_1 \defeq \card\alpha_1$ est le successeur de $\pi_0 \defeq \card\mathcal{C}$, et on a $\pi \geqslant \pi_1$, donc (si $\alpha_0$ est comme explicité dans b$'$ ci-dessus) $\card\alpha_0 > \pi_1$, donc $\alpha_1 < \alpha_0$. Soit (avec les notations de b$'$), pour $\alpha \in I$ \texttt{[plutôt} $\alpha \in I_{\alpha_0}$\texttt{]}, i.e.\ $\alpha < \alpha_0$,
\[
M_\alpha \;=\; \varinjlim_{i<\alpha} \underbrace{M_i}_{=\,M}\ ,
\]
on voit donc, comme les $\alpha_{ji}$ sont injectifs et non surjectifs, que
\[
\card M_\alpha \;\geqslant\; \card\alpha\ ,
\]
\msp{463}
et que pour $\alpha < \beta$, on a
\[
M_\alpha \;\into\; M_\beta \qquad \text{injectif}\ .
\]
Appliquons ceci au cas où $\alpha = \alpha_1$, donc on trouve
\[
\card M_{\alpha_1} \;\geqslant\; \card\alpha_1 \;=\; \pi_1\ ,
\]
d'autre part, en prenant $\beta = \alpha_1 + 1$, d'où $M_\beta \simeq M$, on trouve une injection
\[
M_{\alpha_1} \;\into\; M\ ,
\]
d'où
\[
\pi_1 \;=\; \card M_{\alpha_1} \;\leqslant\; \card M \;=\; \pi_0\ ,
\]
une contradiction, ce qui établit b$'$.

\dgrblock{dgrb1129}{Invertibility is the invariant condition}
\begin{dcomment}{examples}{p462}{Invertibility is the invariant condition}
Grothendieck wants a diagram to stop producing new ind-objects. This is invariant under an isomorphism of diagrams. Literal equality of its transition arrows with identities is not invariant in that way, and the distinction produces a small counterexample to condition~b) as printed.

Let $G=\{1,g\}$ with $g^2=1$, let $\mathcal C=BG$, and let $I=\pi^+$ for any infinite cardinal $\pi$. Choose $h_\alpha\in G$ by taking $h_\alpha=g$ at successor ordinals and $h_\alpha=1$ at zero and nonzero limit ordinals. Both values occur arbitrarily far out. Define
\[
a_{\beta\alpha}=h_\beta h_\alpha^{-1}\qquad(\alpha\le\beta).
\]
Then $a_{\gamma\beta}a_{\beta\alpha}=a_{\gamma\alpha}$ and $a_{\alpha\alpha}=1$, so this is a diagram in $BG$. The maps $h_\alpha$ give a natural isomorphism from the constant diagram to it. Its ind-object is therefore representable.

There is no $\alpha_0$ such that every $a_{\beta\alpha_0}$ is $1$ for $\beta\ge\alpha_0$: that would make $h_\beta=h_{\alpha_0}$ throughout a tail. Thus condition~b) fails. Condition~b$'$ holds automatically, because a group has no noninvertible elements. Conditions~a), a$'$), and a$''$) hold as well: every filtered group-valued diagram trivializes, and the cofinal-colour argument makes it literally constant on a suitable subsystem. The equivalence asserted in 4.15.9 is therefore false in its printed form even for this two-arrow category.

The appropriate replacement for the tail condition is
\[
\text{there exists }i_0\text{ such that }
a_{ii_0}\in (M')^\times\text{ for every }i\ge i_0.
\]
The cocycle identity then makes all later transitions units. If literal identities are desired, one either changes the trivialization or restricts to a cofinal fibre of the finitely or cardinally many trivializing maps. This is exactly the distinction already handled correctly in 4.15.7.2--5.

The edition has $=1$ in both places, so the transcription is retained. Corollary~4.15.10 inherits that literal defect when it says that \emph{all} the preceding conditions hold. Its proof establishes the intended exclusion of nonunit chains, and its substantive conclusion, Corollary~4.15.11, remains valid. The separate subquotient issue noted above must also be respected when using the monoid reformulation.
\end{dcomment}


\textbf{Corollaire 4.15.11.} \textit{Soit $\mathcal{C}$ une catégorie essentiellement petite, $\pi$ un cardinal tel que $\pi > \card\redop\mathcal{C}$. Soit $\mathcal{M} \defeq \Ind_\pi\mathcal{C}$, alors $\mathcal{M} = \mathcal{M}_\pi$, et c'est une enveloppe de Karoubi de $\mathcal{C}$.}

On peut supposer $\mathcal{C}$ réduite, donc $\card\mathcal{C} < \pi$, et on applique le corollaire précédent.

Ainsi, le résultat litigieux 4.8.13 (p.~212) est finalement prouvé sous une forme légèrement plus faible, avec $\pi > \card\redop\mathcal{C}$ au lieu de $\pi \geqslant \card\redop\mathcal{C}$. Du même coup, les applications de 4.8.13 qu'on a fait jusqu'à présent (notamment le lemme 4.15.3).

\msp{464}

\dgrblock{dgrb1130}{Why a larger filteredness bound leaves only retracts}
\begin{dcomment}{exposition}{p463}{Why a larger filteredness bound leaves only retracts}
The counting proof reaches the result needed for accessibility: a sufficiently small category acquires no objects under sufficiently filtered colimits except its missing retracts. We first spell out the source's strict inequality. Write $\mu=|\operatorname{Mor}\mathcal C|<\pi$ and first suppose $\mu$ infinite. A chain of nonunits indexed by $\pi^+$, in any left-cancellative subquotient monoid $M$ of cardinal at most $\mu$, would give a partial colimit at $\mu^+$ with at least $\mu^+$ elements. That partial colimit injects into the next stage, a copy of $M$. The contradiction is
\[
\mu^+\le \left|\varinjlim_{\alpha<\mu^+}M_\alpha\right|
\le |M|\le\mu.
\]
The equalities in the source's last display are unnecessary; these inequalities are the valid argument. If $\mu$ is finite, the assertion is immediate because an injective self-map of a finite set is surjective.

There is also a direct presheaf verification which explains the resulting retract. Let $X$ be a $\pi^+$-filtered colimit of representables and let $\mathcal E_X$ be its category of elements, with objects $(c,x)$ for $x\in X(c)$. A family of at most $\pi$ such elements, including at most $\pi$ relations between them, factors through one stage of the presentation: representables preserve the colimit on Hom, and $\pi^+$-filteredness makes the choices and relations simultaneous. Hence $\mathcal E_X$ is $\pi^+$-filtered.

If $\mathcal E_X$ had more than $\mu$ objects, select $\mu^+$ of them. Their cocone would have distinct underlying arrows of $\mathcal C$: an arrow $f:c\to d$ and an element $y\in X(d)$ determine its source element uniquely as $X(f)(y)$. This contradicts $|\operatorname{Mor}\mathcal C|=\mu$. It follows that $\mathcal E_X$ has at most $\mu$ objects and at most $\mu^2$ arrows. This bound is at most $\pi$, including the finite case, because $\pi$ is infinite.

Apply its filteredness to its entire identity diagram, obtaining a cocone to $(d,y)$. Write its component at $(c,x)$ as $u_{c,x}:c\to d$. These components define a natural transformation
\[
s:X\longrightarrow h_d,\qquad s_c(x)=u_{c,x}.
\]
Naturality is the cocone equation. The element $y$ defines $t:h_d\to X$, and the fact that every component lies in $\mathcal E_X$ says
$t_cs_c(x)=X(u_{c,x})(y)=x$.
Thus $ts=1_X$: $X$ is a retract of a representable. Conversely every split idempotent image occurs as a sufficiently filtered colimit by the constant-object construction in the worked example, using the desired regular ordinal as index. This proves
\[
\operatorname{Ind}_{\pi}(\mathcal C)\simeq
\operatorname{Kar}(\mathcal C)
\qquad\text{when }\pi>|\operatorname{Mor}\mathcal C|.
\]
The same proof sharpens the bound when each Hom set has cardinal less than $\pi$ and $|\operatorname{Ob}\mathcal C|\le\pi$. A family of $\pi$ distinct elements of one $X(c)$ would force $\pi$ distinct arrows in one $\operatorname{Hom}(c,d)$, which is impossible. Hence every $|X(c)|<\pi$, the category of elements has at most $\pi$ arrows, and its identity diagram again has a cocone. This covers some equality cases, without assuming that it covers all of them.
\end{dcomment}


\dgrblock{dgrb1131}{Small accessible categories and Cauchy completion}
\begin{dcomment}{later}{p463}{Small accessible categories and Cauchy completion}
The conclusion belongs to a broader description of accessibility by completion. Makkai and Par\'e's 1989 monograph develops accessible categories through categorical model theory and mixed sketches \citep{MakkaiPare1989}. The result here is especially close to the Cauchy-completion side of that theory: retracts of representables are the objects left after the filteredness bound is raised sufficiently.

A modern ordinary formulation is \citet[Lemmas~2.5--2.6 and Corollary~2.8]{Tendas2022Continuity}: if the Hom sets have cardinal less than $\beta$, a $\gamma$-flat presheaf for $\gamma>\beta$ is small enough to be Cauchy, provided the object set also lies below $\gamma$. This makes the two cardinal bounds explicit. The paper's Theorem~3.7 extends the equivalence between small accessibility and Cauchy completeness to enriched categories.

For $\infty$-categories, \citet[Lemma~5.4.2.4 and Corollaries~5.4.3.5--5.4.3.6]{Lurie2009HTT} identifies compact objects in an ind-completion with the idempotent completion, proves the sufficiently-large-bound collapse, and characterizes small accessible $\infty$-categories by idempotent completeness. Coherent idempotents replace the ordinary projectors. These statements preserve Grothendieck's conclusion while expressing the required redundancy through the sizes of mapping spaces.
\end{dcomment}


\dgrblock{dgrb1132}{The strict bound and the boundary question}
\begin{dcomment}{status}{p463}{The strict bound and the boundary question}
Grothendieck finishes with the assertion he has actually proved: the bound may be chosen strictly larger than the reduced cardinal of the category. That is sufficient for the qualitative assertion that a small idempotent-complete category is accessible. It is also sufficient for earlier applications in which the chosen cardinal can be enlarged.

The stronger assertion at $\pi=|\operatorname{Mor}\mathcal C|$ must be kept separate. The checked modern references give a strict separation between the local Hom bound and the filteredness bound; their general statements do not automatically remove this boundary. The preceding proofs settle it when all Hom sets have cardinal less than $\pi$, and for categories with two-sided cancellative endomorphism monoids. We have not settled the unrestricted equality case, and did not find an exact resolution of that sharper formulation in the sources consulted. This records the limit of this commentary, rather than declaring a new open problem.

There is consequently no need to weaken the valid conclusion of 4.15.11, and no justification for silently strengthening its cardinal inequality. The audit also records the two logically separate defects in 4.15.9: the false literal identity condition, and the use of quotients where the preceding construction supplies subquotients.
\end{dcomment}


\subsection*{4.16\quad Exemples et questions}

\textbf{4.16.0.} J'avais cru que tout foncteur “raisonnable” (ou plus précisément, tout foncteur “logiquement accessible”) entre catégories accessibles devrait être accessible. Ce n'est malheureusement pas vrai, puisqu'un composé de deux foncteurs Hom
\[
\Ens \longrightarrow \Ens\ , \qquad\quad E \longmapsto \underbrace{\Hom(\Hom(E,I),J)}_{J^{(I^E)}}
\]
\printedpage{180}
% head: A. Grothendieck : Dérivateurs, Ch. XVIII, Édition des Universités de Montpellier II et Paris VII
% foot: Version 17 décembre 2025
est non accessible, si $I$, $J$ sont tous deux de cardinal $\geqslant 2$. Prenant p.ex.\ $I = J = \{0,1\}$, on trouve le foncteur
\[
E \longmapsto \mathfrak{P}(\mathfrak{P}(E))\ ,
\]
qui n'est pas accessible.

\textbf{4.16.1.} Deuxième désillusion : j'avais cru que pour toute catégories accessible $\mathcal{M}$, toute % [sic] toute catégories -> toute catégorie, singular required after "toute"
sous-catégorie $\mathcal{N}$ strictement pleine et stable par facteurs directs “raisonnable” (p.ex. “logiquement accessible”) devrait être accessible dans $\mathcal{M}$. Il n'en est rien, même dans le cas particulier de l'image inverse de la sous-catégorie accessible dans Ens
\msp{465}
\[
\underbrace{\Ens^*}_{=\,\{E\,\in\Ens\,\mid\, E\neq\emptyset\}} \quad\subset\ \Ens
\]
par un contrafoncteur accessible \textit{représentable}
\[
\begin{aligned}
\mathcal{M} \quad &\longrightarrow \quad (\Ens)^\circ \\
x \quad &\longmapsto \quad \Hom_{\mathcal{M}}(x,z)\ ,
\end{aligned}
\]
i.e.\ la sous-catégorie ouverte (donc strictement pleine) $\mathcal{M}_z$ de $\mathcal{M}$ formée des objets $x$ tels que $\Hom_{\mathcal{M}}(x,z) \neq \emptyset$, n'est par toujours accessible dans $\mathcal{M}$. Ici, je ne suis pas parvenu à % [sic] n'est par toujours -> n'est pas toujours, misprint "par" for "pas"
un exemple aussi élémentaire (cf.\ appendice) -- j'ai pris $\mathcal{M} = (k\text{-alg.\ comm.})$, $k$ un corps infini et $z$ l'objet initial $k$ de $\mathcal{M}$ -- j'ai trouvé que $\mathcal{M}_z$ n'est pas $\mathrm{L}_\pi$-stable dans $\mathcal{M}$, quelque soit le cardinal $\pi$ % [sic] quelque soit -> quel que soit, standard spelling of the locution
-- mais cet exemple utilisait que $\mathfrak{U} = \mathfrak{U}_1$, plus petit univers contenant $\mathbf{N}$. J'ignore si pour tout autre univers, on peut trouver un corps $k$ de cardinal $\pi \in \mathfrak{U}$ assez grand, de façon que tout $\pi$-ultrafiltre sur un ensemble $E \in \mathfrak{U}$ soit trivial. Peut-être
\msp{466}
un tel exemple, où non la construction du couple $(\mathcal{M},\mathcal{N})$, mais la démonstration du fait que $\mathcal{N}$ n'est pas accessible dans $\mathcal{M}$, fait appel à des hypothèses particulières sur l'univers ambiant, ne doit pas être considéré comme (tout à fait !) “logiquement accessible”. Pour être vraiment probant, il faudrait un exemple de nature plus élémentaire.

\dgrblock{dgrb1133}{A double dual which defeats every rank}
\begin{dcomment}{examples}{p464}{A double dual which defeats every rank}
Grothendieck begins his examples by testing the hope that a familiar formula should automatically define an accessible functor. The double dual is a particularly useful test: each of its two steps is representable with the appropriate variance, but their composite can remember genuinely global information. The issue is visible without cardinal arithmetic beyond the definition of filteredness.

We use a regular cardinal $\kappa$ in the modern convention: a category is $\kappa$-filtered when every diagram with fewer than $\kappa$ arrows has a cocone. Thus a system grand devant $\pi$ in the manuscript is $\pi^+$-filtered. Fix distinct $a,b\in I$ and distinct $c,d\in J$. The functor under discussion is
\[
 D(X)=\operatorname{Set}(\operatorname{Set}(X,I),J),\qquad
 D(u)(\varphi)(g)=\varphi(g\circ u).
\]
This action on maps matters. When $I=J=2$, the notation $\mathfrak P\mathfrak P(X)$ describes its objects, but does not specify its action by itself.

For any infinite regular $\kappa$, choose a set $X$ of cardinal $\kappa$. Its subsets of cardinal less than $\kappa$, ordered by inclusion, form a $\kappa$-directed system with union $X$. Define $\Phi:I^X\to J$ by setting $\Phi(g)=d$ when $g$ is constantly $a$, and $\Phi(g)=c$ otherwise. If $S\subset X$ has cardinal less than $\kappa$, every member of the image of $D(S)\to D(X)$ depends only on $g|_S$. But the constant function $a$ and the function which is $a$ everywhere except at one point of $X\setminus S$, where it is $b$, have the same restriction to $S$ and different values under $\Phi$. Therefore $\Phi$ belongs to none of these images, and
\[
 \operatorname{colim}_{S\in[X]^{<\kappa}}D(S)\longrightarrow D(X)
\]
is not surjective. This proves the assertion for every possible accessibility rank.

The calculation identifies the obstruction more precisely than growth of cardinalities would. An element of the double dual need not have a support of bounded size. Accessibility would supply one uniform bound, because every element would have to come from a stage in the displayed presentation. The failure concerns that uniformity, even though each individual value of the functor is an ordinary set.
\end{dcomment}


% margin: a vertical bar is printed in the left margin beside the three lines from "L_pi-stable dans M pour pi assez grand" to "logiquement accessible) serait accessible dans M"
Il reste la possibilité qu'il y ait un théorème métamathématique qui assure que tout couple $(\mathcal{M},\mathcal{N})$ logiquement accessible, avec $\mathcal{M}$ accessible, et $\mathcal{N}$ sous-catégorie strictement pleine $\mathrm{L}_\pi$-stable dans $\mathcal{M}$ pour $\pi$ assez grand, soit accessible dans $\mathcal{M}$ (\footnote[188]{C'est faux, même pour une sous-catégorie fermée, cf.\ 4.16.7.10 ci-dessous.}). Ceci s'appliquerait en particulier au cas d'une sous-catégorie \textit{fermée} : toute sous-catégorie fermée (qui soit logiquement accessible) serait accessible dans $\mathcal{M}$. L'exemple-test le plus simple serait celui du complémentaire fermé $\mathcal{M}'_z$ de la sous-catégorie ouverte $\mathcal{M}_z$ ci-dessus (\footnote[189]{\textbf{N.B.} elles est accessible dans $\mathcal{M}$ si et seulement si $\mathcal{M}_z$ l'est ! % [sic] elles est -> elle est, agreement with singular verb
}), i.e.\ la catégorie des $x \in \Ob\mathcal{M}$ tels que $\Hom_{\mathcal{M}}(x,z) = \emptyset$, ou encore
\msp{467}
l'image inverse de la sous-catégorie fermée accessible
\[
e \;=\; \{\emptyset\} \;\subset\; (\Ens)
\]
réduite à $\{\emptyset\}$, par le foncteur représentable $x \longmapsto \Hom(x,z) : \mathcal{M} \longrightarrow \Ens$. Il se pourrait donc qu'une telle catégorie soit toujours accessible dans $\mathcal{M}$. Mais j'ai été incapable de le
\printedpage{181}
% head: A. Grothendieck : Dérivateurs, Ch. XVIII, Édition des Universités de Montpellier II et Paris VII
% foot: Version 17 décembre 2025
prouver, même dans le cas particulier ci-dessus, où $\mathcal{M} = (k\text{-alg.\ comm.})$, $k$ étant un corps commutatif infini. La difficulté est (dans tous les cas) la suivante : trouver, pour $z$ fixé, un cardinal $\pi > \pi_0$ tel que pour tout $x \in \Ob\mathcal{M}$, et tout système (ou \textit{un} système) $\pi_0$-$\pi$-adapté $(x_i)_{i\in I}$ de limite inductive $x_0$ ($\pi_0$ un cardinal donné, utile pour $\mathcal{M}$), \textit{si} $\Hom_{\mathcal{M}}(x,z)$ ($\simeq \varprojlim \Hom_{\mathcal{M}}(x_i,z)$) est vide, alors $\exists\ i_0 \in I$ tel que $\Hom_{\mathcal{M}}(x_{i_0},z)$ soit déjà vide. Dans l'exemple $\mathcal{M} = (k\text{-alg.\ comm.})$, il faut donc prouver que l'on peut trouver un cardinal $\pi$
tel que pour tout $k$-algèbre $A$ avec $\Hom_{\mathcal{M}}(A,k) = \emptyset$, il existe une sous-$k$-algèbre $A_i$ de % [sic] pour tout k-algèbre -> pour toute k-algèbre, "algèbre" is feminine
cardinal $\leqslant \pi$, telle que $\Hom_k(A_i,k) = \emptyset$ (\footnote[190]{C'est faux, cf.\ annotation marginale p.~466.}). J'ai tendance à croire que ça doit être vrai (et même sans
\msp{468}
supposer que $k$ soit un corps, et aussi en prenant la condition $\Hom_{\mathcal{M}}(A,Z) = \emptyset$, où $Z$ est un objet quelconque de $\mathcal{M}$). Ce genre de questions semble étonnement délicat \dots % [sic] étonnement -> étonnamment, adverb required

\dgrblock{dgrb1134}{A point requires compatible choices}
\begin{dcomment}{exposition}{p466}{A point requires compatible choices}
Grothendieck next moves the same difficulty into a category of algebraic objects. For a fixed $z\in\mathcal M$, the subcategory
\[
 \mathcal T_z=\{x:\Hom_{\mathcal M}(x,z)\ne\varnothing\}
\]
is a sieve: an arrow $y\to x$ followed by a point $x\to z$ gives a point of $y$. Its complement $\mathcal S_z$ is a cosieve. Both are replete and closed under retracts. These are the meanings of open and closed throughout the closing pages; no topology on a space of objects is being introduced.

The representable contrafunctor sends every existing colimit to a limit of sets,
\[
 \Hom(\operatorname{colim}_i x_i,z)
 \cong\lim_i\Hom(x_i,z).
\]
The existence of an element in every set on the right need not provide a compatible family. Consequently, preservation of these limits by the representable functor does not show that $\mathcal T_z$ is closed under the corresponding colimits. Describing $\Hom(-,z)$ as an accessible functor to $\mathsf{Set}^{op}$ means preservation of sufficiently filtered colimits into that target; it does not assert that $\mathsf{Set}^{op}$ is an accessible category.

The appendix supplies the concrete example with $\mathcal M=\mathsf{CAlg}_k$ and $z=k$. Its hypothesis on complete ultrafilters is part of the example. Under that hypothesis, for every infinite $\pi\geq|k|$ there is an algebra $A$ with no map to $k$, although every subalgebra of cardinal at most $\pi$ has such a map. The explicit reduced-product construction is proved in the appendix commentary. Since the poset of these subalgebras is $\pi^+$-directed, this proves that $\mathcal T_k$ is closed under no sufficiently large filtered-colimit class. Proposition~4.16.7.8 then transfers the failure of accessibility to its complementary cosieve $\mathcal S_k$.

The distinction explains the later marginal corrections. The no-point class is automatically closed under every nonempty colimit: one structure map from an object with no point would contradict a point on the colimit. Nevertheless, that closure does not furnish a small family of generators. The missing condition is eventually isolated in Corollary~4.16.7.3 as a small weakly initial family in the cosieve.
\end{dcomment}


\textbf{4.16.2.} Il est possible, plus généralement, que la question suivante ait une réponse affirmative : Soit
\[
f \;:\; \mathcal{M} \longrightarrow \mathcal{N}^\circ
\]
un foncteur accessible, avec $\mathcal{M}$, $\mathcal{N}$ catégories accessibles, et soit $\underline{T}$ une sous-catégorie de $\mathcal{N}$ accessible dans $\mathcal{N}$, et telle qu'il existe un cardinal $\pi$ ayant les propriétés suivantes : $\mathcal{N}^\circ$ satisfait $\mathrm{L}_\pi$ (i.e.\ $\mathcal{N}$ stable par $\varprojlim$ grandes devant $\pi_1$). Alors $\underline{S} \defeq f^{-1}(\underline{T})$ est-elle accessible dans $\mathcal{M}$ (\footnote[191]{Faux, avec $\mathcal{N} = \Ens$ et un foncteur représentable $\mathcal{M} \longrightarrow \Ens^\circ$, prenant $I = \{\emptyset\}$, cf.\ annotations marginales p.~466, 467.}) ? Ceci généralise le cas $\mathcal{N} = (\Ens)$, $\underline{T} = \{\emptyset\}$, $f$ foncteur représenté par $z \in \Ob\mathcal{M}$, envisagé tout à l'heure. Avec le même type d'exemple, mais en prenant $\mathcal{N} = (\Ens)^*$
\msp{469}
au lieu de $\mathcal{N} = \{\emptyset\}$, on voit que la condition supplémentaire de stabilité sur $\underline{T}$ n'est pas
superflue. (\textbf{N.B.} Elle sont à assurer que $\underline{S}$ est $\mathrm{L}_\pi$-stable dans $\mathcal{M}$ pour $\pi$ grand.) % [sic] Elle sont -> Elles sont, agreement with plural verb

\textbf{4.16.3.} Questions apparentée : j'avais espéré que pour un foncteur accessible entre catégories accessibles % [sic] Questions apparentée -> Question apparentée, agreement
\[
f \;:\; \mathcal{N} \longrightarrow \mathcal{M}\ ,
\]
la $\natural$-image essentielle, formée des $y \in \Ob\mathcal{M}$ qui sont isomorphes à un facteur direct d'un objet de la forme $f(x)$, avec $x \in \Ob\mathcal{N}$, serait accessible dans $\mathcal{N}$. L'exemple ci-dessus (nullement élémentaire !) de la sous-catégorie $\mathcal{M}_z$ de $\mathcal{M}$ (formée des $x \in \Ob\mathcal{M}$ tels que $\Hom_{\mathcal{M}}(x,z) \neq \emptyset$) montre qu'il n'en est rien. En effet, $\mathcal{M}_z$ est l'image essentielle, et aussi l'image $\natural$-essentielle, du foncteur
\[
\mathcal{M}/z \longrightarrow \mathcal{M}\ ,
\]
or $\mathcal{M}/z$ est accessible, et le foncteur précédent accessible, si $\mathcal{M}$ est accessible (vérification par AQT \texttt{[âne qui trotte]}, cf.\ plus bas, énoncé plus général).

\msp{470}

On a l'impression que la raison pour laquelle on a de tels contre-exemples, c'est que l'on a des $\natural$-images essentielles qui ne sont $\mathrm{L}_\pi$-stables dans $\mathcal{M}$ pour aucun cardinal $\pi \in \mathfrak{U}$. Cela amène à modifier la conjecture ainsi : soit $\pi$ un cardinal tel que $\mathcal{M}$ satisfasse à $\mathrm{L}_\pi$,
\printedpage{182}
% head: A. Grothendieck : Dérivateurs, Ch. XVIII, Édition des Universités de Montpellier II et Paris VII
% foot: Version 17 décembre 2025
et soit $\underline{S}(\pi)$ la plus petite sous-catégorie strictement pleine de $\mathcal{M}$, $\mathrm{L}_\pi$-stable dans $\mathcal{M}$, et qui contienne l'image essentielle de $f$. (On pourrait appeler $\underline{S}(\pi)$ la $\mathrm{L}_\pi$-image essentielle.) Alors $\underline{S}(\pi)$ est accessible dans $\mathcal{M}$. De plus, la famille (décroissante) des $\underline{S}(\pi)$, pour $\pi$ variable, est stationnaire, i.e.\ il existe un cardinal $\pi$ tel que $\pi' \geqslant \pi$ implique $\underline{S}(\pi') = \underline{S}(\pi)$. Il est clair alors que $\underline{S}(\pi)$ est la plus petite sous-catégorie de $\mathcal{M}$ qui soit accessible dans $\mathcal{M}$ et qui contienne l'image essentielle de $f$. La même question se pose pour un foncteur $g^\circ : \mathcal{N}^\circ \longrightarrow \mathcal{M}$, où $g : \mathcal{N} \longrightarrow \mathcal{M}^\circ$ est un foncteur accessible, $\mathcal{N}$ et $\mathcal{M}$ accessibles. Je conjecture que la réponse est encore affirmative.

\msp{471}

\dgrblock{dgrb1135}{Filtered closures need not settle at a final rank}
\begin{dcomment}{development}{p470}{Filtered closures need not settle at a final rank}
Grothendieck tries to repair the failure of essential images by closing them under sufficiently filtered colimits. This does recover an accessible category at suitable ranks. The additional prediction that these closures eventually stop changing is stronger, and the appendix gives a way to disprove it.

\textbf{The closure at a common strong rank.}\ours\
Let $F:\mathcal A\to\mathcal B$ be an accessible functor. Choose a regular $\kappa$ such that both categories are $\kappa$-accessible and $F$ preserves both $\kappa$-filtered colimits and $\kappa$-presentable objects. Let $\mathcal C$ be the small full subcategory of $\mathcal B_\kappa$ on one chosen representative of each isomorphism class in $F(\mathcal A_\kappa)$, with no additional objects. Then the smallest replete full $\kappa$-filtered-colimit-closed subcategory $\mathcal S_\kappa$ containing $F(\mathcal A)$ is
\[
 \mathcal S_\kappa\simeq\operatorname{Ind}_\kappa(\mathcal C),
\]
where this Ind notation uses the modern regular-cardinal convention. In particular, it is accessible and accessibly embedded in $\mathcal B$.

Indeed, extend $\mathcal C\hookrightarrow\mathcal B$ by $\kappa$-filtered colimits. For presentations $X=\operatorname{colim}_i c_i$ and $Y=\operatorname{colim}_j d_j$, compactness gives
\[
 \Hom_{\mathcal B}(X,Y)
 \cong\lim_i\operatorname{colim}_j\Hom_{\mathcal C}(c_i,d_j).
\]
This is precisely the Ind hom formula, so the extension is fully faithful and its image is closed under $\kappa$-filtered colimits. Every $F(a)$ lies in that image, because a presentation of $a$ by $\mathcal A_\kappa$ is preserved by $F$. Conversely, any closed full subcategory containing $F(\mathcal A)$ contains $\mathcal C$ and all its $\kappa$-filtered colimits. This proves both the formula and minimality. The existence of arbitrarily large common strong ranks is the usual uniformization theorem; its formulation is recalled in the introduction of \citet{PareRosicky2013}. The proof here applies that theorem rather than asserting the formula at every initially chosen rank.

\textbf{Failure of eventual stabilization.}\ours\
Work under the universe hypothesis of the appendix: some infinite $\pi_0$ excludes all nonprincipal $\pi_0^+$-complete ultrafilters on small sets. Choose an infinite field $k$ with $|k|\ge\pi_0$, put $\mathcal B=\mathsf{CAlg}_k$, and take $F:\mathcal B/k\to\mathcal B$. Its essential image is the sieve $\mathcal T$ of algebras admitting a $k$-point. For every regular rank $\kappa$, the appendix construction, with a cardinal bound at least $\kappa$, supplies a $\kappa$-filtered diagram in $\mathcal T$ whose colimit $A$ has no $k$-point. Thus $A\in\mathcal S_\kappa\setminus\mathcal T$.

Choose a larger regular $\mu$ for which $A$ is $\mu$-presentable. The full subcategory
\[
 \mathcal Z_A=\{B:\Hom(A,B)=\varnothing\}
\]
contains $\mathcal T$: an arrow $A\to B\to k$ would be a forbidden point. It is closed under $\mu$-filtered colimits, because $\Hom(A,-)$ preserves them and a filtered colimit of empty sets is empty. By minimality, $\mathcal S_\mu\subseteq\mathcal Z_A$. But $A\notin\mathcal Z_A$, as its identity is a morphism. Consequently
\[
 A\in\mathcal S_\kappa\setminus\mathcal S_\mu.
\]
There can be no final stationary rank. This also works when both ranks are chosen among common strong ranks, where the individual closures are accessible by the preceding argument.

The two assertions in the manuscript thus have different answers. We checked the accessible-colimit literature cited here and the large-cardinal image theorem discussed in the appendix; we did not find this particular application to the decreasing closures in those sources. The argument is presented as a development of these pages, without a claim of priority. The next development treats the contravariant stabilization assertion separately, using the same reduced products.
\end{dcomment}


\dgrblock{dgrb1136}{The contravariant closures also need not stabilize}
\begin{dcomment}{development}{p470}{The contravariant closures also need not stabilize}
The contravariant version of the conjecture can be tested with the very products used in the appendix. The following argument concerns stationarity; it does not assume, or establish, accessibility of each individual closure.

\textbf{Conditional failure of contravariant stationarity.}\ours\
Under the universe hypothesis of Proposition~4 of the appendix, there are accessible categories $\mathcal N,\mathcal M$ and a functor $g:\mathcal N\to\mathcal M^{op}$ preserving all colimits for which the $\mathrm L_\pi$-closures of the image of $g^{op}$ never become stationary.

Choose an infinite field $k\in\mathfrak U$ with $|k|\ge\pi_0$, where $\pi_0$ is the cardinal in that universe hypothesis. Take $\mathcal N=\mathrm{Set}$, $\mathcal M=\mathsf{CAlg}_k$, and $g(E)=k^E$, with pullback of functions on arrows. The product universal property gives
\[
 \Hom_k(A,k^E)\cong\mathrm{Set}(E,\Hom_k(A,k)).
\]
Consequently $g$ takes every colimit of sets to the corresponding limit of algebras, and therefore preserves colimits as a functor to $\mathcal M^{op}$. This is the accessibility convention required in the manuscript; no accessibility of $\mathcal M^{op}$ is being asserted. Write $\mathcal S(\pi)$ for the smallest full replete subcategory containing all $k^E$ and closed under $\pi^+$-filtered colimits in $\mathcal M$.

Fix an infinite $\pi$, choose $|E|>\pi$, and put
\[
 B_\pi= k^E/I_\pi,
 \qquad I_\pi=\{f:|\operatorname{supp}f|\le\pi\}.
\]
The ideal is proper, and restriction gives
\[
 B_\pi\cong\varinjlim_{F\in[E]^{\le\pi}} k^{E\setminus F}.
\]
The indexing poset is $\pi^+$-directed, so $B_\pi\in\mathcal S(\pi)$. It has no $k$-point. Indeed, the appendix's ultrafilter calculation says that every $k$-algebra map $k^E\to k$ is evaluation at a coordinate: the corresponding ultrafilter is $|k|^+$-complete, hence principal under the hypothesis. A point of $B_\pi$ would therefore give evaluation at some $e\in E$. But the characteristic function of $\{e\}$ maps to zero in $B_\pi$ and to one under that evaluation, a contradiction.

Choose $\pi'>\max\{\pi,|k|,|B_\pi|\}$. Then $B_\pi$ is $\pi'^+$-presentable in $\mathsf{CAlg}_k$: use one generator for each element and at most $\pi'$ relations recording the algebra operations and constants. Let $0$ denote the zero algebra, which is terminal, and consider the full replete class
\[
 \mathcal Z=\{0\}\ \cup\
 \{A:\Hom_k(B_\pi,A)=\varnothing\}.
\]
Every product $k^J$ belongs to $\mathcal Z$. If $J\ne\varnothing$, a map $B_\pi\to k^J$ followed by a coordinate projection would be a forbidden $k$-point; if $J=\varnothing$, the product is $0$.

The class $\mathcal Z$ is closed under $\pi'^+$-filtered colimits. For a diagram with no zero value, presentability of $B_\pi$ makes the hom set into its colimit empty. If a zero value occurs, its zero-valued indices form a nonempty cofinal cosieve, since a unital map from $0$ to an algebra forces that algebra to be zero. The colimit is then zero. By minimality, $\mathcal S(\pi')\subseteq\mathcal Z$. However, $B_\pi$ is nonzero and has an identity map, so $B_\pi\notin\mathcal Z$. Thus
\[
 B_\pi\in\mathcal S(\pi)\setminus\mathcal S(\pi').
\]
As $\pi$ was arbitrary, there is no final stationary rank. The zero-algebra case is necessary because the contravariant image includes the empty product. This supplies a separate negative answer to the stabilization assertion for contravariant images, under the stated universe hypothesis. We did not find this application of the appendix's calculation in the literature checked; no claim of priority is intended.
\end{dcomment}


% margin: a hand-drawn curved stroke is printed in the left margin beside the last three lines of the following paragraph ("pleine M de C^ qui soit ..." to "sans pourtant être accessible.")
Soit $\mathcal{M}$ une catégorie satisfaisant la condition L, et admettant une petite sous-catégorie pleine strictement génératrice $\mathcal{C}$, formée d'objets accessibles. On a donc un foncteur pleinement fidèle et accessible
\[
\mathcal{M} \;\into\; \mathcal{C}^\wedge\ ,
\]
$\mathcal{M}$ $\mathrm{L}_\pi$-stable dans $\mathcal{C}^\wedge$ pour $\pi$ convenable. Si $\overline{\mathcal{M}}$ est l'image essentielle, de sorte que $\mathcal{M} \simeq \overline{\mathcal{M}}$, c'est donc une sous-catégorie strictement pleine $\mathrm{L}_\pi$-stable dans $\mathcal{C}^\wedge$. Donc on s'attend à ce qu'elle soit accessible dans $\mathcal{C}^\wedge$ (si la situation décrite est “logiquement accessible”), ou ce qui revient au même, que $\mathcal{M}$ soit accessible. (Comparer 4.8.1, p.~187.) Mais il ne doit pas être difficile (\footnote[192]{OK en prenant un \textit{fermé} $\underline{S}$ non accessible dans $\mathcal{M}$ d'une catégorie accessible $\mathcal{M}$, qu'on plonge dans un $\mathcal{C}^\wedge$. Mais ces $\underline{S}$ n'a \textit{pas} une petite famille strictement génératrice ! % [sic] ces S n'a -> ce S n'a, plural determiner with singular verb
}), en procédant à un nombre $\pi_{\mathfrak{U}}$ de choix (où $\pi_{\mathfrak{U}} = \card\mathfrak{U}$), de faire quand même des exemples (“logiquement inaccessibles”) d'une sous-catégorie strictement pleine $\mathcal{M}$ de $\mathcal{C}^\wedge$ qui soit $\mathrm{L}_\pi$-stable, mais non accessible dans $\mathcal{C}^\wedge$. Cette $\mathcal{M}$ admet donc une petite sous-catégorie strictement génératrice formée d'objets accessibles (\footnote[193]{Ce n'est pas clair, il faudrait pour cela que le foncteur d'inclusion soit un adjoint à gauche !}) et satisfait à $\mathrm{L}_\pi$, sans pourtant être accessible.

\msp{472}

\dgrblock{dgrb1137}{A dense family of presentables need not suffice}
\begin{dcomment}{development}{p471}{A dense family of presentables need not suffice}
The restricted Yoneda argument at the start of this paragraph is correct. The difficulty in footnotes~192--193 is that an arbitrary full subcategory of a presheaf category need not contain its representables. We can repair the proposed counterexample by adjoining their retracts to a suitable nonaccessible cosieve. This gives an example satisfying the stronger hypothesis actually stated here, under the same universe assumption as the appendix.

\textbf{Conditional counterexample.}\ours\
Assume that some infinite $\pi_0\in\mathfrak U$ has the property that every $\pi_0^+$-complete ultrafilter on a $\mathfrak U$-small set is principal. There is a locally small category $\mathcal M$ with $\kappa$-filtered colimits and a small dense full subcategory consisting of $\kappa$-presentable objects, for one regular $\kappa$, although $\mathcal M$ is not accessible.

Choose an infinite field $k\in\mathfrak U$ with $|k|\ge\pi_0$. Put $\mathcal A=\mathsf{CAlg}_k$, and let $\mathcal C$ be the small full category of polynomial algebras $k[x_1,\ldots,x_n]$, for finite $n$. The restricted nerve
\[
 N:\mathcal A\longrightarrow\mathcal P=\mathrm{Set}^{\mathcal C^{op}}
\]
is fully faithful and preserves filtered colimits. Indeed, its value on the algebra with $n$ generators is the $n$th power of the underlying set. Naturality with respect to polynomial substitutions says exactly that a transformation between two such nerves preserves the algebra operations and constants. Finite powers commute with filtered colimits.

Set $Z=N(k)$ and
\[
 \mathcal S=\{X\in\mathcal P:\Hom_{\mathcal P}(X,Z)=\varnothing\}.
\]
This is a full replete cosieve, closed under every nonempty colimit. It is not accessible in $\mathcal P$. Otherwise its inverse image under the accessible functor $N$ would be accessible in $\mathcal A$, contrary to Proposition~4 of the appendix together with Corollary~4.16.7.9. Here the inverse-image theorem applies without a strict-pullback defect: the inclusion of a full replete subcategory is an isofibration, so its strict pullback is equivalent to the accessible pseudopullback. The latter is constructed by products, inserters and equifiers, using the closure theorem discussed in \citet[Appendix~A]{ChingRiehl2014}. Full faithfulness identifies this inverse image with the algebras having no $k$-point.

Let $\mathcal T$ consist of the retracts of representable presheaves, including all isomorphic copies. It is essentially small, since retracts are classified by idempotents on representables, and is closed under retracts. Every member of $\mathcal T$ maps to $Z$: each representable has such a map by evaluating its polynomial generators at zero, and a retract inherits one. Consequently $\mathcal S\cap\mathcal T=\varnothing$.

We verify the required filtered closure of $\mathcal T$ explicitly. Choose an infinite $\mu\ge|\operatorname{Mor}\mathcal C|$ and put $\kappa=\mu^{++}$. Every object of $\mathcal T$ has at most $\mu$ elements in total across its components. If $X$ is the presheaf colimit of a $\kappa$-filtered diagram in $\mathcal T$, then $|X(c)|\le\mu$ for each $c$: otherwise $\mu^+$ distinct elements would have representatives at one common stage, since $\mu^+<\kappa$, contradicting that stage's bound. Thus $X$ also has at most $\mu$ elements in total. Its category of elements has at most $\mu$ arrows, so its canonical expression as a colimit of representables makes $X$ $\kappa$-presentable. Applying presentability to its original filtered presentation factors $1_X$ through a stage. Hence $X$ is a retract of that stage and belongs to $\mathcal T$. Notice the strict inequality $\mu^+<\kappa$ in this argument.

Now take the full replete subcategory
\[
 \mathcal M=\mathcal S\cup\mathcal T\subseteq\mathcal P.
\]
It is closed under $\kappa$-filtered colimits. A diagram taking values only in $\mathcal T$ was just handled. If it has a value in $\mathcal S$, its $\mathcal S$-valued indices form a nonempty cosieve in the filtered indexing category, hence a cofinal filtered subcategory; its colimit belongs to $\mathcal S$. Thus these colimits in $\mathcal M$ are computed in $\mathcal P$.

The representables form a small dense full subcategory of $\mathcal M$. In fact, the restricted nerve is precisely the full inclusion $\mathcal M\hookrightarrow\mathcal P$, by Yoneda. This proves genuine density, rather than just the weaker generator property. Each representable is $\kappa$-presentable in $\mathcal M$, because its hom functor is evaluation and the relevant colimits are computed pointwise.

Finally, suppose $\mathcal M$ were accessible. Choose a regular $\lambda\ge\kappa$ at which it is accessible, with a small family $\mathcal G$ of $\lambda$-presentable generators. In a $\lambda$-filtered presentation of $X\in\mathcal S$ by $\mathcal G$, not all stages can belong to $\mathcal T$, since $\mathcal T$ is closed under such colimits. The $\mathcal S$-stages are therefore nonempty and cofinal. They present $X$ using the small family $\mathcal G\cap\mathcal S$, whose members remain $\lambda$-presentable in $\mathcal S$. This makes $\mathcal S$ accessible; its inclusion into $\mathcal P$ preserves filtered colimits, contradicting the earlier conclusion. This is the mechanism of Corollary~4.16.7.9 in the present example.

The construction uses the appendix's explicit universe hypothesis, and does not establish failure in every universe. It repairs the dense-family issue without a universe-sized sequence of arbitrary choices. We did not find this particular adjoining-of-retracts construction in the literature checked for this commentary; it is supplied as a development of the manuscript, without a claim of priority.
\end{dcomment}


\textbf{4.16.4} (\footnote[194]{descendre après 4.16.5.}). Soit $\mathcal{M}$ une catégorie telle que $\mathcal{M}$ et $\mathcal{M}^\circ$ soient accessibles. Je conjecture qu'alors $\mathcal{M}$ est essentiellement petite. Plus généralement, supposons $\mathcal{M}$ accessible, et que $\mathcal{M}$ admette une petite sous-catégorie qui soit \textit{strictement co}génératrice. Alors $\mathcal{M}$ est-elle essentiellement petite ? Plus généralement, si $\mathcal{M}$ admet une petite sous-catégorie strictement génératrice, et aussi une petite sous-catégorie strictement cogénératrice, i.e.\ si $\mathcal{M}$ et $\mathcal{M}^\circ$ sont toutes deux équivalentes à des sous-catégories pleines d'une $\mathcal{C}^\wedge$, avec $\mathcal{C}$ petite, alors $\mathcal{M}$ est-elle essentiellement petite ?

\textbf{4.16.5.} Soit $I$ une petite catégorie, $I_\bullet$ le cône droit sur $I$ (obtenu en rajoutant un objet final $e$), $\mathcal{M}$ une catégorie accessible,
\[
\underline{\Hom}_{\square}(I_\bullet,\mathcal{M}) \;\into\; \underline{\Hom}(I_\bullet,\mathcal{M})
\]
la sous-catégorie strictement pleine de la catégorie (accessible) $\underline{\Hom}(I_\bullet,\mathcal{M})$, formée des foncteurs $f$ qui font de $f(e)$ la limite inductive dans $\mathcal{M}$ des $f(i)$ ($i \in \Ob I$). Cette sous-catégorie est $\mathrm{L}_\pi$-stable pour $\pi$ grand, donc il y a lieu
\msp{473}
de conjecturer qu'elle est accessible dans $\underline{\Hom}(I_\bullet,\mathcal{M})$. $\Big($ Je rappelle qu'on a vu dans 4.14 que c'est OK si les $\varinjlim$ de type $I$ sont représentables dans $\mathcal{M}$. Et qu'une réponse affirmative, dans le cas où
\printedpage{183}
% head: A. Grothendieck : Dérivateurs, Ch. XVIII, Édition des Universités de Montpellier II et Paris VII
% foot: Version 17 décembre 2025
\[
\begin{tikzcd}[row sep=small, column sep=small]
 & 1 \\
I = 0 \arrow[ur] \arrow[dr] & \\
 & 2
\end{tikzcd}
\quad ,\qquad\quad \text{donc} \qquad\quad
\begin{tikzcd}[row sep=small, column sep=small]
 & 1 \arrow[dr] & \\
I = 0 \arrow[ur] \arrow[dr] & & 3 \\
 & 2 \arrow[ur] &
\end{tikzcd}
\quad ,
\]
implique que la sous-catégorie $\underline{\text{Ép}}(\mathcal{M})$ dans $\underline{\mathrm{Fl}}(\mathcal{M})$ est accessible dans $\underline{\mathrm{Fl}}(\mathcal{M})$.$\Big)$

\dgrblock{dgrb1138}{Isbell's large category seen from both sides}
\begin{dcomment}{examples}{p472}{Isbell's large category seen from both sides}
Grothendieck's three smallness questions ask whether having enough descriptions by small objects from both directions leaves room for a large category. Isbell had already constructed a decisive example. Let $\mathsf{PInj}$ have sets as objects and partial bijections as morphisms: a map $X\dashrightarrow Y$ is a bijection from a subset of $X$ onto a subset of $Y$. Reversing that bijection defines an isomorphism $\mathsf{PInj}^{op}\cong\mathsf{PInj}$. Isbell's example is recorded by \citet[Example~8.7]{AveryLeinster2021}, who attribute the one-object dense and codense subcategory to \citet[Example~1]{Isbell1968Small}. We verify the properties relevant here directly.

First, filtered colimits exist and finite sets are finitely presentable. To construct the colimit of a filtered diagram, adjoin a basepoint to every set and send an undefined value of a partial map to the basepoint. Take the colimit of this diagram of pointed sets. Each structure map is still injective away from the basepoint: if two elements from one stage become equal to the same nonbasepoint element, they become equal at a common later stage, where the transition is injective away from the basepoint. Removing the colimit basepoint therefore gives a partial bijection from each original stage.

For the universal property, a cocone of partial bijections gives a pointed-set map out of this colimit. If two distinct nonbasepoint classes had the same nonbasepoint image, represent them at a common stage; injectivity of that stage's cocone map would make them equal. Hence the induced map is a partial bijection, and its uniqueness follows from the pointed-set universal property. This proves the colimit assertion without assuming that all transitions are everywhere defined.

A partial bijection from a finite set $K$ to this colimit has finitely many defined values. Represent those values at one common stage; their distinct nonbasepoint images ensure that the representatives are distinct and nonbasepoint. Leave the other elements of $K$ undefined. This lifts the map to that stage. Two such lifts which agree in the colimit agree at a later stage, by applying the filtered equality criterion to their finitely many values, including the basepoint. Thus $K$ is finitely presentable. Finally, every set is the filtered union of its finite subsets along total inclusions. There is a small skeleton of finite sets, so $\mathsf{PInj}$ is finitely accessible. Its self-duality proves the same for its opposite.

To see the stronger density assertion, let $D=\{1,2\}$. Its partial-bijection endomorphism monoid has seven elements: the empty map, four singleton-domain maps, and two permutations. A natural transformation
\[
 \mathsf{PInj}(D,X)\longrightarrow\mathsf{PInj}(D,Y)
\]
means a function equivariant for precomposition by these seven endomorphisms. Commutation with the empty map forces the empty map to be preserved. Commutation with the two partial identity maps shows that the image of a map defined only at $1$ is again defined only at $1$, or empty; this defines a partial function $h:X\dashrightarrow Y$. Commutation with the transposition gives the same $h$ in the second coordinate. For distinct $x,x'\in X$, apply the transformation to the map $D\to X$ taking the two points to $x,x'$. Its image must still be a partial bijection. Hence, whenever both $h(x)$ and $h(x')$ are defined, they are distinct. Thus $h$ is a partial bijection. The partial identities now show that the transformation is exactly postcomposition with $h$. Its uniqueness is detected by singleton-domain maps. The restricted nerve at $D$ is therefore fully faithful, proving density of this one-object full subcategory. Self-duality proves codensity.

Sets of distinct cardinalities are not isomorphic in $\mathsf{PInj}$. There are therefore universe-many isomorphism classes, while the dense and codense test category is finite. All three smallness assertions fail. No large-cardinal hypothesis is used in this example, and the conclusion is about accessibility, not about the much stronger assumption of having all small limits or colimits.
\end{dcomment}


\textbf{N.B.} Si $\mathcal{M}$ admet une petite sous-catégorie \textit{co}génératrice, alors la question est justiciable de la conjecture 4.16.2 relative à un foncteur accessible $\mathcal{M} \longrightarrow \mathcal{N}^\circ$, ici $\underline{\Hom}(I_\bullet,\mathcal{M}) \longrightarrow \underline{\Hom}(I^\circ \times \mathcal{C},\Ens)^\circ \dots$

\textbf{4.16.6.} Il reste la question de l'accessibilité relative d'une catégorie $\mathcal{E}$ sur une petite catégorie $\mathcal{C}$. Je pense à présent, grâce à 4.15.9, et son corollaire 4.15.10, que l'énoncé suivant doit “tomber” par AQT \texttt{[âne qui trotte]} :

\textbf{Proposition 4.16.6.1.} \textit{Soit $f : \mathcal{E} \longrightarrow \mathcal{C}$ un foncteur, avec $\mathcal{C}$ petite catégorie karoubienne. On suppose que $f$ satisfait la condition $H' = H$ (1+2+3) de 4.15.1 (page 415), où il suffit H 2, H 3, pour toute flèche $\lambda$} non identique \textit{de $\mathcal{C}$. On suppose de plus que pour une catégorie $\mathcal{C}'$ d'un des deux types qu'on va décrire, et un foncteur $\mathcal{C}' \longrightarrow \mathcal{C}$, $\mathcal{E}' = \mathcal{E} \times_{\mathcal{C}} \mathcal{C}'$ soit accessible, et $\mathcal{E}' \longrightarrow \mathcal{C}'$ accessible.}

\msp{474}

\dgrblock{dgrb1139}{Existing pushouts can form a nonaccessible class}
\begin{dcomment}{development}{p473}{Existing pushouts can form a nonaccessible class}
Grothendieck wants the collection of universal cocones itself to have a manageable presentation. When all diagrams of a fixed shape have colimits, this is true: restriction identifies the category of colimit cocones with the accessible diagram category. The higher-categorical version has exactly this existence hypothesis in \citet[Example~5.4.7.9]{Lurie2009HTT}. The hypothesis becomes consequential when only some diagrams have colimits.

\textbf{A counterexample from the appendix.}\ours\
Assume the appendix's universe hypothesis with cutoff $\pi_0$, and choose an infinite field $k$ with $|k|\ge\pi_0$. There is a finitely accessible category $\mathcal E$ for which the full replete subcategory of pushout squares is not accessible and accessibly embedded in $\mathcal E^\square$.

To construct it, start with $\mathcal A=\mathsf{CAlg}_k$ and adjoin a new object $*$, with
\[
 \Hom_{\mathcal E}(A,*)=\Hom_k(A,k),\qquad
 \Hom_{\mathcal E}(*,A)=\varnothing,\qquad
 \End_{\mathcal E}(*)=\{1\}.
\]
Composition of crossing maps is composition in $\mathcal A$. Thus $\mathcal A$ is the open piece of a collage and $\{*\}$ its closed piece. The new object and the old algebra $k$ are distinct: in particular, there is no arrow from $*$ to $k$.

We first show that $\mathcal E$ is finitely accessible. A filtered diagram with a $*$-valued stage has a cofinal full subcategory of $*$-valued stages. Indeed, a common successor of such a stage and any other stage must again have value $*$, since no arrow leaves $*$ for the open piece. That subcategory is filtered, and its image diagram is constant with identity transitions. Its colimit is $*$. A filtered diagram entirely in $\mathcal A$ has its usual colimit there, which is also a colimit in $\mathcal E$: cocones to $*$ are precisely compatible maps to $k$, and representability identifies them with maps from the algebraic colimit to $k$. These two cases also show that each finitely presentable algebra remains finitely presentable in $\mathcal E$, and that $*$ is finitely presentable. A small skeleton of finitely presentable algebras together with $*$ generates $\mathcal E$ under filtered colimits. Hence $\mathcal E$ is finitely accessible.

For an algebra $A$, form the commutative square
\[
 \begin{tikzcd}
 k \arrow[r] \arrow[d] & A\times k \arrow[d,"\operatorname{pr}_2"]\\
 * \arrow[r,"1"'] & * .
 \end{tikzcd}
\]
Here the right vertical arrow means the crossing map specified by the algebra homomorphism $A\times k\to k$. Every cocone on the upper-left span has vertex $*$, because its leg from $*$ forces that vertex. Its remaining choice is a map $A\times k\to k$; compatibility with the leg from $k$ is automatic for $k$-algebra maps. Since $k$ is a field, the two complementary central idempotents of $A\times k$ show that
\[
 \Hom_k(A\times k,k)
 \cong\Hom_k(A,k)\amalg\{\operatorname{pr}_2\}.
\]
The category of these cocones is discrete, because $*$ has only its identity endomorphism. The displayed canonical cocone is initial precisely when this set is a singleton, or equivalently when $\Hom_k(A,k)=\varnothing$.

The assignment $A\mapsto$ this square is an accessible functor $\mathcal A\to\mathcal E^\square$. Finite products in $k$-algebras commute with filtered colimits, the inclusion of the open piece preserves them by the preceding calculation, and the two constant $*$ components also preserve filtered colimits. If the full replete subcategory of pushout squares were accessible and accessibly embedded, its inverse image under this functor would be so as well. This last assertion is the accessible pseudopullback theorem, applied to a replete full inclusion, which is an isofibration; see \citet[Remark~4.9]{BourkeLack2024}. The inverse image, however, is exactly the nonaccessible cosieve
\[
 \{A:\Hom_k(A,k)=\varnothing\}
\]
proved nonaccessible in the appendix and Proposition~4.16.7.8. This contradiction establishes the claim.

The example also respects the source's observation that these cocones are closed under sufficiently filtered colimits: universal colimits commute with such pointwise colimits. It is the small generating family which fails. We did not find this collage realization in the checked accounts of accessible categories or in the all-colimits theorem cited above; the proof is given in full as a conditional development, without a priority claim.

The separate implication concerning epimorphisms is correct. For $u:A\to B$, the square with top and left edge $u$ and bottom and right edge $1_B$ is a pushout exactly when $u$ is epic: a pair of maps out of $B$ agreeing after $u$ must then agree. This gives the accessible inverse-image argument mentioned in the manuscript. It does not reverse the implication or remove the existence hypothesis from the preceding positive theorem.
\end{dcomment}


$\Big($ \textit{On prend pour $\mathcal{C}'$ la catégorie}
\[
\mathcal{Q} = \left\{ \quad
\begin{tikzcd}[column sep=2.2em]
\xi_0 = \xi^\natural \arrow[r, shift left, "\sigma"] \arrow[dr, "u_0"'] & \xi \arrow[l, shift left, "\rho"] \arrow[d, "u = u_0\rho"] \arrow[loop right, "p"] \\
 & \eta
\end{tikzcd}
\qquad
\begin{array}{rclll}
\rho\sigma & = & \mathrm{id}_{\xi_0} & & \\[2ex]
\sigma\rho & = & p\ , & \text{\textit{d'où}} & \left\{\begin{array}{rcl} p^2 & = & p \\ p\sigma & = & \sigma \\ \rho p & = & \rho \end{array}\right. \\[4ex]
u & = & u_0\rho\ , & \text{\textit{d'où}} & \left\{\begin{array}{rcl} up & = & u \\ u\sigma & = & u_0 \end{array}\right.
\end{array}
\right.
\]
\textit{(trois objets $\xi$, $\xi_0$, $\eta$ et cinq flèches non identiques $\sigma$, $\rho$, $p$, $u_0$, $u$, liés par les relations ci-dessus), ou la catégorie (isomorphe à $Q^\circ$)} % Doubt resolved against PDF leaf 25: Q is ordinary mathematical italic, retained.
\[
\mathcal{Q}' = \left\{ \quad
\begin{tikzcd}
\xi_0 \arrow[r, shift left, "\sigma"] & \xi \arrow[l, shift left, "\rho"] \arrow[loop right, "p"] \\
 & \eta \arrow[ul, "u_0"] \arrow[u, "u"]
\end{tikzcd}
\qquad
\begin{array}{lcl}
\left.\begin{array}{rcl} \rho\sigma & = & \mathrm{id}_{\xi_0} \\ \sigma\rho & = & p \end{array}\right\} & \Longrightarrow & \left\{\begin{array}{rcl} p^2 & = & p \\ p\sigma & = & \sigma \\ \rho p & = & \rho \end{array}\right. \\[4ex]
\begin{array}{rcl} u & = & \sigma u_0 \end{array} & \Longrightarrow & \left\{\begin{array}{rcl} pu & = & u \\ \rho u & = & u_0 \end{array}\right.
\end{array}
\right.
\]
\textit{On se limite aux foncteurs $\mathcal{Q} \longrightarrow \mathcal{C}$, $\mathcal{Q}' \longrightarrow \mathcal{C}$ tels que $p$ soit transformé en une flèche} non identique \textit{de $\mathcal{C}$ (\footnote[195]{Je ne suis pas sûr si c'est raisonnable -- à vérifier avec soin !}) (donc la condition relative à $\mathcal{Q}$, $\mathcal{Q}'$ devient vide, si $\mathcal{C}$ est sans projecteurs).} $\Big)$

\textit{Sous ces conditions, $\mathcal{E}$ est accessible, et $f : \mathcal{E} \longrightarrow \mathcal{C}$ est accessible.}

% ---- from der-ch18_p186-190.tex
\setcounter{footnote}{195}

\printedpage{184}
% head: A. Grothendieck : Dérivateurs, Ch. XVIII, Édition des Universités de Montpellier II et Paris VII
% foot: Version 17 décembre 2025

\msp{page 475}
On trouve alors le

\dgrblock{dgrb1140}{What the two split-idempotent tests supply}
\begin{dcomment}{development}{p474}{What the two split-idempotent tests supply}
Grothendieck introduces $\mathcal Q$ and $\mathcal Q'$ to isolate the obstruction left after large filtered diagrams have been made constant in the base. One test attaches a target to a split idempotent; the other attaches a source. This is exactly the distinction between constructing a colimit and proving that a chosen object is presentable.

\textbf{The finite-test assertion of Proposition 4.16.6.1.}\ours\
Let $f:\mathcal E\to\mathcal C$ satisfy the hypotheses stated there, with $\mathcal C$ small and Karoubian. Then $\mathcal E$ and $f$ are accessible. In this assertion about the total category itself, it is enough to use the $\mathcal Q$ and $\mathcal Q'$ tests whose distinguished idempotent has nonidentity image in $\mathcal C$.

Here are the rank choices needed before the construction. The hypothesis H supplies accessible fibres, accessible pullbacks over arrows, and preservation of sufficiently filtered fibre colimits in those pullbacks. For a test category $B$ equal to $\Delta_1$, $\mathcal Q$, or $\mathcal Q'$, its fibre inclusions into $\mathcal E_B$ are accessible: maps from a fibre colimit to any other fibre are computed by the contravariant mixed hom functors of H2. This includes arrows of $B$ whose images are identities, where the assertion is ordinary representability inside one fibre. There are only a small number of the indicated tests. Choose a common sufficiently large strong accessibility rank $\kappa$ for their categories $\mathcal E_B$, their projections $\mathcal E_B\to B$, all fibres, and all fibre inclusions. Enlarge it to a suitable successor rank beyond the strict cardinal bound of Corollary~4.15.11. The uniformization used here is the same as in the preceding accessibility arguments; no claim at a prescribed smaller rank is needed.

Let $D:J\to\mathcal E$ be $\kappa$-filtered. The normal form proved near pp.~443--463 permits a cofinal, still $\kappa$-directed reindexing for which all base objects are one object $s$ and every strict base transition is one idempotent $p:s\to s$. If the new index has a last object, its colimit is already a stage and there is nothing to construct. Otherwise every index has a strict successor. If $p=1_s$, H2 shows directly that the colimit in $\mathcal E_s$ is a colimit in $\mathcal E$. Suppose then that $p\ne1_s$, and split it in $\mathcal C$:
\[
 s\xrightarrow{r}t\xrightarrow{i}s,
 \qquad ri=1_t,\quad ir=p.
\]

Write $K$ for the full split-idempotent subcategory of $\mathcal Q$ on $\xi,\xi_0$. The diagram lifts to $\mathcal E_K$ with all index objects labelled $\xi$. Extend $K\to\mathcal C$ to $\mathcal Q\to\mathcal C$ by sending $\eta$ to $t$ and $u_0$ to $1_t$. In $\mathcal Q$ the colimit of this base diagram is $\xi_0$, with structure arrows $\rho$. Indeed, the strict successors force any cocone arrow $a$ to satisfy $ap=a$, and these arrows factor uniquely through $\rho$. Accessibility of $\mathcal E_{\mathcal Q}$ supplies its colimit, and preservation by the projection puts its vertex over $\xi_0$: this finite category has no nontrivial isomorphisms. We obtain $z\in\mathcal E_t$ with cocone maps over $r$. Fullness of $K$ restricts the universal property to $\mathcal E_K$, so this is also a colimit there.

To check all targets in $\mathcal E$, take $y\in\mathcal E_c$. The base arrows of any cocone $D\to y$ are constant, say $a:s\to c$, and satisfy $ap=a$. Hence $a=\mu r$ for the unique $\mu=ai:t\to c$. Use now the $\mathcal Q$ test with $\eta\mapsto c$ and $u_0\mapsto\mu$. Its independently existing colimit again projects to $\xi_0$ and restricts to the same colimit in $\mathcal E_K$. The canonical isomorphism to $z$ is vertical because $\End_K(\xi_0)=\{1\}$. Its universal property gives exactly one arrow $z\to y$ over $\mu$ producing the original cocone. This proves the total universal property and, at the same time, preservation of this colimit by $f$.

We next prove presentability, since existence of colimits alone does not prove accessibility. Let $x$ be $\kappa$-presentable in a fibre $\mathcal E_c$, and let $x\to z$ lie over $\nu:c\to t$. In the $\mathcal Q'$ test put $\eta\mapsto c$, $\xi_0\mapsto t$, $\xi\mapsto s$, $u_0\mapsto\nu$, and $u\mapsto i\nu$. The object labelled $(x,\eta)$ is $\kappa$-presentable in $\mathcal E_{\mathcal Q'}$ by the strong rank choice. Accessibility gives a colimit of the lifted diagram there, and its projection is again $\xi_0$. Restricting to the same full $K$ identifies this colimit vertically with $(z,\xi_0)$. Consequently $x\to z$ factors through some $D_j$, giving surjectivity of the hom-colimit comparison in $\mathcal E$.

For injectivity, consider two maps from $x$ to stages with the same composite to $z$. Their base maps $a,b:c\to s$ have $ra=rb=\nu$. After advancing each stage once, their bases become $pa=pb=i\nu$. They are now maps in the same $\mathcal Q'$ test. Presentability there makes them equal at a common later stage. When $p=1_s$, precisely this argument uses the corresponding $\Delta_1$ test instead. The case with a final index is immediate. Thus every fibrewise $\kappa$-presentable object is $\kappa$-presentable in $\mathcal E$.

The union over the small base of small skeletons of these fibrewise objects is small. Every object of $\mathcal E$ is a $\kappa$-filtered colimit of them inside its own fibre, hence also in $\mathcal E$ by H2. This proves accessibility and completes the finite-test assertion. We did not locate this exact $\mathcal Q,\mathcal Q'$ formulation in the checked literature; the proof supplies the missing argument rather than attributing it to a later author. The restriction on $p$ just justified is specific to this theorem. It cannot be carried unchanged through arbitrary strict base changes, as the next example shows.
\end{dcomment}


\begingroup\itshape
\noindent\textbf{\upshape Corollaire 4.16.6.2.} Soit $f : \cE \longrightarrow \cC$ un foncteur. Conditions équivalentes :
\begin{enumerate}[label={}]
\item[a)] Pour toute catégorie $\cC'$ \textup{petite} et \textup{karoubienne}, et tout foncteur $\cC' \longrightarrow \cC$, $\cE' = \cE \times_{\cC} \cC'$ est accessible, et le foncteur $\cE' \longrightarrow \cC'$ est accessible.
\item[b)] Comme a), mais $\cC'$ étant une des trois catégories $\Delta_1$, $\cQ$, $\cQ'$ (où $\cQ$, $\cQ'$ sont définies dans 4.15.6.1 ci-dessus).
\item[b$'$)] \textup{(}\footnote{?, à vérifier avec soin !}\textup{)}. (Si $\cC$ est sans projecteurs.) Comme b), mais en se limitant au cas $\cC' = \boldsymbol{\Delta}_1$ (à l'exclusion de $\cQ$, $\cQ'$).
\end{enumerate}
\endgroup

\begingroup\itshape
\noindent\textbf{\upshape Corollaire 4.16.6.3.} Supposons $f : \cE \longrightarrow \cC$ fibrant ou cofibrant. Alors les conditions équivalentes précédentes impliquent que :
\begin{enumerate}[label={}]
\item[1$^\circ$)] Les $\cE_s$ sont accessibles ($s \in \Ob \cC$).
\item[2$^\circ$)] Pour toute flèche $\lambda : s \longrightarrow t$ dans $\cC$, le foncteur $\lambda^* : \cE_t \longrightarrow \cE_s$ (resp. $\cE_s \longrightarrow \cE_t$) est accessible.
\end{enumerate}
\noindent D'autre part, l'ensemble de ces conditions 1$^\circ$, 2$^\circ$) équivaut à la validité de b) ci-dessus, dans le cas $\cC' = \boldsymbol{\Delta}_1$, ou encore à la validité
\msp{page 476}
de a) ci-dessus, quand $\cC'$ est \textup{sans projecteurs}.
\endgroup

\dgrblock{dgrb1141}{A strict pullback can expose a hidden projector}
\begin{dcomment}{development}{p475}{A strict pullback can expose a hidden projector}
The uncertain clause b$'$ in Corollary~4.16.6.2 asks to discard both split-idempotent tests when the base has no nonidentity projectors. A strict pullback remembers more than the images of its arrows in that base. A projector upstairs may map to an identity downstairs, and this is exactly the case which the shortened criterion misses.

\textbf{Counterexample to b$'\Rightarrow$a.}\ours\
Let $\mathcal C$ be the walking isomorphism, with objects $0,1$ and exactly one arrow between each ordered pair. It has no nonidentity idempotents. Let $f:e\to\mathcal C$ select $0$. For any functor $\Delta_1\to\mathcal C$, the strict pullback is empty, a point, or $\Delta_1$, according to which endpoints map to $0$. Each is accessible and its projection to $\Delta_1$ preserves sufficiently filtered colimits. Thus the printed condition b$'$ holds.

Let $R$ be the walking split idempotent, with two objects $a,b$ and arrows
\[
 a\xrightarrow{r}b\xrightarrow{i}a,
 \qquad ri=1_b,\quad ir=p\ne1_a.
\]
Its only nonidentity endomorphism is $p$, and it is split, so $R$ is a small Karoubian category. Send $a$ to $0$, $b$ to $1$, and $r,i$ to the inverse isomorphisms of $\mathcal C$. Then $p$ is sent to $1_0$. The strict pullback $e\times_{\mathcal C}R$ has one object with endomorphism monoid $\{1,p\}$, where $p^2=p\ne1$.

This one-object category is not accessible. For any infinite regular $\kappa$, index a diagram by the ordinal $\kappa$, put the sole object at every stage, and put $p$ on every strict transition. A cocone to the sole object must have every component equal to $p$. If it were a colimit cocone, the two endomorphisms $1$ and $p$ of its vertex would induce distinct cocones. They induce the same one, since $1p=pp=p$. Thus the required $\kappa$-filtered colimit does not exist, for any $\kappa$.

Condition a) therefore fails. The unrestricted equivalence a)$\Leftrightarrow$b) is repaired by retaining all $\mathcal Q$ and $\mathcal Q'$ tests, including those which send $p$ to an identity. Indeed, applying Proposition~4.16.6.1 to a pullback over a small Karoubian $\mathcal C'$ may involve a nonidentity projector of $\mathcal C'$ whose image in $\mathcal C$ is an identity. The example was checked directly against the generators and compositions of these finite categories. We found no statement of this particular counterexample in the consulted sources; it is a correction of the queried implication, without a claim of priority.
\end{dcomment}


\textbf{Question 4.16.6.4.} \textit{Les conditions 1$^\circ$ et 2$^\circ$ précédentes impliquent-elles celles de 4.16.6.2 ? En d'autres termes, si la catégorie base $\cC$ est égale à $\cQ$ ou à $\cQ'$, alors ces conditions impliquent-elles que $\cE$ est accessible, et $f : \cE \longrightarrow \cC$ accessible ?}

Ça semble idiot de ne pas connaître la réponse !

\textbf{4.16.6.5.} J'ai envie d'appeler catégorie \textit{accessible sur $\cC$}, une catégorie $\cE$ sur $\cC$ satisfaisant les conditions équivalentes de 4.16.6.1 -- du moins dans le cas où on suppose $\cC$ petite. Pour que cette terminologie soit vraiment satisfaisant, il faudrait que l'on sache prouver que ces conditions impliquent ceci : pour toute catégorie accessible $\cM$ et tout foncteur accessible $\cM \longrightarrow \cC$, la catégorie $\cE \times_{\cC} \cM$ est accessible, et le foncteur $\cE \times_{\cC} \cM \longrightarrow \cM$ est accessible. % [sic] satisfaisant -> satisfaisante, accord avec « terminologie »

\msp{page 476bis}
L'idée, c'est qu'on appellerait (sous hypothèse de petitesse) \textit{catégorie accessible $\cE$} sur une autre $\cC$, une catégorie ayant la propriété qu'on vient de dire -- mais en supposant maintenant de plus que $\cC$ satisfasse à $\mathrm{L}_\pi$, j'imagine, pour tout foncteur accessible $\cM \longrightarrow \cC$, avec $\cM$ accessible, $\cE \times_{\cC} \cM$ est accessible et le foncteur canonique $\cE \times_{\cC} \cM \longrightarrow \cE$ est accessible.

\dgrblock{dgrb1142}{Relative accessibility and lifting isomorphisms}
\begin{dcomment}{development}{p476bis:L'idée, c'est qu'on appellerait}{Relative accessibility and lifting isomorphisms}
Grothendieck wants relative accessibility to survive replacement of the base by an accessible category. The corrected finite tests do give this property. Their strict-pullback formulation also forces an additional feature which explains why a bare accessible functor is insufficient.

\textbf{A criterion over a small Karoubian base.}\ours\
For $f:\mathcal E\to\mathcal C$ with $\mathcal C$ small and Karoubian, the following are equivalent: every strict pullback along a small Karoubian $\mathcal C'\to\mathcal C$ is accessible and has accessible projection to $\mathcal C'$; and $\mathcal E$ and $f$ are accessible and $f$ is an isofibration. Here isofibration means that for every $x\in\mathcal E$ and every isomorphism $u:f(x)\to t$ in $\mathcal C$, there is an isomorphism $\widetilde u:x\to y$ with $f(\widetilde u)=u$ and $f(y)=t$.

For necessity, use $\mathcal C'=\mathcal C$ to obtain accessibility of $\mathcal E$ and $f$. For the lifting assertion, use the walking split idempotent $R$ from the preceding block, mapping its two objects to $f(x)$ and $t$, its arrows $r,i$ to $u,u^{-1}$, and its projector $p$ to $1_{f(x)}$. In the accessible strict pullback there is an idempotent $(1_x,p)$ on $(x,a)$. An accessible category splits every idempotent: a sufficiently long regular-ordinal diagram with that idempotent on all strict transitions has a colimit, and its universal property gives the splitting. The only object of $R$ at which $p$ can split is $b$, since the monoid $\{1,p\}$ at $a$ does not split it internally. The splitting object is therefore $(y,b)$ for some $f(y)=t$. Its two splitting arrows project in $\mathcal E$ to inverse isomorphisms, because the first component of the original idempotent was $1_x$. The arrow over $r$ is the required lift of $u$.

For sufficiency, a small Karoubian category is accessible at a sufficiently large rank, and any functor between two such categories preserves sufficiently filtered colimits: those colimits reduce to splittings of idempotents, which every functor preserves. The iso-comma category of the accessible functors $\mathcal E\to\mathcal C\leftarrow\mathcal C'$ is accessible and its projections are accessible, by the pseudopullback theorem. Because $f$ is an isofibration, its strict pullback is equivalent to this iso-comma category. Explicitly, an object $(x,c',\alpha:f(x)\cong g(c'))$ can be replaced, using a lift of $\alpha$, by one with equal base objects and identity comparison. The inclusion of the strict pullback is fully faithful and this replacement proves essential surjectivity. This proves sufficiency.

The same argument now allows any accessible $\mathcal M\to\mathcal C$ in place of $\mathcal C'$. Its strict pullback is accessible, and both projections are accessible. This answers the intended base-change question of p.~476, and also reconciles the two different projections mentioned on p.~476 bis. The assertion here is over the stated small Karoubian base; the proposal for a general nonsmall base is a further definition.

The standard closure theorem used in this proof is discussed precisely in \citet[Examples~4.8 and Remark~4.9]{BourkeLack2024}. The identification of Grothendieck's strict test property with isomorphism lifting is our deduction; we did not find this formulation in those references. For actual Grothendieck fibrations and opfibrations, isomorphisms already lift, and the total-category proof supplied near pp.~438--440 gives the remaining accessibility. Thus Question~4.16.6.4 has an affirmative answer in both orientations. The distinction between strict pullbacks and pseudopullbacks accounts for the failure of b$'$ without making that structured result problematic.
\end{dcomment}


\msp{page 477}
\subsection*{4.16.7. Retour sur les sous-catégories fermées d'une catégorie accessible (4.16.1).}

\textbf{Proposition 4.16.7.1.} \textit{Soient $\cM$ une catégorie accessible, $(\uS_i)_{i\in I}$ une famille de sous-catégories fermées accessibles dans $\cM$. On ne suppose que \textup{card}\,$I$ soit petit. Soit $\uS$ la sous-catégorie strictement pleine de $\cM$ telle que $\Ob \uS = \bigcup_{i\in I} \Ob \uS_i$. Alors pour que $\uS$ % [sic] On ne suppose que -> On ne suppose pas que, « pas » apparemment omis
\printedpage{185}
% head: A. Grothendieck : Dérivateurs, Ch. XVIII, Édition des Universités de Montpellier II et Paris VII
% foot: Version 17 décembre 2025
soit accessible dans $\cM$, il faut et il suffit qu'il existe une \textup{petite} partie $J \subset I$, telle que $\Ob \uS = \bigcup_{i\in J} \uS_i$. (En particulier, si $I$ est ordonnée, et $i \longmapsto \uS_i$ croissante et grande devant tout petit cardinal, il faut et il suffit qu'il existe un $i \in I$ tel que $\uS = \uS_i$, i.e. tel que $j \geqslant i$ implique $\uS_j = \uS_i$.)}

C'est évidemment suffisant, en vertu de 4.13\dots{} Prouvons inversement que si $\uS$ est accessible dans $\cM$, alors la condition précédente est satisfaite. Soit $\pi$ un cardinal $\geqslant \pi_0$ tel que $\uS$ soit $\pi$-accessible dans $\cM$, donc
\[
\uS \ \xleftarrow{\ \sim\ }\ \Ind_\pi \uS(\pi) \ , \qquad\qquad \text{où } \uS(\pi) = \uS \cap \cM_\pi.
\]

\msp{page 478}
Soit $\uS(\pi)_0$ une sous-catégorie réduite de $\uS(\pi)$, telle que $\uS(\pi)_0 \hookrightarrow \uS(\pi)$ soit une équivalence, de sorte que $\uS(\pi)_0$ est petite. Pour tout $x \in \Ob \uS(\pi)_0$, soit $i_x \in I$ tel que $x \in \Ob \uS_{i_x}$. Soit $J$ l'ensemble de tous les $i_x$, que est une petite partie de $I$, et soit $\uS_J$ la sous-catégorie strictement pleine de $\cM$ ayant pour ensemble d'objets $\bigcup_{j\in J} \Ob \uS_j$. Donc $\uS(\pi)_0$ est contenue dans $\uS_J$, donc aussi $\uS(\pi)$, puisque $\uS_J$ est strictement pleine : % [sic] que est -> qui est, pronom relatif sujet
\[
\uS(\pi) \ \hookrightarrow\ \uS_J \ .
\]
D'ailleurs $\uS_J$ est accessible dans $\cM$ (loc. cit.), et est fermée dans $\cM$ comme réunion de sous-catégories fermées, \textit{a fortiori elle satisfait à $L_\pi$}. Donc, comme $\uS(\pi) \subset \uS_J$, le foncteur essentiellement surjectif $\Ind_\pi(\uS(\pi)) \longrightarrow \uS$ se factorise par $\uS_J$. Donc $\uS_J \hookrightarrow \uS$ est essentiellement surjectif, et comme $\uS_J$ est strictement pleine, on en conclut $\uS_J = \uS$, q.e.d.

\msp{page 479}
\textbf{Proposition 4.16.7.2.} \textit{Soient $\cM$ une catégorie accessible, $x$ un objet de $\cM$, et ${}_x\cM$ la sous-catégorie fermée de $\cM$ engendrée par $x$, formée des $y \in \Ob \cM$ tels que $\Hom(x,y) \neq \emptyset$, i.e. l'image essentielle de $x\backslash\cM \longrightarrow \cM$. Alors ${}_x\cM$ est une sous-catégorie de $\cM$ accessible dans $\cM$.}

En effet, c'est l'image inverse de la sous-catégorie $\mathrm{Ens}^*$ accessible dans $\mathrm{Ens}$ par le foncteur représentable (donc accessible)
\[
h_x \ :\ y \longmapsto \Hom(x,y) \ :\ \cM \ \longrightarrow\ \mathrm{Ens} \ ,
\]
donc est accessible dans $\cM$ en vertu de 4.13\dots

\textbf{Corollaire 4.16.7.3.} \textit{Soit $\cM$ une catégorie accessible, $\uS$ une sous-catégorie fermée. Pour que $\uS$ soit accessible, il faut et il suffit qu'il existe une \textup{petite} partie $J$ de $\uS$, telle que $\Ob \uS = \bigcup_{x\in J} \Ob\,(\,{}_x\cM)$.}

En effet, on a
\[
\Ob \uS \ =\ \bigcup_{x\in\Ob\uS} \Ob\,(\,{}_x\cM) \ ,
\]
i.e. $\uS$ apparaît comme réunion de
\msp{page 480}
la famille (pas nécessairement petite) de sous-catégories fermées ${}_x\cM$ de $\cM$ (indexées par $x \in \Ob \uS$), accessibles dans $\cM$ en vertu de 4.16.7.2. Donc notre assertion est un cas particulier de 4.16.7.1.

\dgrblock{dgrb1143}{Closed subcategories and small weakly initial families}
\begin{dcomment}{exposition}{p479}{Closed subcategories and small weakly initial families}
Grothendieck's return to closed subcategories isolates the generation condition which the earlier examples lacked. A cosieve $\mathcal S$ is closed under any existing nonempty colimit of its objects: choosing one stage already gives an arrow from $\mathcal S$ into the colimit. Its accessibility therefore depends on how many different starting objects are required to reach every object of $\mathcal S$.

A family $A\subseteq\operatorname{Ob}\mathcal S$ generates the cosieve precisely when every $s\in\mathcal S$ receives a map from some $a\in A$. This is a weakly initial family; no uniqueness of the map is required. If it is small, choose a common useful rank at which all its objects are presentable. Let $\mathcal T$ be the complementary sieve. A sufficiently filtered diagram in $\mathcal T$ cannot have its colimit in $\mathcal S$: a map from one of the presentable generators to that colimit would factor through a stage, placing that stage in the cosieve. Thus $\mathcal T$ is closed under those colimits, and the canonical presentations used below prove accessibility of both parts.

Conversely, if $\mathcal S$ is accessible in $\mathcal M$, take a small skeleton of its presentable generators at a common useful rank. Every object of $\mathcal S$ is a nonempty filtered colimit of these generators, so receives a map from one of them. This proves Corollary~4.16.7.3 directly and explains why its small family is the right condition.

The union criterion follows with the same economy. If a union of accessible cosieves is accessible, choose one member of the family containing each object of a small generating skeleton of the union. These choices use only a small subfamily. Its union is itself a cosieve, hence contains every colimit generated by that skeleton and is the whole union. Conversely, the union of the small generating families of a small subfamily is again small. In particular, a union indexed by an order directed beyond every small cardinal is accessible exactly when a single stage already contains it. No finiteness or countability assertion is hidden in the word small: it is always relative to the fixed universe.
\end{dcomment}


\printedpage{186}
% head: A. Grothendieck : Dérivateurs, Ch. XVIII, Édition des Universités de Montpellier II et Paris VII
% foot: Version 17 décembre 2025

On peut aussi, pour toute partie petite $J$ de $\Ob \uS$, considérer la sous-catégorie fermée $\uS_J \subset \uS$ de $\cM$ engendrée par $J$. L'ensemble $I$ des parties $J$ en question, ordonné par inclusion, est grand devant tout cardinal $\pi \in \U$, et $\uS$ est réunion croissante des $\uS_J$. D'autre part, la proposition nous dit que $\uS$ est accessible dans $\cM$ si et seulement s'il existe un $J \in I$, tel que $\uS_J = \uS$, i.e. tel que $\uS_J = \uS_{J'}$ pour $J \subset J'$.

\begingroup\itshape
\noindent\textbf{\upshape Proposition 4.16.7.4.} Soit $\cM$ une catégorie accessible. Les conditions suivantes sont équivalentes :
\begin{enumerate}[label={}]
\item[a)] Il existe une sous-catégorie $\uS$ de $\cM$ fermée dans $\cM$, et non accessible dans $\cM$.
\item[b)] Soit $I$ l'ensemble des ordinaux $\alpha$ tels que $\card \alpha \in \U$, i.e. des
\msp{page 481}
ordinaux $\alpha$ qui sont dans $\U$. Donc $I = I_{\alpha_0}$, où $\alpha_0$ est le plus petit ordinal qui n'appartienne pas à $\U$, et $I$ est grand devant tout petit cardinal $\pi$, i.e. devant tout cardinal dans $\U$ \textup{(}\footnote{introduire $I = I_{\U}$ avant l'énoncé de la proposition.}\textup{)}. La condition envisagée est la suivante :

Il existe une suite transfini % [sic] transfini -> transfinie, accord avec « suite »
\[
\underline{x} \ =\ (x_\alpha)_{\alpha\in I}
\]
d'éléments de $\Ob \cM$, telle que l'on ait
\[
\Hom(x_\beta, x_\alpha) \ =\ \emptyset \qquad\qquad \textit{si } \beta < \alpha.
\]
De façon précise, pour une telle suite transfinie, soit $\uS$ ou ${}_{\underline{x}}\cM$ la sous-catégorie fermée de $\cM$ engendrée par la famille des $x_\alpha$, i.e. telle que $\Ob \uS = \bigcup_{\alpha\in I} \Ob\ {}_{x_\alpha}\cM = \bigcup_{\alpha\in I} x_\alpha\backslash\Ob \cM$. Alors ${}_{\underline{x}}\cM$ est une sous-catégorie fermée non accessible dans $\cM$, et toute sous-catégorie fermée non accessible dans $\cM$ s'obtient de cette façon.
\end{enumerate}
\endgroup

\textsc{Démonstration}. Considérons une suite $\underline{x} = (x_\alpha)_{\alpha\in I}$ comme dessus. Si $J$ est une petite partie de $I$, il existe un $\alpha \in I$ qui majore les $\beta \in J$, donc $\Hom(x_\beta, x_\alpha) = \emptyset$, pour $\beta \in J$,
\msp{page 482}
ce qui signifie que $x_\alpha$ n'appartient à aucune des $\uS_\beta = {}_{x_\beta}\cM$, pour $\beta \in J$. Donc $\uS$ n'est par réunion d'une petite sous-famille de la famille des ${}_{x_\alpha}\cM$, donc en vertu de 4.16.7.1, $\uS$ n'est par accessible dans $\cM$. % [sic] n'est par réunion -> n'est pas réunion ; n'est par accessible -> n'est pas accessible, « par » pour « pas » (2 fois)

Inversement, soit $\uS$ une telle sous-catégorie dans $\cM$. Considérons un bon ordre sur $\Ob \uS$ (objets $s_\alpha$), et construisons la suite des $x_\alpha$ ($\geqslant s_\alpha$) par récurrence transfinie (\footnote{\textbf{N.B.} On a $\card \uS \notin \U$, sinon, $\uS$ serait petite, donc accessible dans $\cM$, contrairement à l'hypothèse.}).

1) On choisit $x_0 =$ premier élément de $\Ob \uS$ (existe, car $\Ob \uS \neq \emptyset$, i.e. $\uS \neq \emptyset$, car sinon, $\uS$ serait accessible dans $\cM$).

2) Si les $x_\beta$ ($\beta < \alpha$) sont construits, on prend pour $x_\alpha$ le premier élément de $\uS$ majorant les $x_\beta$ et qui n'appartient pas à la sous-catégorie fermée de $\cM$ engendrée par les $x_\beta$ ($\beta < \alpha$). (Cette sous-catégorie fermée est $\subset \uS$, et $\neq \uS$, car $\uS$ n'est pas accessible dans $\cM$.) On a donc bien
\[
\Hom(x_\beta, x_\alpha) \ =\ \emptyset \qquad\qquad \text{si } \beta < \alpha.
\]

\printedpage{187}
% head: A. Grothendieck : Dérivateurs, Ch. XVIII, Édition des Universités de Montpellier II et Paris VII
% foot: Version 17 décembre 2025

\msp{page 483}
Ainsi, la récurrence transfinie se poursuit sur $I$ tout entier. On voit par récurrence transfinie qu'on a, pour tout $\alpha \in I$,
\[
s_\alpha \ \leqslant\ x_\alpha
\]
(card $x_\alpha > x_\beta \geqslant s_\beta$ pour $\beta < \alpha$, donc $x_\alpha > s_\beta$ pour tout $\beta < \alpha$, ce qui équivaut à $x_\alpha \geqslant s_\alpha$). De plus, on trouve aussi, par récurrence transfinie, que
\[
\beta \ <\ \alpha \qquad\quad \Longrightarrow \qquad\quad
\text{\parbox[c]{5.5cm}{$s_\beta$ appartient à la sous-catégorie fermée $\uS_\alpha$ de $\cM$ engendrée par les $x_{\beta'}$ avec $\beta' < \alpha$.}}
\quad ,
\tag*{$(*)$}
\]
puisque par construction $s_\alpha$ est le \textit{premier} des éléments de $\uS$ qui n'est \textit{pas} dans ladite sous-catégorie $\uS_\alpha$ de $\cM$. Cela montre que tout $s_\beta$ appartient à la sous-catégorie fermée ${}_{\underline{x}}\cM$ de $\cM$ engendrée par la famille de \textit{tous} les $x_\alpha$, donc on trouve que $\uS \subset {}_{\underline{x}}\cM \subset \uS$, donc $\uS = {}_{\underline{x}}\cM$, q.e.d.

\begingroup\itshape
\noindent\textbf{\upshape Proposition 4.16.7.5.} Pour un univers $\U$ donné (contenant $\mathcal{N}$ comme élément), les conditions suivantes sont équivalentes :
\begin{enumerate}[label={}]
\item[1)] Il existe une catégorie $\U$-accessible $\cM$ ayant une sous-catégorie fermée non accessible % [sic] 1) -> a), la démonstration renvoie à a) et b)
\msp{page 484}
dans $\cM$.
\item[b)] Il existe une catégorie $\cC \in \U$, et une suite transfinie d'éléments de
\[
\cC^\wedge \ \overset{\text{\upshape déf}}{=}\ \underline{\mathrm{Hom}}(\cC, \mathrm{Ens}_{\U}) \ ,
\]
$\underline{x} = (x_\alpha)_{\alpha\in I}$, indexée par l'ensemble de tous les ordinaux $\alpha \in \U$, et telle que
\[
\Hom(x_\alpha, x_\beta) \ =\ \emptyset \qquad\qquad \textit{si } \alpha < \beta.
\]
\end{enumerate}
\endgroup

\dgrblock{dgrb1144}{The transfinite construction and its size bound}
\begin{dcomment}{exposition}{p483}{The transfinite construction and its size bound}
Grothendieck turns failure of a small weakly initial family into a sequence with no arrows pointing forward. The order used in the recursive construction is an external well order of representatives. It is not the preorder defined by existence of a morphism; confusing them would make the displayed condition contradict the instruction to choose a larger object.

Write $\mathfrak U=V_\theta$ for the fixed infinite Grothendieck universe. An accessible $\mathfrak U$-category has at most $\theta$ isomorphism classes. Indeed, a small dense subcategory embeds it fully faithfully in a presheaf category on a $\mathfrak U$-small category. For every cardinal bound below $\theta$, only a $\mathfrak U$-small set of isomorphism classes of such presheaves has values within that bound; taking all bounds gives at most $\theta$ classes. We may therefore work with a skeleton of a cosieve $\mathcal S$, rather than a possibly much larger set of duplicate objects. If $\mathcal S$ is nonaccessible, this skeleton cannot have cardinal less than $\theta$, by the small-family criterion. Enumerate it as $(s_\alpha)_{\alpha<\theta}$.

At stage $\alpha$, let $\mathcal S_\alpha$ be the cosieve generated by the already chosen $x_\beta$, $\beta<\alpha$. This is generated by a small set, hence is not all of $\mathcal S$. Choose $x_\alpha$ to be the first element of the well order outside $\mathcal S_\alpha$. Then
\[
 \Hom(x_\beta,x_\alpha)=\varnothing\qquad(\beta<\alpha).
\]
The indices of the chosen objects are strictly increasing, since every object preceding the least previously uncovered object was already covered. By induction, before stage $\alpha$ every $s_\beta$ with $\beta<\alpha$ is covered: if one were still uncovered, the next least-uncovered index would be less than $\alpha$, contrary to the strictly increasing sequence of prior choices. Thus the resulting family generates all of $\mathcal S$. In the printed explanation of this step, the least uncovered object is $x_\alpha$, not $s_\alpha$.

Conversely, let $(x_\alpha)_{\alpha<\theta}$ have no forward morphisms and let $\mathcal S$ be its generated cosieve. If it had a small weakly initial family $(a_j)$, each $a_j$ would receive a map from some $x_{\beta_j}$. The small set of ordinals $\beta_j$ has a bound below $\theta$. A later $x_\alpha$ would have to receive a map from one $a_j$, giving the forbidden composite $x_{\beta_j}\to a_j\to x_\alpha$. This proves nonaccessibility and both directions of Proposition~4.16.7.4.

The same full presheaf embedding proves Proposition~4.16.7.5, since it preserves the emptiness of hom sets. It does not prove the additional assertion about a sequence consisting entirely of fields: the algebras constructed in the appendix need not be fields. The following development at p.~485 answers that stronger question by a separate graph-to-field construction and identifies its exact dependence on the universe boundary.
\end{dcomment}


\textsc{Démonstration}. On a b) $\Longrightarrow$ a) en prenant $\cM = \cC^\wedge$, en vertu de la proposition précédente. On voit de même a) $\Longrightarrow$ b), en plongeant $\cM$ dans une catégorie $\cC^\wedge$ ($\cC$ étant une petite sous-catégorie strictement génératrice de $\cM$).

Je n'ai pas été capable (\footnote{Mais si, finalement, cf. 4.16.7.10 plus bas.}) de construire une suite $(x_\alpha)_{\alpha\in I}$ comme exigée dans cette proposition, ni de prouver qu'une telle suite n'existe pas, même dans un cas particulier tel que $\cM = (k\text{-alg. comm.})$, $k$ un anneau commutatif. Je n'ai pas
\msp{page 485}
même prouvé qu'on ne puisse trouver une suite transfinie de \textit{corps} $k_\alpha$, ayant les propriétés dites.

A fortiori, je n'ai pu tirer au clair la question si une sous-catégorie $\underline{\mathrm{L}}$-stable d'une catégorie $\cM$ accessible est toujours accessible dans $\cM$ (\footnote{C'est faux, cf. plus bas.}). Il se pourrait qu'il n'y ait pas de tel cas. Il se pourrait aussi qu'il y en ait, même avec une sous-catégorie fermée (laquelle est stable par tout type de $\varinjlim$ représentable dans $\cM$, sauf la $\varinjlim$ sur une catégorie d'indices vide $\ldots$).

\printedpage{188}
% head: A. Grothendieck : Dérivateurs, Ch. XVIII, Édition des Universités de Montpellier II et Paris VII
% foot: Version 17 décembre 2025

Pour les sous-catégories \textit{ouvertes}, je rappelle que j'ai construit (dans l'appendice) un exemple d'une telle sous-catégorie, non accessible dans $\cM$. Au sujet des sous-catégories ouvertes, je signale ici l'énoncé (quasi tautologique) :

\msp{page 486}
\textbf{Proposition 4.16.7.6} (\footnote{Cette proposition et son corollaire sont contenus dans 4.16.7.8 ci-dessous, il n'y a pas lieu d'en faire un énoncé séparé.}). \textit{Soient $\cM$ une catégorie accessible, $\uS$ une sous-catégorie \textup{ouverte} de $\cM$. Pour que $\uS$ soit accessible dans $\cM$, il faut et il suffit qu'elle soit $\underline{L}$-stable dans $\cM$, i.e. qu'il existe un cardinal $\pi_0$ tel que $\cM$ satisfasse à $L_{\pi_0}$, et que $\uS$ soit stable dans $\cM$ par $\varinjlim$ de type $L_{\pi_0}$. De façon plus précise, si cette condition est satisfaite, alors pour tout cardinal $\pi$ tel que $\pi \geqslant \pi_0$ et que $\pi$ soit utile pour $\cM$, $\uS$ est $\pi$-accessible dans $\cM$. Ou encore : si $\pi$ est un cardinal utile pour $\cM$, alors $\uS$ est $\pi$-accessible dans $\cM$ si et seulement s'il est $L_\pi$-stable dans $\cM$.}

\dgrblock{dgrb1145}{The field question detects a Vopenka universe}
\begin{dcomment}{development}{p485}{The field question detects a Vopenka universe}
The question about fields has an answer depending on the chosen universe. The later example with arbitrary $k$-algebras does not by itself answer this narrower question. A field functor from graph theory provides the missing passage.

\textbf{Universe-relative field criterion.}\ours\
Let $\mathfrak U=V_\kappa$, with $\kappa$ strongly inaccessible and $\omega\in\mathfrak U$. There exists a sequence $(k_\alpha)_{\alpha<\kappa}$ of fields in $\mathfrak U$ such that
\[
\operatorname{Hom}(k_\alpha,k_\beta)=\varnothing\quad(\alpha<\beta)
\]
if and only if $\kappa$ is not a Vop\v{e}nka cardinal. If such a sequence exists, one can choose characteristic zero throughout and make the Hom sets empty for \emph{all} distinct indices.

Here a Vop\v{e}nka cardinal means that every $\kappa$-element family of structures in $V_\kappa$ in a fixed small language contains distinct members with an elementary embedding between them. All external subsets of $V_\kappa$ are allowed as families. This is stronger than merely asking $V_\kappa$ to satisfy an internally definable axiom scheme. A fixed language with one binary and one unary relation suffices \citep[Definitions~1--3]{BrookeTaylor2011Vopenka}.

First suppose $\kappa$ is Vop\v{e}nka. The proposed fields are pairwise nonisomorphic. For every infinite $\mu<\kappa$, at most $2^\mu<\kappa$ isomorphism types of fields have size at most $\mu$: put their operations on subsets of $\mu$ and count the finite-arity tables. Regularity of $\kappa$ therefore lets us thin the sequence to length $\kappa$ with strictly increasing field cardinalities and indices. At each stage, the previous suprema are still below $\kappa$, whereas fewer than $\kappa$ indices have bounded field size. Backwards homomorphisms in this subsequence are impossible because field homomorphisms are injective. Forward ones are excluded by hypothesis. The Vop\v{e}nka property nevertheless supplies an elementary embedding between two distinct fields, hence a field homomorphism, a contradiction.

Conversely, suppose $\kappa$ is not Vop\v{e}nka. Choose $\kappa$ structures $M_\alpha\in V_\kappa$ in the fixed finite relational language with no elementary embedding between distinct members. For every formula $\varphi(\bar x)$ add the relation $R_\varphi$ of its true tuples, including negated formulas and sentences. This countable relational expansion has the property that any injection preserving all its relations is elementary: $R_\varphi$ preserves truth and $R_{\neg\varphi}$ preserves falsity.

We give the graph encoding explicitly, to match the precise graph morphisms needed below. Make a core graph with an element vertex $v_a$ for every $a\in M$, a tuple vertex $t_{\varphi,\bar a}$ for each true $R_\varphi(\bar a)$, and a position vertex $s_{\varphi,\bar a,j}$ for each coordinate. The core edges are exactly
\[
t_{\varphi,\bar a}\;-\;s_{\varphi,\bar a,j}\;-\;v_{a_j}.
\]
This core is bipartite: elements and tuple vertices form one part, positions the other. Assign one type to elements, a distinct type to each formula's tuple vertices, and a distinct type to each formula-coordinate pair. There are countably many types; enumerate them by $n\in\omega$.

At every core vertex of type $n$, attach two cycles of lengths $4n+5$ and $4n+7$, identifying one vertex of each cycle with that core vertex and otherwise using fresh vertices. Every odd simple cycle in the resulting graph is one of these attached cycles. A simple cycle cannot enter and leave an attached block through its sole attachment vertex without repeating that vertex; a cycle confined to the core is even.

An induced-subgraph embedding must send each tag cycle to a tag cycle of the same length. The two cycles at a source core vertex meet exactly in that vertex; their images meet in its image. As the target tag cycles of these two lengths intersect only when they are attached at the same core vertex, that image is a core vertex of the same type. Thus every core type is preserved. The image of a tuple vertex has the same formula type, and its unique position neighbor of type $(\varphi,j)$ forces the corresponding coordinate element to be the image of $a_j$. The induced injection on elements preserves every $R_\varphi$, and is elementary. The graphs $G(M_\alpha)$ consequently have no induced-subgraph embeddings between distinct members.

Their sizes are at most $\max(|M_\alpha|,\aleph_0)$: finite tuples, countably many relation symbols, and the finite tags do not exceed that bound. Now apply \citet[Theorem~1.1, Section~5.A and Remark~5.1]{MillerPoonenSchoutensShlapentokh2018}. They construct a fully faithful functor from symmetric irreflexive graphs with induced-subgraph embeddings to fields with field homomorphisms. Its values have characteristic zero and size $\max(|G|,\aleph_0)$. One detail of their proof must be read with finite support in view: the real-specialization argument in Lemma~5.3(ii) is applied to the finite-vertex subfield containing the parameters of an alleged point and the indicated vertices $i,j$. An arbitrarily large whole field need not embed in $\mathbb R$, but that finite-parameter argument proves the stated exclusion at every cardinality. Therefore
\[
k_\alpha=F(G(M_\alpha))
\]
belongs to $\mathfrak U$ and has no homomorphism to $k_\beta$ for $\alpha\ne\beta$. Standard presentations of cardinality below $\kappa$ may be chosen inside $V_\kappa$. This completes the proof.

In particular, such a field sequence exists for the smallest universe containing $\omega$. Its first inaccessible height cannot be Vop\v{e}nka. Indeed the characterization by extendibility below $\kappa$ gives an elementary embedding $j:V_\eta\to V_\zeta$ of critical point $\delta<\kappa$, with $\eta>\delta+\omega$ \citep[Proposition~7]{BrookeTaylor2011Vopenka}. The derived ultrafilter
\[
\{X\subseteq\delta:\delta\in j(X)\}
\]
is nonprincipal and $\delta$-complete, so $\delta$ is measurable and hence inaccessible, contradicting first inaccessibility. At a Vop\v{e}nka universe the first half of the argument excludes the sequence.

This is a deduction from established theorems, with the reduction supplied in detail and no claim of priority. The field functor is a substantial algebraic construction: its cardinal preservation and full faithfulness, rather than merely the existence of fields with few automorphisms, are what the argument needs. The manuscript's question should accordingly be recorded as dependent on the universe, rather than left as an unqualified currently open field problem.
\end{dcomment}


Cette dernière assertion résulte aussitôt du corollaire suivant, lui-même tautologique.

\textbf{Corollaire 4.16.7.7.} \textit{Soit $\pi$ un cardinal comme indiqué ci-dessus, soit $\uS(\pi) = \uS \cap \cM_\pi$, et soit $x \in \Ob \cM$ et
\[
x \ \simeq\ \varinjlim_{i\in I} x_i \qquad\qquad
\begin{cases}
x_i \in \Ob \cM_\pi \\[1ex]
I \ \textit{grande devant}\ \pi
\end{cases}
\]
une représentation de $x$ comme image essentielle d'un objet de $\Ind_\pi \cM_\pi$. Alors $x \in \Ob \uS$ $\Longleftrightarrow$ $x_i \in \Ob \uS(\pi)$ $\forall\, i \in I$.}

\msp{page 487}
J'ai finalement tiré au clair la question de l'accessibilité des sous-catégories ouvertes et des sous-catégories fermées d'une catégorie accessible :

\begingroup\itshape
\noindent\textbf{\upshape Proposition 4.16.7.8.} Soient $\cM$ une catégorie accessible, $\uS$ une sous-catégorie fermée de $\cM$, $\uT$ la sous-catégorie ouverte complémentaire de $\uS$, $\pi$ un cardinal infini utile pour $\cM$. Les conditions suivantes sont équivalentes :
\begin{enumerate}[label={}]
\item[a)] $\uS$ est $\pi$-accessible dans $\cM$.
\item[a$'$)] $\uS$ est le fermé engendré par $\uS(\pi) \overset{\text{\upshape déf}}{=} \uS \cap \cM_\pi$, i.e. pour $x \in \Ob \cM$, on a $x \in \Ob \uS$ si et seulement si $\exists\, \xi \in \Ob \uS(\pi)$ tel que $\Hom_{\cM}(\xi, x) \neq \emptyset$.
\item[b)] $\uT$ est $\pi$-accessible dans $\cM$.
\item[b$'$)] $\uT$ est $\underline{L}_\pi$-stable dans $\cM$.
\end{enumerate}
\endgroup

\noindent Notons tout de suite le

\begingroup\itshape
\noindent\textbf{\upshape Corollaire 4.16.7.9.} Pour $\cM$, $\uS$, $\uT$ comme dessus, conditions équivalentes :
\begin{enumerate}[label={}]
\item[a)] $\uS$ est accessible dans $\cM$.
\item[a$'$)] Il existe une \textup{petite} partie $A$ de $\Ob \cM$, telle que $\uS$ soit la sous-catégorie fermée ${}_A\cM$ engendrée par $A$, i.e. $\uS = \bigcup_{\xi\in A} {}_\xi\cM$, où ${}_\xi\cM$ est la sous-catégorie fermée de $\cM$ formée des $x \in \Ob \cM$ tel que $\Hom(\xi, x) \neq \emptyset$.
\item[b)] \msp{page 488}$\uT$ est accessible dans $\cM$.
\end{enumerate}
\endgroup

\dgrblock{dgrb1146}{One canonical slice proves the open–closed criterion}
\begin{dcomment}{exposition}{p487}{One canonical slice proves the open--closed criterion}
Grothendieck's final criterion makes the algebraic obstruction usable. It says that the inaccessible no-point cosieve and the failure of filtered-colimit closure of the point-admitting sieve are two manifestations of one phenomenon. The proof takes place in a single category of presentable objects over a fixed target; several omitted slice denominators in the printed proof obscure this otherwise exact argument.

Fix a useful $\pi$, write $\kappa=\pi^+$, and let $\mathcal A$ be a small skeleton of $\mathcal M_\pi$. For $x\in\mathcal M$, the canonical category $\mathcal A/x$ is $\kappa$-filtered, and its tautological diagram has colimit $x$. If $\mathcal S$ is a replete cosieve and $\mathcal T$ its complementary sieve, put
\[
 J_x=(\mathcal S\cap\mathcal A)/x,
 \qquad K_x=(\mathcal T\cap\mathcal A)/x.
\]
These are a cosieve and its complementary sieve in $\mathcal A/x$, not in $\mathcal S(\pi)$. A nonempty full cosieve in a filtered category is cofinal. Given an object $j$ in the cosieve and any object $a$ of the ambient category, a common successor belongs to the cosieve; common successors and equalizations similarly verify connectedness of the relevant comma categories. For the stronger $\kappa$-filtered assertion, first cocone any fewer-than-$\kappa$ diagram in the ambient category and then move its vertex into the cosieve. In contrast, a proper complementary sieve cannot be cofinal, since an object in the cosieve has no arrow into it. Hence
\[
 J_x\ne\varnothing\ \Longleftrightarrow\ J_x\text{ cofinal},
 \qquad
 K_x\text{ cofinal}\ \Longleftrightarrow\ J_x=\varnothing.
\]

For completeness, the Ind-image criterion used here has a direct proof. If $x$ is a $\kappa$-filtered colimit of objects of a full $\mathcal B\subseteq\mathcal A$, every map $a\to x$ with $a\in\mathcal A$ factors through a stage by presentability. Equalities between two factorizations hold at a later stage. Therefore the presentation functor is cofinal in $\mathcal A/x$, and so is the full subcategory $\mathcal B/x$ through which it factors. Conversely, if $\mathcal B/x$ is cofinal, restriction of the canonical presentation expresses $x$ as its colimit. A full cofinal subcategory of this $\kappa$-filtered category is again $\kappa$-filtered: cocone a small diagram in the ambient category and use fullness after moving the vertex into the subcategory. This proves Lemma~4.16.7.11 with its slice variables restored.

Apply the criterion to $\mathcal B=\mathcal S\cap\mathcal A$. The condition that the Ind-image be exactly $\mathcal S$ becomes
\[
 x\in\mathcal S
 \quad\Longleftrightarrow\quad
 J_x\ne\varnothing,
\]
which says precisely that $\mathcal S$ is generated by its $\pi$-presentable objects. Taking complements turns this into
\[
 x\in\mathcal T
 \quad\Longleftrightarrow\quad
 K_x\text{ is cofinal in }\mathcal A/x,
\]
the identical Ind-image criterion for $\mathcal T$. Finally, if $\mathcal T$ is closed under the relevant filtered colimits, every presentation of an object in $\mathcal T$ has all its stages in $\mathcal T$, by the sieve condition, and conversely a presentation with all stages there has colimit there by closure. This proves its accessibility and gives the equivalence with b$'$.

Passing to a common larger useful rank yields Corollary~4.16.7.9. This proof explains both why the theorem is valid and why it settles the appendix's abandoned question about the complementary cosieve. The only changes required in the printed argument are the missing slice or presentability subscripts, the names of the Ind functor, and the final reference to Proposition~4.16.7.6.
\end{dcomment}


% ---- from der-ch18_p191-195.tex
\printedpage{189}
% head: A. Grothendieck : Dérivateurs, Ch. XVIII, Édition des Universités de Montpellier II et Paris VII

\begin{enumerate}
\item[\textit{b$'$)}] \textit{$\underline{T}$ est $\underline{L}$-stable dans $\cM$.}
\end{enumerate}

\noindent\textbf{Corollaire 4.16.7.10.} \textit{Soit $\cM$ une catégorie accessible. Les conditions suivantes sont équivalentes} (\footnote[202]{Exemple où ces conditions sont satisfaites : $\cM = (k\text{-alg. comm.})$, $k$ un corps commutatif infini, $\mathfrak{U}$ étant le plus petit univers tel que $\mathbf{N}\in\mathfrak{U}$.}).
\begin{enumerate}
\item[\textit{a)}] \textit{Il existe une sous-catégorie fermée de $\cM$, non accessible dans $\cM$.}
\item[\textit{b)}] \textit{Il existe une sous-catégorie ouverte dans $\cM$, non accessible dans $\cM$.}
\item[\textit{c)}] \textit{Il existe une suite transfinie}
\[
\underline{x} \ = \ (x_\alpha)_{\alpha\in I}
\]
\textit{d'objets de $\cM$, indexée par l'ensemble $I$ des ordinaux de l'univers de référence $\mathfrak{U}$, telle que l'on ait}
\[
\Hom(x_\alpha,x_\beta) \ = \ \emptyset \qquad \text{\textit{pour} } \alpha<\beta .
\]
\end{enumerate}

\noindent En effet, l'équivalence de a) et b) est claire par le corollaire précédent, et celle de a) et c) n'est autre que 4.16.7.4.

\noindent\textsc{Démonstration de} 4.16.7.8. Nous aurons à utiliser le

\noindent\textbf{Lemme 4.16.7.11.} \textit{Soient $\cM$ une catégorie accessible, $\pi$ un cardinal utile pour $\cM$, $\underline{S}$ une sous-catégorie pleine de $\cM$, $\underline{S}(\pi)\defeq\underline{S}\cap\cM_\pi$, d'où un foncteur pleinement fidèle}
\[
\varphi \ : \ \Ind_\pi(\underline{S}(\pi)) \longrightarrow \cM \ .
\]

\msp{page 489}
\begin{enumerate}
\item[\textit{a)}] \textit{Soit $x\in\Ob\cM$. Pour que $x$ appartienne à l'image essentielle du foncteur $\psi$, il faut et il suffit que la sous-catégorie pleine $\underline{S}(\pi)/x$ de la catégorie $\cM_\pi/x$ (qui est filtrante et même grande devant $\pi$) soit cofinale (donc aussi grande devant $\pi$).} % [sic] \psi -> \varphi, the functor of the lemma is called \varphi in the display above
\item[\textit{b)}] \textit{Pour que $\underline{S}$ soit $\pi$-accessible dans $\cM$, il faut et il suffit que pour tout $x\in\Ob\cM$, on ait}
\[
x \ \in \ \Ob\underline{S} \qquad\iff\qquad \underline{S}(\pi)/x \text{ \textit{est cofinale dans} } \cM_\pi/x
\]
\textit{(ou encore, que $\underline{S}$ satisfasse à $L_\pi$, et que pour tout $x$ dans $\cM$, on ait}
\[
x \ \in\Ob\underline{S} \qquad\Longrightarrow\qquad \underline{S}(\pi)/x \text{ \textit{cofinale dans} } \cM_\pi/x\text{\textit{.)}}
\]
\end{enumerate}

\dgrblock{dgrb1147}{The criterion in sets and in a pointed category}
\begin{dcomment}{examples}{p488}{The criterion in sets and in a pointed category}
The distinction between an arbitrary full subcategory and a sieve or cosieve is restrictive enough to make the final criterion concrete. In $\mathsf{Set}$, a nonempty cosieve containing $\varnothing$ is all of $\mathsf{Set}$, because $\varnothing$ maps to every set. If a cosieve contains a nonempty set $X$, the map $X\to1$ puts $1$ in the cosieve, and the maps $1\to Y$ put every nonempty $Y$ there. Thus the only replete cosieves are
\[
 \varnothing,\qquad\mathsf{Set}_{\ne\varnothing},\qquad\mathsf{Set}.
\]
The middle one is generated by $1$ and its complementary sieve consists only of $\varnothing$. The former has the usual filtered presentations by nonempty finite subsets; the latter is the terminal category and is closed under filtered colimits. This computes both sides of the criterion at the smallest relevant ranks.

In a pointed category with zero morphisms, any two objects admit a map. Hence every nonempty cosieve, and every nonempty sieve, is the whole category. For example, the no-map construction cannot give a nontrivial open--closed decomposition of abelian groups: every group has its zero map to every fixed group. This is why commutative unital $k$-algebras are a useful test here. Their morphisms must preserve the unit and the specified copy of $k$, so the absence of a map to $k$ is a genuine algebraic condition.
\end{dcomment}


\noindent\textsc{Démonstration}. La partie b) est une conséquence tautologique de a), que je vais prouver. Dire que $x$ est dans l'image essentielle de $\underline{S}$, signifie que
\[
x \ \simeq \ \varinjlim_{i\in I} x_i \ , \tag*{$(*)$}
\]
avec $I$ grande devant $\pi$ et les $x_i$ dans $\underline{S}(\pi)$. On sait que $(*)$ implique ($I$ étant filtrante, et $\cM$ étant équivalente à $\Ind_\pi(\cM)$ puisque $\pi$ utile pour $\cM$) que $I\longrightarrow\cM_\pi/x$ est cofinale (4.4.2, page 79). Or ce foncteur se factorise par le foncteur d'inclusion pleinement fidèle $\underline{S}(\pi)/x\incl\cM_\pi/x$, lequel est cofinal si et seulement s'il satisfait la condition F 1. Il est immédiat que cette condition est satisfaite quand elle est satisfaite pour le foncteur composé en question, ce qui est bien le cas
\printedpage{190}
% head: A. Grothendieck : Dérivateurs, Ch. XVIII, Édition des Universités de Montpellier II et Paris VII
\msp{page 490}
par hypothèse.

\noindent Inversement, supposons que $\underline{S}(\pi)/x$ soit cofinale dans $\cM_\pi/x$. Comme on a
\[
x \xleftarrow{\ \sim\ } \varinjlim_{\cM_\pi/x} x_i \ ,
\]
$\pi$ étant utile pour $\cM$, on aura aussi par cofinalité
\[
x \xleftarrow{\ \sim\ } \varinjlim_{\underline{S}(\pi)/\pi} x_i \ , % [sic] \underline{S}(\pi)/\pi -> \underline{S}(\pi)/x, the cofinal subcategory is \underline{S}(\pi)/x
\]
or $\cM_\pi$ étant grande devant $\pi$, il en est de même de la sous-catégorie pleine cofinale $\underline{S}(\pi)/x$. Donc $x$ est dans l'image essentielle de $\Ind_\pi\underline{S}(\pi)$, q.e.d.

\noindent\textbf{Lemme 4.16.7.12.} \textit{Soient $I$ une catégorie filtrante, $J$ une sous-catégorie fermée, $K$ la sous-catégorie ouverte complémentaire. On a alors les équivalences}
\[
J \ \neq \ \emptyset \iff J \text{ \textit{cofinale} } \iff K \text{ \textit{non cofinale} } \iff K \ \neq \ I \ .
\]

\noindent\textsc{Démonstration} immédiate.

\noindent Je reviens à la démonstration de 4.16.7.8. On applique le lemme 4.16.7.11 à $\cM$, $\pi$, et aux deux
\msp{page 491}
sous-catégories $\underline{S}$, $\underline{T}$ de $\cM$, en notant que pour $x\in\Ob\cM$, $\underline{S}(\pi)/x$ est une sous-catégorie fermée de $\underline{S}(\pi)$, dont le complémentaire ouvert est $\underline{T}(\pi)/x$, de sorte qu'on peut appliquer le lemme 4.16.7.12. On trouve :
\[
\begin{aligned}
\underline{S}\ \pi\text{-accessible} \quad &\overset{4.16.7.11}{\iff} \quad \forall\ x\ \in\ \Ob\cM,\ \text{on a } (x\ \in\ \Ob\underline{S}\iff\underline{S}(\pi)/x \text{ cofinale dans}\\
&\qquad\quad \cM_\pi/x)\\
&\iff \quad \forall\, x\in\Ob\cM,\ \text{on a } \Bigl( \qquad \underbrace{x\notin\Ob\underline{S}}_{\text{(i.e. $x\in\Ob\underline{T}$, par 4.16.7.12)}}\\
&\qquad\qquad\qquad \iff \qquad \underbrace{\underline{S}(\pi)/x \text{ non cofinale dans } \cM_\pi/x}_{\text{(i.e. $\underline{T}(\pi)/x$ cofinale dans $\cM_\pi$, par 4.16.7.12)}} \ \Bigr)\\
&\overset{4.16.7.11}{\iff} \quad \underline{T} \text{ est } \pi\text{-accessible.}
\end{aligned}
\]
On a donc l'équivalence de a) et b). D'autre part, $\underline{S}(\pi)/x$ cofinale dans $\cM_\pi/x$ signifie simplement $\underline{S}(\pi)/x\neq\emptyset$ (4.16.7.12), donc $\underline{S}$ $\pi$-accessible signifie que $\underline{S}$ est formée des $x\in\Ob\cM$ tels que $\underline{S}(\pi)\neq\emptyset$, ce qui signifie aussi que $\underline{S}$ est le fermé dans $\cM$ engendré par $\underline{S}(\pi)$, d'où a $\iff$ a$'$. D'autre part, l'équivalence b $\iff$ b$'$ n'est autre que 4.6.7.6 (dernière formulation). % [sic] 4.6.7.6 -> 4.16.7.6, the surrounding references are all to 4.16.7.x

\msp{page 492}

\section*{4.17\quad Retour sur les paratopos, et sur les opérations $\cM\boxtimes\cN$, $\underline{\operatorname{Hom}}_{!}(\cM,\cN)$}

\noindent\texttt{[Ce paragraphe n'existe pas.]}

\printedpage{191}
% head: A. Grothendieck : Dérivateurs, Ch. XVIII, Édition des Universités de Montpellier II et Paris VII

\msp{page 1}

\section*{Appendice}

\subsection*{Ultraproduits de corps, et catégories accessibles}

\noindent Soit $(k_\xi)_{\xi\in E}$ une famille de corps, indexé par un ensemble $E$. $\Bigl($ Soit $\Phi$ un filtre non trivial dans $E$, i.e. une partie $\Phi\subset\mathfrak{P}(E)$ telle que :
\begin{enumerate}
\item[a)] $\emptyset\notin\Phi$, i.e. $A\in\Phi\Longrightarrow A\neq\emptyset$.
\item[b)] $A,B\in\Phi\Longrightarrow A\cap B\in\Phi$.
\item[c)] Si $A\in\Phi$, $A\subset B$, alors $B\in\Phi$.
\item[d)] $\Phi\neq\emptyset$ ($\iff E\in\Phi$).
\end{enumerate}
(Cela revient à dire que $\Phi$ est une partie fermée de l'ensemble ordonné $\mathfrak{P}(E)$, \textit{filtrante à gauche}, et ne contient pas le plus petit élément de $\mathfrak{P}(E)$, i.e. non égale à $\mathfrak{P}(E)$ tout entier.) $\Bigr)$ Pour toute partie $J$ de $E$, posons
\[
k_J \ = \ \prod_{\xi\in J} k_\xi \ ,
\]
on a donc un système projectif
\[
(k_J)_{J\in\mathfrak{P}(E)} \ , \qquad \text{i.e. } J\subset J' \text{ donne } k_{J'}\longrightarrow k_J \ ,
\]
et le système
\[
(k_J)_{J\in\Phi}
\]
est un système inductif d'annaux, indexé par l'ensemble ordonné filtrant $\Phi^o$. Posons % [sic] d'annaux -> d'anneaux, misspelling
\[
k_\Phi \ = \ \varinjlim_{J\in\Phi} k_J \ .
\]
Si $\Psi$ est un filtre plus fin que $\Phi$, i.e.
\msp{page 2}
tel que $\Psi\supset\Phi$, on trouve un homomorphisme canonique
\[
k_\Phi \longrightarrow k_\Psi \ ,
\]
donc si $\cF(E)=\cF$ est l'ensemble des filtres sur $E$, ordonnés par inclusion, on trouve un système inductif
\[
\Phi \longmapsto k_\Phi \ : \ \cF \longrightarrow \Ann \ .
\]
On n'a pas utilisé encore que les $k_\xi$ sont des corps, en fait il suffirait que ce soient des objets d'une catégorie $\cA$ stable par petits produis et par petits $\varinjlim$ filtrantes. Prenons maintenant $\cA=\Ann$ (catégorie des anneaux unitaires, non nécessairement commutatifs), et supposons que les $k_\xi$ soient des corps. % [sic] produis -> produits, misspelling

\printedpage{192}
% head: A. Grothendieck : Dérivateurs, Ch. XVIII, Édition des Universités de Montpellier II et Paris VII

\noindent\textbf{Proposition 1.} \textit{Supposons que les $k_\xi$ ($\xi\in E$) soient des} corps (\footnote[203]{NB Cette condition est nécessaire si on veut que $k_\Phi$ soit un corps pour tout ultrafiltre.}). \textit{Pour que $k_\Phi$ soit un corps, il faut et il suffit que $\Phi$ soit un ultrafiltre, i.e. que pour toute décomposition $E=A\amalg B$, on ait $A\in\Phi$ ou $B\in\Phi$.} (NB Si $A\in\Phi$, on a $B\notin\Phi$, puisque sinon, on aurait $A\cap B\in\Phi$, donc $\emptyset\in\Phi$, absurde.)

\noindent Nécessité : Soit, pour toute partie $A$ de $E$,
\[
f_A\in k_E \qquad \text{défini par} \qquad f_A(\xi) \ = \ \left\{
\begin{array}{lll}
1_{k_\xi} & \text{si} & \xi\in A\\[1ex]
0_{k_\xi} & \text{si} & \xi\in E\backslash A \ .
\end{array}\right.
\]

\msp{page 3}
\noindent Soit
\[
f_{A,\Phi} \ \in k_\Phi
\]
l'image de $f_A$ par l'homomorphisme canonique $k_E\longrightarrow k_\Phi$. Si $E=A\amalg B$, on aura
\[
f_A \ + \ f_B \ = \ 1 \ , \qquad f_A\, f_B \ = \ 0 \ ,
\]
d'où
\[
f_{A,\Phi} \ + \ f_{B,\Phi} \ = \ 1 \ , \qquad f_{A,\Phi}\, f_{B,\Phi} \ = \ 0 \ .
\]
Si $k_\Phi$ est un corps, cela implique qu'on a $f_{A,\Phi}=0$ ou $f_{B,\Phi}=0$. Mais si p.ex. $f_{A,\Phi}=0$, on voit que cela signifie qu'il existe un $J\in\Phi$ tel que $f_A|J=0$, ce qui signifie aussi que $J\subset B$. De même, si $f_{B,\Phi}=0$, on aura $J\subset A$. Donc $\Phi$ est un ultrafiltre.

\dgrblock{dgrb1148}{Reduced products as germs on a filter}
\begin{dcomment}{exposition}{p2}{Reduced products as germs on a filter}
Grothendieck turns to fields because they offer a concrete way of separating the existence of many local solutions from the existence of one compatible solution. The categorical construction comes first: a filter specifies which parts of a product may be discarded, and a filtered colimit makes those restrictions into identifications. This is the reduced product; it is called an ultraproduct only when the filter is an ultrafilter.

Write $R=\prod_{\xi\in E}k_\xi$. Every element of $k_\Phi$ has a representative in $R$, since a section on $J\in\Phi$ extends by zero outside $J$. Two representatives $f,g$ have the same image precisely when their equalizer belongs to $\Phi$. Consequently
\[
 k_\Phi\cong R/I_\Phi,
 \qquad
 I_\Phi=\{f\in R:Z(f)\in\Phi\},
 \qquad Z(f)=\{\xi:f(\xi)=0\}.
\]
This formula includes the order reversal in the indexing category: shrinking $J$ moves forward in $\Phi^{\mathrm{op}}$. It also proves directly that $I_\Phi$ is a two-sided ideal. Intersect two sets on which two functions vanish to handle addition; multiplying a vanishing function by an arbitrary function preserves its zero set. Properness of the filter gives $1\notin I_\Phi$.

The proof of Proposition~1 uses exactly the dichotomy supplied by an ultrafilter. If $[g]\ne0$, its zero set is not in $\Phi$. For an ultrafilter the complementary support $A$ is therefore in $\Phi$. Define $h(\xi)=g(\xi)^{-1}$ on $A$, and put $h(\xi)=1$ elsewhere. Both $gh$ and $hg$ agree with $1$ on $A$, so $[h]$ is a two-sided inverse. This proof applies to division rings. In the reverse direction, the images of the complementary characteristic functions $e_A,e_{E\setminus A}$ are orthogonal idempotents summing to $1$. A division ring has no idempotents other than $0,1$, so one of the two sets belongs to the filter.

Several formulas in this short proof look like failures of transcription but are already printed in the edition. The conditions involving $f\ne\emptyset$ and $g(\xi)\ne\emptyset$ concern nonzero ring elements. The representative to invert is $g(\xi)$, and the concluding identity must mean $ff'=f'f=1$. The later sentence saying that every $k_\Phi$ is a field needs $\Phi\in\mathcal F_0$. For example, $E=\{0,1\}$ and $\Phi=\{E\}$ give $k_\Phi=k\times k$. These slips are retained in the source and interpreted by the argument above.

The field assertion is also an instance of \L o\'s's theorem: the axioms for division rings are first-order sentences, and ultraproducts preserve them. The modern construction and the theorem, including their proofs, can be read in \citet[Definitions~2.4--2.5 and Theorem~3.1]{Keisler2010Ultraproduct}. Grothendieck's elementary proof is especially useful here because it keeps visible the characteristic functions which will classify all maps to the ground field.
\end{dcomment}


\noindent Inversement, supposons que $\Phi$ soit un ultrafiltre. Prouvons que $k_\Phi$ est un corps, i.e. que pour $f\in k_\Phi$, $f\neq\emptyset\Longrightarrow f$ inversible. Notons que $f$ provient d'un élément d'un $k_J$ ($J\in\Phi$), donc est l'image canonique d'un $g\in k_E$. Dire que $f=0$ signifie que $\exists\ J\in\Phi$ tel que $g|J=0$. L'inverse $f\neq0$ signifie donc que pour tout $J\in\Phi$, on a $g|J\neq\emptyset$.

\msp{page 4}
\noindent Soit alors $A$ l'ensemble des $\xi\in E$ tel que $g(\xi)\neq\emptyset$, $B=E\backslash A$ l'ensemble des $\xi\in E$ tels que $g(\xi)\neq\emptyset$. On a donc $B\notin\Phi$ puisque $f\neq0$, et comme $\Phi$ est un ultrafiltre, on a donc $A\in\Phi$. Mais soit $g'\in k_E$ défini par $g'(\xi)=f(\xi)^{-1}$ si $\xi\in A$, $g'(\xi)=1$ si $\xi\in B$, et soit $f'$ l'image de $g'$ dans $k_\Phi$. On voit alors qu'on a % [sic] B = E\A ... tels que g(\xi) \neq \emptyset -> g(\xi) = 0, the complement of A is the set where g vanishes
\[
f\ f' \ = \ f' \ = \ f \ = \ 1 \ ,
\]
ce qui montre bien que $k_\Phi$ est un corps.

\noindent Supposons maintenant que tous les $k_\xi$ ($\xi\in E$) soient égaux à un même corps $k$. Alors les $k_J$ forment un système projectif de $k$-algèbres, les $k_\Phi$ un système inductif de $k$-algèbres. Si $k$ est commutatif, il en est de même des $k_J$, $k_\Phi$, et en particulier pour $\Phi$ un ultrafiltre, $k_\Phi$ est une \textit{extension} du corps $k$.

\noindent Revenons au cas général, mais les $k_\xi$ commutatifs. Soit $\cF_0\subset\cF$ l'ensemble des ultrafiltres sur $E$. Notons que pour tout $\Phi\in\cF$, $k_E\longrightarrow k_\Phi$ est surjectif. Donc si $\Phi\in\cF$, $k_\Phi$ s'identifie à un corps quotient de l'anneau $k_E$.

\printedpage{193}
% head: A. Grothendieck : Dérivateurs, Ch. XVIII, Édition des Universités de Montpellier II et Paris VII

\msp{page 5}
\noindent Pour tout anneau $\cA$, soit $\Max(\cA)$ le spectre maximal, i.e. l'ensemble des idéaux maximaux de $\cA$. On trouve donc une application canonique
\[
\left\{
\begin{array}{lll}
\cF_0 & \longrightarrow & \Max(k_E)\\[1ex]
\ \Phi & \longmapsto & \Ker(k_E\longrightarrow k_\Phi) \ \defeq \ \mathfrak{m}_\Phi \ .
\end{array}\right. \tag*{$(*)$}
\]

\noindent\textbf{Proposition 2.} \textit{L'application précédente est} bijective.

\noindent Prouvons d'abord qu'elle est injective, cela résulte aussitôt du

\noindent\textbf{Corollaire 1} (\footnote[204]{\textbf{NB} Pour $\Phi\in\cF$ quelconque, $\mathfrak{m}_\Phi\defeq\Ker(\chi_\Phi : k_E\longrightarrow k_\Phi)$ détermine $\Phi$, puisque
\[
A \ \in \ \Phi \qquad\iff\qquad \chi_\Phi(f_A) \ = \ 1
\]
(mais il faut supposer les $K_\xi\neq0$, i.e. $0_\xi\neq1_\xi\ \ldots$).}). \textit{Soit, pour $\Phi\in\cF_0$,}
\[
\chi_\Phi \ : \ k_E \longrightarrow k_\Phi
\]
\textit{l'application canonique. Alors pour $A\in\mathfrak{P}(E)$, on a}
\[
A \ \notin \ \Phi \qquad\iff\qquad \chi_\Phi(f_A) \ = \ 0 \ , \quad \text{\textit{i.e.} } f_A\in\mathfrak{m}_\Phi
\]
\textit{(i.e. $A\in\Phi\iff\chi_\Phi(f_A)\neq0$, i.e. $f_A\notin\mathfrak{m}_\Phi$)} (\footnote[205]{\textbf{NB} On a donc
\[
(\chi_\Phi(f_A) \ \neq \ 0) \qquad\iff\qquad (\chi_\Phi(f_A) \ = \ 1) \ .
\]}).

\noindent En effet, $\chi_\Phi(f_A)=0$ signifie qu'il existe $B\in\Phi$ tel que $f_A|B=0$, i.e. tel que $B\subset E\backslash A$ ; mais cela équivaut à $E\backslash A\in\Phi$, donc à $A\notin\Phi$ puisque $\Phi$ est un ultrafiltre, q.e.d.

\msp{page 6}
\noindent Reste à prouver que l'application $(*)$ est surjective. Soit donc $\mathfrak{m}$ un idéal maximal de $k_E$, et soit
\[
\Phi \ \defeq \ \{A\in\mathfrak{P}(E)\mid\chi(f_A) \ = \ 1\} \ ,
\]
où $\chi : k_E\longrightarrow k_E/\mathfrak{m}=\kappa(\mathfrak{m})$ est l'application canonique. Je dis que $\Phi$ est un ultrafiltre. En effet :
\begin{enumerate}
\item[a)] $\emptyset\notin\Phi$ car $f_\emptyset=0$, d'où $\chi(f_\emptyset)=0\neq1$ dans $k_\Phi$.
\item[b)] Si $A,B\in\Phi$, on a $A\cap B\in\Phi$, car $f_{A\cap B}=f_A\cdot f_B$, d'où $\chi(f_{A\cap B})=\chi(f_A)\chi(f_B)=1$.
\item[c)] Si $A\subset B$ et $A\in\Phi$, alors $B\in\Phi$. En effet, on a $f_A=f_Af_B$, d'où $\chi(f_A)=\chi(f_A)\chi(f_B)$, et comme $\chi(f_A)=1$, on trouve $\chi(f_B)=1$.
\item[d)] On a $E\in\Phi$, i.e. $\chi(f_E)=1$, or $f_E=1$, donc $\chi(f_E)=1$.
\end{enumerate}
Ainsi, $\Phi$ est un filtre -- on n'a utilisé que le fait que $A/\mathfrak{m}$ est un anneau $\neq0$ (pour conclure que $\emptyset\notin\Phi$). Prouvons que si $A/\mathfrak{m}$ est intègre, $\Phi$ est un \textit{ultrafiltre}, i.e. :
\begin{enumerate}
\item[e)] Si $A\notin\Phi$, alors $B\defeq E\backslash A\in\Phi$.
\end{enumerate}
En effet, on a $f_A+f_B=1$, $f_Af_B=0$, d'où

% ---- from der-ch18_p196-200.tex
\printedpage{194}
% head: A. Grothendieck : Dérivateurs, Ch. XVIII, Édition des Universités de Montpellier II et Paris VII

\noindent\hspace*{1em}\textcircled{a}\quad $\chi(f_A)+\chi(f_B)=1$, et\\
\hspace*{1em}\textcircled{b}\quad $\chi(f_A)\chi(f_B)=0$.\\
$\chi(f_A)\neq 1$, on conclut de \textcircled{a} $\chi(f_B)\neq 0$, d'où par \textcircled{b} $\chi(f_A)=0$, d'où de \textcircled{a} $\chi(f_B)=1$, d'où $B\in\Phi$.

\msp{page 7}
Prouvons enfin que $\mathfrak{m}=\mathfrak{m}_\Phi$. Comme $\mathfrak{m}_\Phi$ est un idéal maximal, il suffit de prouver
\[
\mathfrak{m}_\Phi\ \subset\ \mathfrak{m}\ .
\]
Or si $A\in\mathfrak{P}(E)$ est telle que $A\notin\Phi$, ou $f_A\in\mathfrak{m}_\Phi$ (cor.~1), on a aussi $f_A\in\mathfrak{m}$. Donc il suffit de prouver le

\textbf{Lemme.} \textit{Soit $\Phi\in\mathcal{F}_0$. Alors $\mathfrak{m}_\Phi$ est l'idéal engendré par les $f_A$ pour $A\in\mathfrak{P}(E)\backslash\Phi$. Plus précisément, les $f_A$ sont dans $\mathfrak{m}_\Phi$, et tout $f\in\mathfrak{m}_\Phi$ est de la forme $ff_A$, avec $A$ comme dessus.}

En effet, $f\in\mathfrak{m}_\Phi$ signifie qu'il existe $B\in\Phi$ tel que $f|B=0$. Mais cela signifie aussi que $f=ff_A$, où $A=E\backslash A$ \texttt{[plutôt} $A=E\backslash B$\texttt{]}, d'où $A\notin\Phi$.

Cela prouve la proposition 2, et en même temps :

\textbf{Corollaire 2.} \textit{Tout anneau quotient intègre de $k_E=\prod_{\xi\in E}k_\xi$ est un corps.}

Notons maintenant que si $\Phi\in\mathcal{F}$ est un filtre quelconque sur $E$, alors $k_\Phi$ est un anneau quotient de $k_E$, donc
\[
\operatorname{Max}(k_\Phi)\ \hookrightarrow\ \underbrace{\operatorname{Max}(k_E)}_{\simeq\,\mathcal{F}_0}\ . \tag*{$(*)$}
\]
Ceci posé, on trouve :

% a vertical bar is printed in the left margin beside the statement of Corollaire 3
\textbf{Corollaire 3} (\footnote[206]{Marche pour tout système d'anneaux $(k_\xi)_{\xi\in E}$ \textit{non nuls}, et $\Phi,\Psi\in\mathcal{F}$, $\Psi$ plus fin que $\Phi$, i.e. $\Phi\subset\Psi$ si et seulement si $\mathfrak{m}_\Phi\subset\mathfrak{m}_\Psi$ (pas la peine que $\Psi$ soit un ultrafiltre).}). \textit{L'inclusion précédente identifie} $\operatorname{Max}(k_\Phi)$ \textit{à l'ensemble des ultrafiltres $\Psi$ sur $E$ qui sont} plus fins \textit{que $\Phi$.}

En effet, si $\Psi$ est un tel ultrafiltre, on
\msp{page 8}
a le diagramme commutatif
\[
\begin{tikzcd}
k_E \arrow[r, "\chi_\Phi"] \arrow[rr, bend right=15, "\chi_\Psi"'] & k_\Phi \arrow[r] & k_\Psi
\end{tikzcd}\ , \tag*{$(**)$}
\]
ce qui montre que $\Psi$ \texttt{[plutôt} $\mathfrak{m}_\Psi$\texttt{]} est dans l'image de $(*)$. Inversement, supposons que $\Psi\in\mathcal{F}_0$ provienne de $\operatorname{Max}(k_\Phi)$, i.e. que l'on ait $\mathfrak{m}_\Phi\subset\mathfrak{m}_\Psi$ (où $\mathfrak{m}_\Phi$ est le noyau de $\chi_\Phi$), ou encore qu'on ait une factorisation $(**)$ de $\chi_\Psi$. Prouvons que $\Psi$ est plus fin que $\Phi$, i.e. que $\Psi\supset\Phi$. Soit donc $A\in\Phi$, prouvons que l'on a $A\in\Psi$. Mais $A\in\Phi$ implique \texttt{[le]} fait que $\chi_\Phi(f_A)=1$, donc son image dans $k_\Psi$ est égale à 1, donc on a $A\in\Psi$ par le corollaire 1.

\dgrblock{dgrb1149}{Idempotents recover the spectrum of a product}
\begin{dcomment}{context}{p7}{Idempotents recover the spectrum of a product}
Grothendieck needs more than the existence of some residue fields: he needs to know exactly which quotients can return the original field. The characteristic idempotents make this classification possible. They form a Boolean algebra inside the product ring, and each prime quotient must decide every Boolean alternative.

Let $R=\prod_E k_\xi$ with commutative field factors. For $f\in R$, put $f^*(\xi)=f(\xi)^{-1}$ when $f(\xi)\ne0$, and $f^*(\xi)=0$ otherwise. Then
\[
 ff^*=e_{\operatorname{supp}(f)},
 \qquad f=f^2f^*.
\]
In particular $R$ is a commutative von Neumann regular, or absolutely flat, ring. If $R\to D$ is a nonzero domain quotient, a nonzero image $\bar f$ satisfies $\bar f=\bar f^{\,2}\bar f^*$; cancellation gives $\bar f\bar f^*=1$. Thus every domain quotient is a field. This is a second proof of Corollary~2 which isolates the ring-theoretic reason that all prime ideals are maximal.

For an ultrafilter $U$, the ideal $I_U$ is maximal by Proposition~1. Conversely, if $\mathfrak p$ is prime, put
\[
 U_{\mathfrak p}=\{A\subseteq E:e_A\notin\mathfrak p\}.
\]
In the domain $R/\mathfrak p$, each idempotent is $0$ or $1$. The identities $e_{A\cap B}=e_Ae_B$ and $e_{E\setminus A}=1-e_A$ show that $U_{\mathfrak p}$ is an ultrafilter. If $Z(f)\in U_{\mathfrak p}$, then $e_{\operatorname{supp}(f)}\in\mathfrak p$ and $f=fe_{\operatorname{supp}(f)}\in\mathfrak p$. Hence $I_{U_{\mathfrak p}}\subseteq\mathfrak p$; maximality of the former gives equality. This proves the asserted bijection without assuming what it is meant to establish about the residue field.

There is a useful topological interpretation. Under this bijection the principal open subset $D(f)\subseteq\operatorname{Spec}R$ is
\[
 \{U:\operatorname{supp}(f)\in U\}.
\]
It is both open and closed, with complement determined by the zero set of $f$. Thus $\operatorname{Spec}R$ is the Stone space of the Boolean algebra $\mathcal P(E)$; its points are ultrafilters and its principal points are the coordinates of $E$. This is the Boolean representation viewpoint originating with \citet{Stone1936Boolean}. The proof just given supplies the particular identification used here; the reference is historical, not a substitute for the proof.

For a proper filter $\Phi$, passing from $R$ to $R/I_\Phi$ restricts this space to
\[
 \{U:U\supseteq\Phi\}
   =\bigcap_{A\in\Phi}\{U:A\in U\}.
\]
Indeed $I_\Phi\subseteq I_U$ if and only if $\Phi\subseteq U$, as one sees by applying ideal inclusion to $1-e_A$. This also explains the scope of footnote~206. The order-recovery argument works for arbitrary nonzero coefficient rings. The assertion that all prime quotients are fields uses the field hypothesis and is not part of that generalization.
\end{dcomment}


\printedpage{195}
% head: A. Grothendieck : Dérivateurs, Ch. XVIII, Édition des Universités de Montpellier II et Paris VII

\msp{page 9}
Supposons à présent que tous les $k_\xi$ sont égaux \texttt{[à]} un même anneau \fbox{$\ k\neq 0\ $}\,. Donc $k_J$, $k_\Phi$ sont des `anneaux sous $k$' $k\longrightarrow k_J\longrightarrow k_\Phi$, d'ailleurs $k\longrightarrow k_J$ (si $J\neq\emptyset$) et $k\longrightarrow k_\Phi$ sont injectifs. On se pose la question quand $k\longrightarrow k_\Phi$ est \textit{bijectif}, i.e. surjectif. Comme $k_\Phi\simeq k_E/\mathfrak{m}_\Phi$, cela signifie donc que $k_E=k+\mathfrak{m}_\Phi$, i.e. que tout $f\in k_E$ s'écrit
\[
f\ =\ c\ +\ g\qquad
\begin{cases}
c\ \in\ k\ , & \text{i.e. } c(\xi)=\underbrace{c_0}_{\in\,k}\ \forall\,\xi\in E\ ,\\[2ex]
g\ \in\ \mathfrak{m}_\Phi & \text{i.e. } \exists\,A\in\Phi \text{ avec } g|A=0\ .
\end{cases}
\]
Mais pour $f\in k_E$, $c\in k$ et $A\in\Phi$ fixés, l'existence de $g\in k_E$ telle que $g|A=0$ et $f=c+g$ équivaut à $f|A=c_A$ (fonction constante dans $A$, de valeur $c$), i.e. $A\subset f^{-1}(\{c\})$. Pour $f$, $c$ fixés, l'existence de $A\in\Phi$ satisfaisant la relation équivaut à $f^{-1}(\{c\})\in\Phi$. Et pour $f$ fixé, l'existence de $c$ équivaut à celle d'une fibre $E_c$ de $f:E\longrightarrow k$, qui soit dans $\Phi$. Ainsi :

\textbf{Proposition 3.} \textit{Soit $\Phi\in\mathcal{F}$. Pour que $k\longrightarrow k_\Phi$ soit surjectif, il faut et il suffit que pour toute famille de parties} essentiellement disjointes \textit{$(E_\xi)_{\xi\in k}$ de $E$ indexée par $k$, de réunion $E$, il existe un $\xi\in k$ tel que $E_\xi\in\Phi$. A fortiori, il} faut \textit{que $\Phi$ soit un ultrafiltre (puisque} $\operatorname{card}k\geqslant 2$ \textit{par hypothèse).}

\msp{page 10}
En effet, la donnée d'une telle famille $(E_\xi)_{\xi\in k}$ équivaut à la donnée d'une application $E\longrightarrow k$, i.e. d'un élément de $f_E$. % [sic] $f_E$ -> $k_E$, an application $E\to k$ is an element of $k_E$

\textbf{Scholie.} La propriété qu'on vient d'expliciter pour un filtre $\Phi$ sur $E$, ne dépend de $k$ que par l'intermédiaire de $\pi=\operatorname{card}k$. On peut aussi l'énoncer en disant que pour toute \textit{partition} $(E_i)_{i\in I}$ de $E$ en un ensemble $I$ des parties tel que $\operatorname{card}I\leqslant\pi$, il existe un (et un seul) des $E_i$ qui soit dans $\Phi$. On dit alors que $\Phi$ est un $\pi$-\textit{ultrafiltre.}

Si $\pi\geqslant 2$ est fini, cette propriété équivaut à dire que $\Phi$ est un ultrafiltre. Par ailleurs, quel que soit $\pi$, les ultrafiltres $\Phi$ \textit{triviaux} (i.e. tels qu'il existe $\xi\in E$ avec $\{\xi\}\in\Phi$, ou encore tel que $\Phi=\{A\in\mathfrak{P}(E)\mid\xi\in A\}$) sont des $\pi$-ultrafiltres. Si $\pi$ est infini, il n'est pas clair qu'il existe des $\pi$-ultrafiltres non triviaux. On sait qu'il n'en existe sur $E$ \textit{que} si $\operatorname{card}E\notin\mathfrak{U}_1$, où $\mathfrak{U}_1$ est le plus petit univers tel que $\mathbf{N}\in\mathfrak{U}_1$, i.e. tel
\msp{page 11}
que $\mathfrak{U}$ contienne comme élément un ensemble infini.

\dgrblock{dgrb1150}{The successor hidden in a pi-ultrafilter}
\begin{dcomment}{exposition}{p10}{The successor hidden in a pi-ultrafilter}
The ultrapower becomes trivial precisely when every function to the coefficient ring is almost constant. Grothendieck packages this condition as a property of the ultrafilter alone, so that the size of the field can later be chosen to exclude it. A small shift in cardinal convention is essential in comparing his terminology with modern set theory.

For an infinite cardinal $\pi$, his $\pi$-ultrafilter is an ultrafilter closed under intersections of \emph{at most} $\pi$ sets. With the usual convention that $\kappa$-complete means closed under intersections of fewer than $\kappa$ sets, it is therefore a $\pi^+$-complete ultrafilter. To verify this equivalence, suppose first that the intersection of a family $(A_i)_{i<\delta}$ in $U$, with $\delta\le\pi$, is not in $U$. Partition $E$ into the first failures
\[
 (E\setminus A_i)\cap\bigcap_{j<i}A_j
 \quad(i<\delta),
 \qquad\text{and}\qquad \bigcap_{i<\delta}A_i.
\]
Every cell lies outside $U$, contradicting the partition condition. Conversely, if a partition has at most $\pi$ cells and none belongs to $U$, all their complements belong to $U$; completeness would put the empty intersection of the complements in $U$. This proves both implications, including the case of singular $\pi$.

Let $\theta=|k|$. Proposition~3 says exactly that
\[
 k\longrightarrow k^E/U\text{ is onto}
 \quad\Longleftrightarrow\quad
 U\text{ selects a fibre of every map }E\longrightarrow k.
\]
For infinite $k$ this is $\theta^+$-completeness. For finite nonzero $k$ every ultrafilter has the property, because a finite partition always has one selected cell. Consequently every ultrapower of a finite field is canonically the field itself, while a nonprincipal ultrapower of an infinite field can be a proper extension. The elementary embedding supplied by \L o\'s's theorem does not decide surjectivity; compare \citet[Corollary~3.3]{Keisler2010Ultraproduct}.

A concrete infinite example keeps the two questions apart. Take $k=\mathbb Q$, $E=\mathbb N$, and a nonprincipal ultrafilter $U$. The class $t=[n\mapsto n]$ differs from every constant, because each fibre of the function $n\mapsto n$ is finite. More strongly, for every nonzero polynomial $P\in\mathbb Q[X]$, the zero set of $P(n)$ is finite, so $P(t)\ne0$. Thus $\mathbb Q^\mathbb N/U$ contains a transcendental element. It is an elementary extension of $\mathbb Q$, but it admits no $\mathbb Q$-algebra map back to $\mathbb Q$: any field homomorphism would be injective and would have to send the transcendental $t$ to an element of $\mathbb Q$.

The later filter-extension terminology used by \citet[Definition~2.3]{BrookeTaylorRosicky2017} follows the fewer-than convention. We will keep $\theta^+$ visible when moving from Grothendieck's partition condition to those compactness principles.
\end{dcomment}


Si $\mathfrak{U}=\mathfrak{U}_1$, et si $k$ est un corps infini, p.ex. $k=\mathbf{Q}$, alors pour tout ensemble $E\in\mathfrak{U}$ et pour tout ultrafiltre \textit{non trivial} $\Phi$ sur $E$, $k_E$ est une extension \textit{non triviale} de $k$. Mais si $\mathfrak{U}=\mathfrak{U}_\alpha$ est un univers quelconque, et $\pi$ le premier cardinal dans $\mathfrak{U}$ qui n'appartienne pas à un univers $\mathfrak{U}'\subset\mathfrak{U}$ strictement plus petit, je présume que pour tout $E\in\mathfrak{U}$, tout $\pi$-ultrafiltre sur $E$ est trivial. % [sic] $k_E$ est une extension -> $k_\Phi$ est une extension, the extension of $k$ attached to the ultrafilter $\Phi$ is $k_\Phi$

\printedpage{196}
% head: A. Grothendieck : Dérivateurs, Ch. XVIII, Édition des Universités de Montpellier II et Paris VII

\subsection*{Application à un exemple d'une sous-catégorie non accessible.}

Soit $\mathcal{M}$ une catégorie $\mathfrak{U}$-accessible, et $a\in\operatorname{Ob}\mathcal{M}$, d'où un foncteur
\[
\begin{aligned}
f^o\ :\ \mathcal{M}^o\ &\longrightarrow\ \mathrm{Ens}\\
x\ &\longmapsto\ \operatorname{Hom}_{\mathcal{M}}(x,a)\ \overset{\text{déf}}{=}\ a(x)\ ,
\end{aligned}
\]
foncteur commutant aux $\limproj$ quelconques, donc tel que $f:\mathcal{M}\longrightarrow(\mathrm{Ens})^o$ commute aux $\limind$ quelconques, et en particulier $f$ est accessible. Considérons la sous-catégorie strictement pleine
\msp{page 12}
\[
\underbrace{\mathrm{Ens}^*}_{=\,\mathrm{Ens}\backslash\{\emptyset\}}\ \subset\ \mathrm{Ens}
\]
des ensembles non vides, on sait qu'elle est accessible, mais elle n'est par stable par $\limproj$ de type $\mathrm{L}_\pi$, quelque grand qu'on prenne $\pi$. % [sic] n'est par stable -> n'est pas stable, typo

\dgrblock{dgrb1151}{The universe hypothesis and its precise boundary}
\begin{dcomment}{status}{p11}{The universe hypothesis and its precise boundary}
Grothendieck's first universe is chosen to exclude complete nonprincipal ultrafilters. This is a mathematical restriction on the available cardinals, not a harmless change in the notation for smallness. The link with measurable cardinals can be established directly, and it also clarifies the conjecture at the end of the page.

Suppose that $U$ is a nonprincipal $\pi^+$-complete ultrafilter on $E$, with $\pi$ infinite. Let $\mu$ be the least cardinal for which $E$ has a partition into $\mu$ sets, none belonging to $U$. Such a partition exists, by singletons, and $\mu>\pi$. Pushing $U$ forward along this partition gives an ultrafilter $V$ on $\mu$. It contains no singleton. Minimality of $\mu$ shows that $U$, and hence $V$, is closed under intersections of fewer than $\mu$ members: a failed intersection would yield a partition into fewer than $\mu$ null pieces by the first-failure construction. Thus $V$ is a nonprincipal $\mu$-complete ultrafilter, and $\mu$ is measurable.

For completeness, the inaccessibility needed here follows from the same argument. Every subset of $\mu$ of size less than $\mu$ is $V$-null. If $\mu$ were singular, a partition into $\operatorname{cf}(\mu)$ bounded pieces would contradict completeness. If $\mu\le2^\lambda$ for some $\lambda<\mu$, assign distinct subsets of $\lambda$ to the elements of $\mu$. For each coordinate of $\lambda$, $V$ chooses one of the two membership alternatives. The intersection of these fewer than $\mu$ chosen sets belongs to $V$, but contains at most one element, another contradiction. Hence $\mu$ is regular and a strong limit. This is the classical measure obstruction associated with \citet{Ulam1930Measure}; compare the complete-ultrafilter discussion in \citet[Section~5]{Keisler2010Ultraproduct}. The argument is included because the universe boundary is decisive for the example.

If $\mathfrak U=V_\Omega$ is the least Grothendieck universe containing an infinite set as an element, $\Omega$ is the first inaccessible cardinal. A set $E\in\mathfrak U$ has cardinality less than $\Omega$, whereas the construction above would produce an inaccessible $\mu\le|E|$. Therefore every such complete ultrafilter is principal. More generally, the hypothesis of Proposition~4 holds in $V_\Omega$ exactly when the measurable cardinals below $\Omega$ are bounded. In one direction an unbounded measurable cardinal supplies a sufficiently complete nonprincipal measure. In the other, a nonprincipal $\pi_0^+$-complete ultrafilter would produce a measurable cardinal above the chosen bound $\pi_0$.

The final conjecture on this page has a positive answer \emph{when its proposed cardinal exists}. Suppose $\pi$ belongs to no strictly smaller infinite universe inside $\mathfrak U$. A nonprincipal $\pi$-ultrafilter would produce a measurable $\mu>\pi$ in $\mathfrak U$. Then $V_\mu$ is a smaller universe containing $\pi$, which is impossible. There need not, however, be a first cardinal with the stated property: when the smaller inaccessible cardinals are cofinal in $\Omega$, every cardinal below $\Omega$ belongs to a smaller universe. This separates the valid conditional assertion from an unstated existence assumption.

No conclusion about all universes follows from the first-universe example. A universe containing a sufficiently large strongly compact cardinal permits the opposite compactness behavior, as the final commentary block will show. The elementary algebraic construction remains the same; the supply of ultrafilters changes.
\end{dcomment}


Considérons
\[
\underbrace{\underline{S}\ =\ \underline{S}_a}_{\overset{\text{déf}}{=}\,f^{-1}(\mathrm{Ens}^*)}\ \subset\ \mathcal{M}\ ,
\]
sous-catégorie strictement pleine de $\mathcal{M}$ formée des $x\in\operatorname{Ob}\mathcal{M}$ tels que
\[
a(x)\ \neq\ \emptyset\ ,\qquad\quad \text{i.e. } x\leqslant a \text{ pour le préordre canonique de } \mathcal{M}\ .
\]
C'est donc une sous-catégorie \textit{ouverte} de $\mathcal{M}$. Son complémentaire fermé $\underline{S}'=\underline{S}'_a$ est la sous-catégorie fermée image inverse, par $f$, de la sous-catégorie strictement pleine de Ens réduite au seul objet $\{\emptyset\}$. Cette sous-catégorie de Ens est la catégorie finale type. Elle est stable dans Ens par $\limind$ quelconques, et par $\limproj$ filtrantes quelconques. Ainsi, $\underline{S}'$ est stable dans $\mathcal{M}$ par $\limind$ filtrantes quelconques. Ainsi, $\underline{S}'$ est stable dans $\mathcal{M}$ par $\limind$ filtrantes quelconques, en particulier, elle est $\mathrm{L}_\pi$-stable pour $\pi$ assez grand. Explicitons :

\msp{page 13}
\[
\begin{array}{cl}
\underbrace{\underline{S}'\ =\ \underline{S}'_a}_{}\ \subset\ \mathcal{M} & \\
\Big\|\ \text{\scriptsize déf} & \text{sous-catégorie fermée de } \mathcal{M}\\
f^{-1}(\{\emptyset\})\ = & \text{formée des } x\ \in\ \operatorname{Ob}\mathcal{M} \text{ tels}\\
 & \text{que } a(x)=\emptyset.
\end{array}
\]

\dgrblock{dgrb1152}{Existence of points is the missing compatibility condition}
\begin{dcomment}{exposition}{p12}{Existence of points is the missing compatibility condition}
Grothendieck is testing a tempting inference: a representable functor behaves perfectly with respect to limits, so perhaps the property that its values are nonempty should also behave well. The obstruction is that an inverse limit asks for compatible elements. Separate nonemptiness supplies no such compatibility.

For any diagram with colimit $A$, the universal property gives
\[
 \operatorname{Hom}_{\mathcal M}(A,a)
   \cong\varprojlim_i\operatorname{Hom}_{\mathcal M}(A_i,a).
\]
The full subcategory of objects admitting a map to $a$ is therefore a sieve: if $B\to A$ and $A\to a$ exist, so does $B\to a$. Its complement is a cosieve. A nonempty colimit of objects of the complement cannot admit a map to $a$, since composing with even one structure map would contradict membership of that stage. Dually, a nonempty limit of objects in the sieve lies in the sieve, by composing any projection with a map to $a$. These observations prove the closure and retract assertions without choosing a compatible family of points.

Even highly directed inverse systems of nonempty sets can have empty limit. Given an infinite $\pi$, choose $|E|>\pi$ and let $I=[E]^{\le\pi}$, ordered by inclusion. It is $\pi^+$-directed. The inverse system $X_F=E\setminus F$, with inclusion maps for $F\subseteq G$, has nonempty terms, while
\[
 \varprojlim_{F\in I}X_F
    =\bigcap_{F\in I}(E\setminus F)=\varnothing.
\]
For the categorical counterexample, Grothendieck realizes the same disappearance of coordinates inside products of fields. Complete nonprincipal ultrafilters would create additional points not given by coordinates; the universe hypothesis removes precisely those additional points.

The statement that the contravariant representable is accessible needs its original convention. It sends existing colimits in $\mathcal M$ to limits of sets, or equivalently preserves them as a functor to $\mathrm{Set}^{\mathrm{op}}$. This does not assert that $\mathrm{Set}^{\mathrm{op}}$ is an accessible category in the modern sense. The example concerns the sieve of existence of a point, not a failure of representability to preserve its appropriate limits.
\end{dcomment}


Comme $\underline{S}'$, étant fermée dans $\mathcal{M}$, est stable par $\limind$ non vides et en particulier $\mathrm{L}_\pi$-stable pour $\pi$ grand, il y a plus de raison de s'attendre à ce que $\underline{S}'$ soit accessible dans $\mathcal{M}$ que pour son complémentaire $\underline{S}$, dont il n'est pas clair qu'il ait cette propriété de stabilité, mais seulement stabilité par $\limproj$ non vides, et en particulier par facteurs directs.

En fait, nous allons voir qu'il peut arriver que $\underline{S}$ et $\underline{S}'$ soient tous deux non accessibles dans $\mathcal{M}$ :

\textbf{Proposition 4.} \textit{Soient $\mathfrak{U}$ un univers, contenant un cardinal infini $\pi_0$ tel que pour tout $E\in\mathfrak{U}$, tout $\pi_0$-ultrafiltre sur $E$ soit trivial. Soit alors $k\in\mathfrak{U}$ un corps de cardinal $\pi'_0\geqslant\pi_0$, et soit $\mathcal{M}$ la catégorie}
\printedpage{197}
% head: A. Grothendieck : Dérivateurs, Ch. XVIII, Édition des Universités de Montpellier II et Paris VII
\msp{page 14}
\textit{des $k$-algèbres commutatives. Donc $\mathcal{M}$ est un $\mathfrak{U}$-paratopos, admettant $k$ comme objet initial. Soit $\underline{S}'$ la sous-catégorie strictement pleine, stable par petites $\limind$ non vides, formée des objets $A$ de $\mathcal{M}$ tels que $\operatorname{Hom}_{\mathcal{M}}(A,k)=\emptyset$, et $\underline{S}$ la sous-catégorie strictement pleine complémentaire, formée des objets $A$ tels que $\operatorname{Hom}_{\mathcal{M}}(A,k)=\emptyset$. Alors $\underline{S}'$ est une sous-catégorie} fermée \textit{dans $\mathcal{M}$, et $\underline{S}$ la sous-catégorie ouverte complémentaire, donc $\underline{S}$ (tout comme $\underline{S}'$) est stable par facteurs directs. Ceci posé :} % [sic] second $\operatorname{Hom}_{\mathcal{M}}(A,k)=\emptyset$ -> $\operatorname{Hom}_{\mathcal{M}}(A,k)\neq\emptyset$, $\underline{S}$ is the complement of $\underline{S}'$
\begin{enumerate}[label={}]
\item[\textit{a)}] \textit{Il n'existe pas de cardinal $\pi$ tel que $\underline{S}$ soit $L_\pi$-stable dans $\mathcal{M}$, a fortiori $\underline{S}$ n'est pas accessible dans $\mathcal{M}$.}
\item[\textit{b)}] (\footnote[207]{?, pas prouvé.}). \textit{Bien que $\underline{S}'$ soit $L_0$-stable dans $\mathcal{M}$ (i.e. stable par $\limind$ filtrantes, et même, on l'a dit, par toute $\limind$ non vide), $\underline{S}'$ n'est pourtant pas accessible dans $\mathcal{M}$.}
\end{enumerate}
\noindent\textsc{Démonstration}. L'assertion a) va résulter du

\textbf{Corollaire 1.} \textit{Sous les conditions de la proposition 4, soient $E\in\mathfrak{U}$ un ensemble infini, $\Phi_0$ le filtre des complémentaires des parties finies de $E$. Alors $\operatorname{Hom}_{\mathcal{M}}(k_{\Phi_0},k)=\emptyset$, a fortiori, pour tout filtre $\Phi$ sur $E$ plus fin que $\Phi_0$, on a $\operatorname{Hom}_{\mathcal{M}}(k_\Phi,k)=\emptyset$.}

\msp{page 15}
On a posé ici, bien sûr, pour tout filtre $\Phi$ sur $E$,
\[
k_\Phi\ =\ \limind_{J\in\Phi}k_J\ .
\]
On sait que pour toute $k$-algèbre commutative $\mathcal{A}$, $\operatorname{Hom}(\mathcal{A},k)$ est en correspondance 1-1 avec la partie de $\operatorname{Max}(\mathcal{A})$ formée des idéaux maximaux $\mathfrak{m}$ de $\mathcal{A}$ tels que $k\longrightarrow\mathcal{A}/\mathfrak{m}$ soit un isomorphisme. Si $\mathcal{A}$ est de la forme $k_\Phi$, $\Phi$ un filtre sur $E$, cet ensemble est en correspondance biunivoque avec l'ensemble des $\pi'_0$-ultrafiltres $\Psi$ sur $E$ plus fins que $\Phi$, où $\pi'_0=\operatorname{card}(k)$ (proposition 3, et cor.~3 de prop.~2, page 7). Comme $\pi'_0\geqslant\pi_0$, un tel ultrafiltre est trivial par hypothèse. Donc on trouve
\[
\operatorname{Hom}_{\mathcal{M}}(k_\Phi,k)\ \simeq\ \text{ensemble des } \xi\in E \text{ tels que } \{\xi\}=\textstyle\bigcap_{A\in\Phi}A\ . \tag*{$(*)$}
\]
Mais dire que $\Phi$ est plus fin que $\Phi_0$ signifie aussi que pour tout $\xi\in E$, on a $E\backslash\{\xi\}\in\Phi$, i.e. $\exists\,A\in\Phi$ avec $A\subset E\backslash\{\xi\}$, i.e. $\{\xi\}\notin\bigcap_{A\in\Phi}A$. Donc on a alors $\operatorname{Hom}_{\mathcal{M}}(k_\Phi,k)=\emptyset$.

La conclusion du corollaire 1 s'exprime par la relation
\[
k_\Phi\ \notin\ \operatorname{Ob}\underline{S}\ ,\qquad\quad \text{i.e. } k_\Phi\in\operatorname{Ob}\underline{S}'\ ,
\]
d'autre part, on a par définition
\[
k_\Phi\ =\ \limind_{J\in\Phi}k_J\ ,
\]
\msp{page 16}
et il est clair (comme ces $J\in\Phi$ sont $\neq\emptyset$), que les $k_J$ en question sont dans $\underline{S}$. Ainsi, pour un filtre $\Phi$ comme dans le corollaire 1, on trouve un système inductif d'objets de $\underline{S}$, indexé par $\Phi^o$, dont la limite inductive n'est pas dans $\underline{S}$. Donc le a) de la proposition 4 se déduit en notant que pour tout cardinal $\pi\in\mathfrak{U}$, on peut trouver un ensemble $E\in\mathfrak{U}$, et un filtre $\Phi$ sur $E$ plus fin que $\Phi_0$, tel que l'ensemble ordonné $\Phi^o$ soit grand devant $\pi$. Pour trouver un tel $\Phi$, il suffit de prendre $E$ de cardinal $\pi'>\pi$, et pour $\Phi$ l'ensemble
\printedpage{198}
% head: A. Grothendieck : Dérivateurs, Ch. XVIII, Édition des Universités de Montpellier II et Paris VII
des complémentaires des parties de $E$ de cardinal $\leqslant\pi$, de sorte que $\Phi^o$ est isomorphe à l'ensemble ordonné des parties de $E$ de cardinal $\leqslant\pi$, lequel est bien grand devant $\pi$. Cela prouve le a) de la proposition 4.

Prouvons maintenant le b). Notons le

\textbf{Lemme.} \textit{Soit $\mathcal{M}$ une catégorie accessible stable par $\limind$ de type $I_\alpha$, $\alpha$ un ordinal limite donné, $\underline{S}'$ une sous-catégorie fermée, $\pi$ un cardinal utile}
\msp{page 17}
\textit{pour $\mathcal{M}$, tel que} $\operatorname{card}\alpha\leqslant\pi$\textit{. Pour tout $x\in\operatorname{Ob}\underline{S}'$, choisissons une représentation $\omega$-$\pi$-adaptée ($\omega$ étant le premier ordinal infini)}
\[
x\ \simeq\ \limind_{i\in I_x}x_i
\]
\textit{(donc $I_x$ grande devant $\pi$, stable par $\limind$ de type $I_\alpha$, $i\longmapsto x_i:I\longrightarrow\mathcal{M}$ commute auxdites limites, enfin les $x_i\in\operatorname{Ob}\mathcal{M}_\pi$). Pour que $\underline{S}'$ soit $\pi$-accessible dans $\mathcal{M}$, il faut et il suffit que pour tout $x\in\operatorname{Ob}\underline{S}'$, il existe un $i\in I_x$ tel que $x_i\in\operatorname{Ob}\underline{S}'$.}

\dgrblock{dgrb1153}{A directed family of products loses every rational point}
\begin{dcomment}{examples}{p16}{A directed family of products loses every rational point}
The proof of Proposition~4(a) becomes quite explicit once the filter is written in terms of the coordinates it allows us to discard. Fix an infinite $\pi$, choose $E$ of cardinality greater than $\pi$, and set
\[
 \Phi_\pi=\{E\setminus F:F\subseteq E,\ |F|\le\pi\},
 \qquad
 A_\pi=k^E/I_\pi,
 \quad
 I_\pi=\{f:|\operatorname{supp}(f)|\le\pi\}.
\]
The set $\Phi_\pi$ is a proper filter, and its opposite order is $\pi^+$-directed: the union of at most $\pi$ exceptional subsets still has size at most $\pi$. Its reduced product is exactly $A_\pi$. Restriction of functions gives
\[
 A_\pi\cong
 \varinjlim_{F\in[E]^{\le\pi}} k^{E\setminus F}.
\]
Every stage has a $k$-point, evaluation at any surviving coordinate.

Now let $\theta=|k|$, and assume the hypothesis of Proposition~4. A $k$-algebra map $A_\pi\to k$ is necessarily surjective, since it fixes all constants. Its kernel therefore gives a maximal ideal, hence an ultrafilter $U$ extending $\Phi_\pi$. The residue field is the ultrapower $k^E/U$, and the structural map from $k$ to this field must be onto. Proposition~3 says that $U$ is $\theta^+$-complete. Since $\theta\ge\pi_0$, the universe hypothesis makes $U$ principal. But $\Phi_\pi$ contains $E\setminus\{\xi\}$ for every $\xi$, so no principal ultrafilter extends it. This contradiction proves that $A_\pi$ has no $k$-point.

The argument proves failure of closure for every allowed $\pi$; it does not merely show a failure of ordinary filtered-colimit closure. This is why it rules out accessibility of the inclusion at every eventual rank. It also identifies the exact role of the field hypothesis: a map to $k$ has a field quotient, whose maximal ideal can be classified by an ultrafilter.

The formula $(*)$ on manuscript page~15 should be interpreted as
\[
 \operatorname{Hom}_{k\text{-}\mathrm{alg}}(k_\Phi,k)
       \cong\bigcap_{A\in\Phi}A
\]
under the principal-ultrafilter hypothesis. Each member of the intersection gives an evaluation map. The printed condition $\{\xi\}=\bigcap\Phi$ would fail already for $\Phi=\{E\}$ with two coordinates. The corrected membership condition, rather than singleton equality, is exactly what the proof uses.
\end{dcomment}


En effet, notons que $\underline{S}'$ étant fermée, est stable dans $\mathcal{M}$ par toute $\limind$ non vide représentable dans $\mathcal{M}$, et en particulier par $\limind$ de type $I_\alpha$. Donc on peut appliquer 4.13.22 (p.~342) à $\mathcal{M}$, $\underline{S}'$, $\pi$ et à $\Phi=\{I_\omega\}$. Alors a) $\Longrightarrow$ c) implique que si $\underline{S}'$ est $\pi$-accessible dans $\mathcal{M}$, alors pour tout $x\in\operatorname{Ob}\underline{S}'$, l'ensemble $I_{x,S'}$ des $i\in I_x$ tels que $x_i\in\operatorname{Ob}\underline{S}$ est cofinal dans $I_x$. Cela implique qu'il est non vide. (Et réciproquement d'ailleurs, $\underline{S}'$ étant un fermé dans $\mathcal{M}$, puisque si $x_i\in\underline{S}'$, alors pour $j\in\operatorname{Ob}I_x$ avec $i\leqslant j$, on a $x_j\in\underline{S}'$, i.e. $I_{x,\underline{S}'}$ contient l'image essentielle de $i\backslash I_x$, laquelle est cofinale dans $I$ \dots) Cela prouve le `il faut'. Prouvons le `il suffit'. Donc prouvons que $\underline{S}'$ est égale à l'image essentielle de $\operatorname{Ind}_\pi(\underline{S}'(\pi))$. Elle la contient, puisque $\underline{S}'$ est stable par $\limind$ de type $\mathrm{L}_\pi$. % Doubt resolved against PDF leaf 40: the first S-prime subscript is not underlined, retained. % [sic] $x_i\in\operatorname{Ob}\underline{S}$ -> $x_i\in\operatorname{Ob}\underline{S}'$, the set $I_{x,\underline{S}'}$ concerns the $x_i$ in $\underline{S}'$

\msp{page 18}
Donc il reste à prouver que tout $x\in\operatorname{Ob}\underline{S}'$ est dans l'image essentielle. Or l'hypothèse nous assure que $x\simeq\limind\limits_{i\in I_x}x_i$ est isomorphe à $\limind\limits_{i\in I_{x,\underline{S}'}}x_i$ (car $I_{x,\underline{S}'}$ est cofinale dans $I_x$), où cette fois les $x_i$ sont dans $\underline{S}'$. Comme $I_x$ grande devant $\pi$, de même la sous-catégorie pleine cofinale $I_{x,\underline{S}'}$, donc on gagne.

\textbf{Corollaire.} \textit{Sous les conditions précédentes, si $\underline{S}'$ est $\pi$-accessible, il en est de même de l'ouvert complémentaire $\underline{S}$,} pourvu \textit{que $\underline{S}$ soit stable dans $\mathcal{M}$ par $\limind$ de type $L_\pi$.}

Il faut prouver que $\underline{S}$ est égale à l'image essentielle de $\operatorname{Ind}_\pi\underline{S}(\pi)$. L'hypothèse `pourvu' assure que cette image essentielle est $\subset\underline{S}$. Inversement, tout [\,?] $x\in\operatorname{Ob}\underline{S}$ est dans cette image essentielle. On aura
\[
x\ \simeq\ \limind_{i\in I}x_i\ ,\qquad\quad x_i\in\operatorname{Ob}\mathcal{M}_\pi,\ I=\mathcal{M}_\pi/x,\ \text{grande devant } \pi\ ,
\]
et comme $x\in\operatorname{Ob}\underline{S}$ et $\underline{S}$ ouverte, les $x_i$ sont dans $\underline{S}$, donc dans $\underline{S}(\pi)$, OK.

Dans le cas qui nous occupe, où $\mathcal{M}=(k\text{-alg.})$, comme $\underline{S}$ n'est justement \textit{pas} stable dans $\mathcal{M}$ par $\mathrm{L}_\pi$ (quel que soit $\pi$) (c'est cela qui nous faisait conclure que
\msp{page 19}
$\underline{S}$ n'est pas accessible dans $\mathcal{M}$), on ne peut donc conclure de cette inaccessibilité de $\underline{S}$, celle de $\underline{S}'$. Notons le

\dgrblock{dgrb1154}{The main text completes the abandoned argument}
\begin{dcomment}{status}{p18}{The main text completes the abandoned argument}
Grothendieck distinguishes the two possible failures with care. The open class of algebras with points fails a closure condition. The closed class of algebras without points already has every nonempty colimit, so its possible failure of accessibility must concern the absence of a small family of generators. The appendix reaches this distinction but does not finish the second argument; footnote~207 records that fact.

The main text has already supplied the missing implication. Proposition~4.16.7.8 and Corollary~4.16.7.9 identify accessibility of a cosieve with accessibility of its complementary sieve. Hence Proposition~4(a), under its stated universe hypothesis, proves Proposition~4(b) as well. The final question in the appendix must be read with this later result in view; it is not an unresolved assertion under those hypotheses.

Here is the mechanism, in the notation needed for the last page. Let $\pi\ge|k|$ be infinite. A $k$-algebra of cardinality at most $\pi$ has a presentation with at most $\pi$ generators and relations: take one generator for each element and record its addition, multiplication and scalar relations. It is consequently $\pi^+$-presentable. Every algebra is the $\pi^+$-directed union of its subalgebras of cardinality at most $\pi$.

If the closed class $S'$ were generated by its ambient $\pi^+$-presentable objects, an object $A\in S'$ would receive a map from one such object $B\in S'$. That map factors through a sufficiently large stage of any $\pi^+$-filtered presentation of $A$. Since $S'$ is a cosieve, that stage lies in $S'$. Conversely, one stage in $S'$ makes the full subcategory of $S'$-stages a nonempty cosieve in the filtered indexing category, hence cofinal; these stages then present $A$ inside $S'$. This proves the appendix's criterion without its additional ordinal-continuity apparatus.

The corollary on manuscript page~18 does not perform this argument: its explicit proviso assumes closure of $S$, and that closure is known to fail. Grothendieck correctly refuses the inference based on that corollary alone. The stronger equivalence in the main text is what repairs the proof. The next block will exhibit the missing small-generator obstruction directly in the very reduced products he proposes.
\end{dcomment}


% ---- from der-ch18_p201-201.tex
\printedpage{199}
% head: A. Grothendieck : Dérivateurs, Ch. XVIII, Édition des Universités de Montpellier II et Paris VII
% foot: Version 17 décembre 2025

\noindent\textbf{Corollaire.} \textit{Soit $\pi$ un cardinal infini. Pour que la sous-catégorie fermée $\underline{S}'$ de $\mathcal{M}$ (formée de $k$-algèbres commutatives $A$ telles que $\Hom_{\mathcal{M}}(A,k)=\emptyset$) soit $\pi$-accessible, il faut et il suffit que pour tout $k$-algèbre $A$ (de cardinal ${}>\pi$), telle que $\Hom_{\mathcal{M}}(A,k)=\emptyset$, existe une sous-algèbre $A_i$} de cardinal $\leqslant\pi$, \textit{telle que $\Hom_{\mathcal{M}}(A_i,k)=\emptyset$.} % [sic] tout k-algèbre -> toute k-algèbre, « algèbre » est féminin

\noindent(On applique le lemme, en prenant pour $\alpha$ le premier ordinal infini.)

\noindent\textbf{Question.} \textit{Cela est-il vérifié, en prenant simplement $\pi\geqslant\card k$~?} La question de l'existence d'un cardinal $\pi$, ayant la propriété indiquée plus haut, me paraît délicate, et j'ignore la réponse. Si on veut construire un contre-exemple pour tout $\pi$, les développements précédents suggèrent de prendre pour $A$ un produit infini d'extensions non triviales de $k$, ou un anneau $k_{\Phi}$ comme dans la démonstration de la proposition 4 a).


\dgrblock{dgrb1155}{The proposed reduced products answer every bound}
\begin{dcomment}{development}{p19}{The proposed reduced products answer every bound}
Grothendieck's final suggestion is especially effective because the reduced product can have many local evaluations even when it has no global point. The following calculation realizes that suggestion under the exact hypothesis of Proposition~4. We did not find this explicit small-subalgebra calculation in the sources checked here; its conclusion about nonaccessibility already follows from his own Corollary~4.16.7.9.

\textbf{Small subalgebras of a co-small reduced power.}\ours
Let $k$ be an infinite field with $|k|\ge\pi_0$ in the universe of Proposition~4, and let $\pi\ge|k|$ be an infinite cardinal in that universe. Take $E=\pi^+$ and
\[
 A=k^E/I_\pi,
 \qquad I_\pi=\{f\in k^E:|\operatorname{supp}(f)|\le\pi\}.
\]
Then $A$ has no $k$-point, but every $k$-subalgebra $B\subseteq A$ of cardinality at most $\pi$ has a $k$-point.

\emph{Proof.} The preceding product construction proves that $A$ has no point. The ideal is proper because $|E|>\pi$, and it is an ideal because finite unions of supports of size at most $\pi$ again have that size.

To construct a point on $B$, choose a representative $g_b\in k^E$ of each $b\in B$. For the embedded constants $c\in k$, choose the constant representative $g_c(\xi)=c$; constants are distinct in $A$, so these choices are consistent. Since addition and multiplication hold in the quotient, each of
\[
 g_{b+b'}-g_b-g_{b'},\qquad
 g_{bb'}-g_bg_{b'}
 \quad(b,b'\in B)
\]
has support of cardinality at most $\pi$. There are at most $|B|^2\le\pi$ such exceptional supports. Their union $F$ has cardinality at most $\pi$. Choose $\xi\in E\setminus F$. Then
\[
 h:B\longrightarrow k,\qquad h(b)=g_b(\xi)
\]
preserves addition and multiplication, and it fixes every constant. In particular $h(1)=1$, and $h$ is a $k$-algebra homomorphism. No compatibility between the initially chosen representatives was required; discarding the common exceptional set makes them compatible at one coordinate. This proves the assertion.\hfill$\square$

The poset of these subalgebras is $\pi^+$-directed: the algebra generated by the union of at most $\pi$ of them still has cardinality at most $\pi$. Their union is $A$. Thus the same algebra is both a directed colimit of products with points and a directed union of small subalgebras with points. The first presentation defeats closure of $S$; the second defeats the small-obstruction criterion for $S'$.

In particular, under Proposition~4's hypothesis the answer to the final question is negative for \emph{every} cardinal bound, not only for $\pi=|k|$. This establishes the statement marked ``pas prouv\'e'' in footnote~207 by a direct argument. The conditional nature of the conclusion is retained: the proof that $A$ has no point uses the absence of sufficiently complete nonprincipal ultrafilters. The local evaluation argument itself needs no such hypothesis.
\end{dcomment}


\dgrblock{dgrb1156}{The proposed first threshold fails without a universe hypothesis}
\begin{dcomment}{development}{p19}{The proposed first threshold fails without a universe hypothesis}
The previous example depends on the universe only when excluding a global point. There is a different obstruction at the proposed smallest threshold which works for every infinite field and in any ambient universe containing the construction. It simply asks a field to accommodate more distinct elements than it has.

\textbf{Pairwise distinct variables.}\ours
Let $k$ be an infinite field, $\theta=|k|$, and let $J$ have cardinality $\theta^+$. Put
\[
 B_J=k[x_j\mid j\in J]
       [(x_i-x_j)^{-1}\mid i\ne j].
\]
Then $B_J$ has no $k$-point, while every $k$-subalgebra of $B_J$ of cardinality at most $\theta$ has one.

\emph{Proof.} A homomorphism $B_J\to k$ would assign values $a_j\in k$ with $a_i-a_j$ invertible for every distinct pair. Thus $j\mapsto a_j$ would inject $J$ into $k$, which is impossible. The localization is nevertheless nonzero: it localizes the polynomial domain at nonzero elements.

Let $C\subseteq B_J$ have cardinality at most $\theta$. Each element of $C$ is represented by a polynomial fraction involving finitely many variables and finitely many inverted differences. Taking the union of these finite supports gives $J_0\subseteq J$ with $|J_0|\le\theta$ and $C\subseteq B_{J_0}$. Choose an injection $J_0\hookrightarrow k$. Evaluation at these distinct values makes every inverted difference nonzero, hence extends to $B_{J_0}\to k$. Its restriction is the required point of $C$.\hfill$\square$

We checked this formulation against the reduced-product and accessible-image sources cited in the appendix; we did not locate this particular presentation there. The construction is an elementary cardinality obstruction, with no claim of priority. It disproves the suggestion that the bound $\pi=|k|$ should suffice. It does not exclude a larger bound: when $\pi\ge\theta^+$, the displayed algebra itself is a small obstruction of cardinality $\theta^+$.
\end{dcomment}


\dgrblock{dgrb1157}{A compact cardinal restores a uniform obstruction bound}
\begin{dcomment}{later}{p19}{A compact cardinal restores a uniform obstruction bound}
Grothendieck's example anticipates a central feature of later accessibility theory: the existence of small witnesses can depend on compactness principles stronger than the existence of Grothendieck universes. Here the positive direction can be proved directly, so that the point of contact with the general theory remains explicit.

Let $\kappa>|k|$ be a regular uncountable cardinal such that every $\kappa$-complete filter extends to a $|k|^+$-complete ultrafilter. A strongly compact $\kappa>|k|$ satisfies this condition. In the terminology of \citet[Definition~2.3]{BrookeTaylorRosicky2017}, the weaker assumption is $L_{|k|^+,\omega}$-compactness. Assume the cardinal and the sets used below belong to the chosen universe.

We claim that if every $k$-subalgebra of an algebra $A$ of cardinality less than $\kappa$ has a $k$-point, then $A$ has a point. Let $I$ be the set of those subalgebras. Regularity of $\kappa$ and $|k|<\kappa$ imply that $I$ is $\kappa$-directed: the subalgebra generated by fewer than $\kappa$ of its members again has cardinality less than $\kappa$. The filter generated by the upper cones in $I$ is $\kappa$-complete and proper. Extend it to a $|k|^+$-complete ultrafilter $U$. For each $B\in I$, choose a homomorphism $h_B:B\to k$.

For $a\in A$, define a function on $I$ by $B\mapsto h_B(a)$ when $a\in B$, and give it the arbitrary value $0$ otherwise. The upper cone of subalgebras containing $a$ belongs to $U$. Since the function has at most $|k|$ values, completeness gives a unique value $h(a)\in k$ attained on a member of $U$. For $a,b\in A$, intersect the upper cone containing both with the three sets selecting $h(a),h(b),h(a+b)$. This intersection belongs to $U$ and is nonempty. At any $B$ in it,
\[
 h(a+b)=h_B(a+b)=h_B(a)+h_B(b)=h(a)+h(b).
\]
The same argument proves multiplicativity. For a constant $c\in k$, all $h_B$ take the value $c$, so $h(c)=c$. Hence $h:A\to k$ is a $k$-point. Taking the contrapositive yields a subalgebra of size less than $\kappa$ witnessing the absence of a point.

For a finite field the same proof needs only the ordinary ultrafilter lemma. Use the directed poset of finitely generated subalgebras and an ultrafilter extending its upper-cone filter. Every function to the finite set $k$ is constant on a member of this ultrafilter, without any additional completeness assumption. Therefore every algebra without a $k$-point has a finitely generated subalgebra without a $k$-point. Such a subalgebra has cardinality at most $\aleph_0$, so the final question has a positive answer with $\pi=\aleph_0$ for finite $k$, in every universe under consideration.

This is the usual complete-ultrafilter compactness argument specialized to Grothendieck's question, not a new compactness principle. It makes the relation to accessible images precise. The class $S$ of algebras with a $k$-point is the full image of the domain functor
\[
 (\mathrm{CAlg}_k)/k\longrightarrow\mathrm{CAlg}_k.
\]
It is closed under subobjects, so it is also that functor's powerful image. Theorem~3.4 of \citet{BrookeTaylorRosicky2017}, extending the work of Makkai and Par\'e, gives accessibility and accessible embedding of powerful images under appropriately large compactness assumptions. The elementary argument above proves the needed witness property in this particular case; Corollary~4.16.7.9 then also gives accessibility of the complementary cosieve.

The alternatives established here have definite hypotheses. In a universe with the cutoff assumed by Proposition~4, no uniform obstruction bound exists. In a universe containing a sufficiently large compact cardinal, such a bound exists. The exact weakest set-theoretic assumption for this one fixed field is not identified by these arguments, and the general powerful-image theorem should not be cited as an exact converse for this particular example.
\end{dcomment}
\begin{otherlanguage}{english}
\dgroverview{XVIII-e}{Chapter XVIII \textperiodcentered\ \textup{XVIII-e} \textperiodcentered\ Overview}
\begin{otherlanguage*}{english}
These final pages of Chapter~XVIII examine how far accessibility survives constructions which are elementary at the level of categories but delicate at the level of cardinal bounds \citep{Grothendieck1990XVIII}. The passage begins inside the proof of a gluing theorem, moves through categories of sections and categories over a small base, and ends by testing the limits of the theory against subcategories defined by the existence of a morphism. The appended calculation with products of fields supplies the decisive examples. The announced section~4.17 on paratoposes contains no mathematical text in the edition; the reader should not look to that heading for an additional theorem.

One convention governs almost every argument. An indexing category grand devant an infinite cardinal $\pi$ admits cocones for diagrams with at most $\pi$ arrows. Thus the source's $\mathrm L_\pi$ means existence or preservation of $\pi^+$-filtered colimits in the modern regular-cardinal convention. The same successor occurs in a $\pi$-ultrafilter: its intersections may have at most $\pi$ factors. A useful cardinal also controls a small family of presentable objects. These are separate requirements. An accessible category has arbitrarily large suitable ranks, but a proof using a rank chosen in advance must verify the additional uniformity it needs.

The gluing construction explains why these distinctions matter. A category over the walking arrow has an open fibre $\mathcal E_0$, a closed fibre $\mathcal E_1$, and a bifunctor $H(x_0,x_1)$ describing arrows between them. Presentability in a fibre need not be presentability in the glued category: the crossing Hom sets have to preserve the same filtered colimits. The calculations in SGA~4, Expos\'e~I, 9.6 and 9.8, provide the small-limit interchange used in the proof \citep{SGA4I1972}. The audit identifies two necessary repairs. Corollary~4.14.25(1) needs existence of the relevant colimits in the closed fibre. Corollary~4.15.2 needs a common initial rank controlling the stated functors before its subsequent cardinal enlargement. Explicit counterexamples show that these are hypotheses with mathematical content. The existential accessibility theorem survives the repaired choice of ranks.

The categories of sections express the same problem in a more organized form. Products collect the objects, inserters collect the transition maps, and equifiers impose their equations. This relates the source directly to the accessible-category closure results of Makkai and Par\'e, whose book appeared in 1989, and to the contemporary theory of flexible limits. The detailed modern construction in \citet[Appendix~A]{ChingRiehl2014} also explains why this machinery became useful in homotopical algebra. For the total category of a fibration or an opfibration, the commentary gives separate proofs of colimit existence and presentability. The two orientations require different uses of cartesian or cocartesian factorization; merely reversing a statement about filtered colimits would not supply either proof.

The small-base argument is the technical centre of the passage. A sufficiently filtered diagram in a small category can be reduced to one object and one recurrent idempotent. Its formal ind-object is the image of that idempotent, so it becomes representable exactly when the surviving projector splits. The walking unsplit projector, with monoid $\{1,p\}$ and $p^2=p$, is the simplest example. The source eventually proves the strict-bound conclusion
\[
 \operatorname{Ind}_{\pi}(\mathcal C)\simeq
 \operatorname{Kar}(\mathcal C)
 \qquad\bigl(\pi>|\operatorname{Mor}\mathcal C|\bigr).
\]
Here the notation on the left retains the source's successor convention. The commentary supplies both its growth argument and a presheaf proof. The literal identity condition in 4.15.9 is false, even for a two-element group; eventual invertibility and an appropriate trivialization are the invariant conclusions. A separate monoid reduction produces a quotient of a submonoid, not necessarily a quotient of the whole monoid. Neither defect invalidates the strict-bound theorem. The general equality case remains unsettled here, although two-sided cancellation and a stricter local Hom bound give affirmative special cases. Modern comparisons with Cauchy completion and coherent idempotents are provided by \citet{Tendas2022Continuity} and \citet[Sections~5.4.2--5.4.3]{Lurie2009HTT}.

The later questions expose several different ways in which accessibility can fail. The double-dual functor on sets fails to preserve filtered colimits at every rank. Isbell's category of partial bijections is self-dual and finitely accessible, yet has a one-object finite dense and codense subcategory and is not essentially small; this answers all three questions in 4.16.4 negatively \citep[Example~8.7]{AveryLeinster2021}. Other failures depend on the universe. Under the appendix's hypothesis, the commentary proves that the decreasing filtered-colimit closures of an accessible functor's image need not stabilize. A separate argument treats contravariant images. Adjoining retracts of representables to a nonaccessible cosieve also produces a category with a small dense family of presentable objects and sufficiently filtered colimits which is nevertheless not accessible. This repairs the precise defect which Grothendieck notices in his own proposed construction.

The question about existing colimits needs equal care. If all diagrams of a fixed shape have colimits, their colimit cocones form an accessible category, as in \citet[Example~5.4.7.9]{Lurie2009HTT}. The commentary constructs, under the appendix's hypothesis, a finitely accessible category whose existing pushout squares form a nonaccessible full subcategory. It is a collage obtained by adjoining one object to the category of $k$-algebras. Its pushout condition detects exactly the absence of a $k$-point. Thus closure under filtered colimits alone does not produce a small generating family.

The finite tests in 4.16.6 provide a positive result of a different kind. We prove the two split-idempotent tests by first constructing a colimit over the split image and then checking maps out of fibrewise presentable objects. A strict pullback can nevertheless retain an unsplit projector whose image in the base is an identity; the walking isomorphism gives a counterexample to the proposed shortcut b$'$. With all the tests retained, the intended relative property has a concise description. Over a small Karoubian base $\mathcal C$, it is equivalent to accessibility of $\mathcal E$ and $f:\mathcal E\to\mathcal C$, together with the lifting of every prescribed base isomorphism. Such an $f$ is an isofibration. Its strict pullbacks are equivalent to accessible pseudopullbacks, so the proposed change from small test bases to accessible test bases is valid. The standard closure theorem and its applications are discussed in \citet[Examples~4.8 and Remark~4.9, published version]{BourkeLack2024}.

The closing open--closed calculus supplies the conceptual explanation of the counterexamples. Open means a full sieve, closed means a full cosieve. A cosieve is accessible precisely when it has a small weakly initial family, and precisely when its complementary sieve is closed under sufficiently filtered colimits. The canonical slice of presentable objects over $x$ proves these equivalences: a nonempty cosieve in that filtered slice is cofinal, whereas a proper complementary sieve cannot be. This argument repairs the omitted slice denominators in the edition and completes the implication left unfinished in the appendix. The transfinite no-forward-morphism construction is also justified with an explicit universe bound. For the additional question restricted to fields, the answer is exact: in $V_\kappa$ such a $\kappa$-long sequence exists if and only if $\kappa$ is not a Vop\v{e}nka cardinal, with external families allowed. The commentary proves the reduction using a typed graph construction and the fully faithful functor of \citet{MillerPoonenSchoutensShlapentokh2018}.

The appendix identifies every maximal ideal of a product of fields with an ultrafilter. A map from the product back to an infinite ground field $k$ requires $|k|^+$-completeness. Its universe hypothesis says that sufficiently complete nonprincipal ultrafilters do not exist on small sets, equivalently that measurable cardinals below the universe boundary are bounded. Under that hypothesis, for every infinite $\pi\ge |k|$ the algebra
\[
 k^{\pi^+}/\{f:|\operatorname{supp}(f)|\le\pi\}
\]
has no $k$-point, while every subalgebra of size at most $\pi$ has one. A complete proof evaluates representatives outside their common exceptional set. A second construction disproves the proposed bound $\pi=|k|$ for every infinite field without the universe hypothesis. Conversely, a sufficiently large compact cardinal restores a uniform obstruction bound; for finite fields the ordinary ultrafilter lemma already gives a finitely generated obstruction. The comparison with powerful images uses \citet{BrookeTaylorRosicky2017}, without asserting an exact large-cardinal converse for this one fixed-field problem.

\end{otherlanguage*}
\end{otherlanguage}
\dgrchunkend
